\documentclass[11pt,reqno]{amsart}
% Shared preamble for the manuscript. Loaded by main.tex AFTER \documentclass.
% Section files must NOT redefine any of this.

\usepackage[margin=2.5cm]{geometry}
\usepackage{amsmath,amssymb,amsthm,mathtools}
\usepackage{setspace}
\setstretch{1.2}
\usepackage{enumitem}
\usepackage{microtype}
\usepackage[hidelinks]{hyperref}
\usepackage[capitalise,nameinlink]{cleveref}

% ---- theorem environments (amsart, numbered within section) ----
\theoremstyle{plain}
\newtheorem{theorem}{Theorem}[section]
\newtheorem{proposition}[theorem]{Proposition}
\newtheorem{lemma}[theorem]{Lemma}
\newtheorem{corollary}[theorem]{Corollary}
\theoremstyle{definition}
\newtheorem{definition}[theorem]{Definition}
\newtheorem{assumption}[theorem]{Assumption}
\theoremstyle{remark}
\newtheorem{remark}[theorem]{Remark}

% ---- macros ----
\newcommand{\R}{\mathbb{R}}
\newcommand{\N}{\mathbb{N}}
\newcommand{\Z}{\mathbb{Z}}
\newcommand{\C}{\mathbb{C}}
\newcommand{\Om}{\Omega}
\newcommand{\HH}{\mathcal{H}}
\newcommand{\eps}{\varepsilon}
\newcommand{\dT}{\delta_T}
\newcommand{\diso}{\delta_{\mathrm{iso}}}
\newcommand{\dstar}{\delta_{*}}
\newcommand{\AF}{\mathcal{A}}                 % Fraenkel asymmetry
\newcommand{\gd}{\mathrm{gd}}
\newcommand{\Etil}{\widetilde{E}}
\newcommand{\Exc}{\mathrm{Exc}}
\newcommand{\Per}{P}                          % perimeter
\newcommand{\dist}{\operatorname{dist}}
\newcommand{\supp}{\operatorname{supp}}
\newcommand{\diam}{\operatorname{diam}}
\newcommand{\sgn}{\operatorname{sgn}}
\newcommand{\norm}[1]{\left\lVert#1\right\rVert}
\newcommand{\abs}[1]{\left\lvert#1\right\rvert}
\newcommand{\set}[1]{\left\{#1\right\}}
\newcommand{\restr}[1]{\!\restriction_{#1}}
\DeclareMathOperator{\divg}{div}
\renewcommand{\d}{\,d}

% cleveref names
\crefname{equation}{}{}
\Crefname{equation}{}{}


\title[Sharp planar Saint--Venant stability without selection]
{Sharp planar quantitative Saint--Venant stability\\
by reduction to a graph inset}

\author{}
\date{\today}

\begin{document}

\begin{abstract}
We prove the sharp planar quantitative Saint--Venant inequality
$\AF(\Om)^2\le C\,\dT(\Om)$, where $\AF$ is the Fraenkel asymmetry and
$\dT$ is the torsional (Saint--Venant) deficit, with an absolute constant
$C$ and no regularity assumption on $\Om$. The proof reduces a
near-extremal $\Om$ to a small-amplitude radial graph by a purely
level-set construction governed by the two level deficits --- the
Cauchy--Schwarz/Serrin defect $D_H$ and the isoperimetric defect $D_I$ ---
and concludes with the perturbative stability of nearly spherical sets. In
particular the argument uses neither the Brasco--De~Philippis--Velichkov
selection principle, nor mass transport, nor symmetry-type (Serrin /
Magnanini--Poggesi) estimates. Two sharp ingredients drive the reduction:
a fold-content bound free of the quantitative-isoperimetric square-root
loss, obtained from a point-source flux identity and a thin-annulus
cancellation; and the observation that excising the low-torsion boundary
caps to manufacture a graph \emph{decreases} the scale-invariant torsional
deficit. The concluding nearly-spherical coercivity --- the perturbative
estimate that here replaces the selection principle --- is proved
self-containedly for graphs of small \emph{amplitude} and \emph{arbitrary
slope}, through an exact energy identity and a slope-free annular Steklov
lower bound; this is the regime the level-set reduction produces and the one
in which classical (slope-controlled) nearly-spherical theorems do not
apply. The reduction is carried out unconditionally except for two precisely
isolated analytic points --- the localisation of a harmonic boundary-layer
term in the cap-cutting step, and a quadratic mode-one smallness of the graph
profile at the optimal centre (replacing a barycentre normalisation) --- on
which the final estimate is conditional and which are stated as explicit
scoped hypotheses where used; every other step is proved in full, the only
external inputs being classical theorems quoted with precise references.
\end{abstract}

\maketitle

\tableofcontents

\section{Introduction}
\label{sec:intro}

\subsection*{The Saint--Venant inequality and its quantitative form}

Let $\Omega\subset\R^2$ be a bounded open set of finite perimeter with
$|\Omega|=\pi$.  The \emph{torsion function} of $\Omega$ is the unique
solution $u\in H^1_0(\Omega)$ of $-\Delta u=1$ in $\Omega$, $u=0$ on
$\partial\Omega$; its integral $T(\Omega):=\int_\Omega u$ is the
\emph{torsional rigidity}.  The classical Saint--Venant conjecture,
proved by P\'olya and Szeg\H{o} \cite{PolyaSzego1951} via the
Schwarz symmetrization of the torsion potential, asserts that the disc
maximises torsional rigidity at fixed area:
\[
  T(\Omega)\ \le\ T(B_1)\ =\ \frac{\pi}{8},
\]
with equality if and only if $\Omega$ is a disc.  A more precise, and
substantially more difficult, question is to quantify stability: how
large must $\Omega$ deviate from a disc when the deficit
$\dT(\Omega):=\tfrac\pi8-T(\Omega)$ is small?  The natural measure of
deviation is the \emph{Fraenkel asymmetry}
\[
  \AF(\Omega):=\min_{x_0\in\R^2}
    \frac{|\Omega\triangle B(x_0,r_\Omega)|}{|\Omega|},\qquad
  r_\Omega:=\sqrt{|\Omega|/\pi},
\]
which is zero if and only if $\Omega$ is a disc.  The sharp quantitative
Saint--Venant inequality is the statement
\begin{equation}\label{eq:sharp-main}
  \AF(\Omega)^2\ \le\ C\,\dT(\Omega)
\end{equation}
for some \emph{absolute} constant $C$ (independent of $\Omega$, of $|\Omega|$,
and of any regularity of $\partial\Omega$).  The exponent $2$ on the
left and the exponent $1$ on the right are both optimal: the
ratio $\AF(\Omega)^2/\dT(\Omega)$ is bounded and bounded away from
zero for a one-parameter family of ellipses.

Talenti \cite{Talenti1976} proved a pointwise comparison between the
decreasing rearrangement of $u_\Omega$ and the torsion function of the
equimeasure disc, giving $T(\Omega)\le T(B_1)$ and the quantitative
comparison $\mu(v)\le\mu_B(v)$ between the distribution functions.
This Talenti comparison, which we state precisely as \cref{thm:talenti},
is a key input in the present work.  Talenti's result also implies the
profile identity $\int_0^{1/4}(\mu_B-\mu)\d v=\dT$ (see \cref{lem:profile}),
which localizes the deficit to a single integral of a non-negative
integrand.

\subsection*{Prior work on quantitative stability}

The Faber--Krahn problem for the first Dirichlet eigenvalue $\lambda_1$
is the structural analogue of the Saint--Venant problem, and the two
share much of their quantitative theory.  For the quantitative
isoperimetric inequality the sharp form
$\AF(E)\le C_{\mathrm{F}}\,\diso(E)^{1/2}$ (where $\diso(E)$ denotes the
isoperimetric deficit and $C_{\mathrm F}$ is an absolute constant; see
\cref{thm:fmp}) was established by Fusco, Maggi, and Pratelli
\cite{FMP2008}; this result, which we use as a black box, is the
\emph{only} external quantitative inequality the present paper imports.

For Saint--Venant/Faber--Krahn stability, Brasco and Pratelli
\cite{BrascoPratelli2012} proved a quantitative stability result for
both the torsional rigidity and the first eigenvalue, establishing
\eqref{eq:sharp-main} with a \emph{non-sharp} exponent.  The sharp
inequality \eqref{eq:sharp-main} --- with exponent $2$ and an absolute
constant, valid with no regularity on $\partial\Omega$ --- was
established by Brasco, De Philippis, and Velichkov \cite{BDV2015}.
Their argument proceeds by a \emph{selection principle}: from a
near-extremal sequence one extracts a subsequence that, after
appropriate normalisation, converges in a strong topology to a
$C^{1,\gamma}$ domain; the stability of the latter is then proved by a
second-variation analysis.  The selection principle is a powerful and
elegant tool, but it is non-constructive, non-quantitative in the
convergence, and it relies on the regularity theory for minimizers of
shape functionals in a fundamental way.  A survey of quantitative
spectral inequalities, including the state of the art up to~2017, is
given by Brasco and De Philippis \cite{BrascoDePhilippis2017}.

\subsection*{The main theorem}

The main result of this paper is the following, proved in \cref{sec:main}
as \cref{thm:main}.

\begin{theorem}[Main theorem; see \cref{thm:main}]
  There exists an absolute constant $C$ such that every bounded open set
  $\Omega\subset\R^2$ with $|\Omega|=\pi$ satisfies
  \[
    \AF(\Omega)^2\ \le\ C\,\dT(\Omega).
  \]
\end{theorem}

This is the same conclusion as \cite{BDV2015}, but obtained by an
entirely different route.  The proof requires no selection principle,
no mass transport, and no symmetry-type argument (Serrin overde\-termined
problem / Magnanini--Poggesi estimates).  The only external inputs are:
classical function-space and geometric-measure-theory tools
(the coarea formula, the planar isoperimetric inequality, the Talenti
comparison, and the sharp quantitative isoperimetric inequality of
\cite{FMP2008}), all quoted with precise references in
\cref{sec:prelim}.

\subsection*{Strategy and what is new}

The argument replaces the selection--regularity--second-variation
architecture of \cite{BDV2015} by a purely level-set construction
driven by the two \emph{level deficits}:
\begin{align*}
  D_H(v) &:= \mu(v)\bigl(-\mu'(v)\bigr) - P(\{u>v\})^2\ \ge\ 0,\\
  D_I(v) &:= P(\{u>v\})^2 - 4\pi\mu(v)\ \ge\ 0,
\end{align*}
where $\mu(v)=|\{u>v\}|$ is the distribution function and
$P(\{u>v\})$ is the perimeter of the superlevel set.  The central
identity (\cref{thm:level-identity}, proved in \cref{sec:identity})
\[
  \int_0^m\bigl(D_H(v)+D_I(v)\bigr)\d v\ =\ 4\pi\,\dT(\Omega)
\]
encodes the entire deficit in a non-negative integral over the level
parameter.  All the estimates in the paper are consequences of this
identity and the classical Talenti comparison.

\smallskip
\noindent\textbf{The two sharp ingredients.}
Two novel observations make the reduction exact:

\begin{enumerate}
\item \emph{Fold-content bound with no square-root loss.}
  At the optimal inset level $\hat v\sim\sqrt{\dT}$, the connected
  component $\Etil$ of $\{u>\hat v\}$ of largest measure is trapped in a
  thin annulus $B_{\rho_i}(z)\subset\Etil\subset B_{\rho_e}(z)$ of
  width $\eta=\rho_e-\rho_i\lesssim\dT^{1/8}$ (established in
  \cref{sec:annulus}).  The total angular measure of rays along which
  $\Etil$ is \emph{not} a radial graph is bounded, via the
  point-source flux identity and the cancellation
  $|I_{\Etil}-2\pi\hat r|\le a\eta/(\rho_i\hat r)$
  (where $a=|\Etil\triangle B_{\hat r}(z)|$), by a quantity of order
  $\Exc+a\eta$ — without the factor of $\sqrt{\diso}$ that the
  naive application of the quantitative isoperimetric inequality
  \cite{FMP2008} would introduce.  This is the content of
  \cref{prop:fold,cor:gd} in \cref{sec:fold}.

\item \emph{Cutting caps decreases the deficit.}
  To manufacture a genuine radial graph $G$ from $\Etil$ one must excise
  the ``caps'' (the portions of $\Etil$ that cause each ray to intersect
  $\partial\Etil$ more than once).  A potential concern is that this
  excision increases the scale-invariant deficit
  $\dstar(G):=(T(B_G)-T(G))/T(B_G)$.  In fact it does not: the
  equal-area comparison ball shrinks faster than torsion is lost, so
  $\dstar$ \emph{decreases}.  This is \cref{lem:cutdeficit} in
  \cref{sec:cut}, proved via the Green-function representation of
  torsional rigidity and a barrier estimate in the collar.
\end{enumerate}

With $\tau:=\sqrt{\dT}$, the two observations combine to produce a
continuous radial graph $G$ with amplitude $\lesssim\dT^{1/8}$,
scale-invariant deficit $\dstar(G)\le C\dstar(\Omega)$, and symmetric
difference $|\Omega\triangle G|\le C\sqrt{\dT}$ (the precise reduction
is \cref{thm:reduction} in \cref{sec:main}).  The
perturbative stability of nearly spherical sets
(\cref{thm:graphG} in \cref{sec:nearly}) then gives $\AF(G)^2\le C_G\dstar(G)$,
and the triangle inequality for the symmetric difference transfers this
to $\AF(\Omega)^2\le C\dT(\Omega)$.

\subsection*{Organization of the paper}

\Cref{sec:prelim} collects the classical prerequisites: the theory of
sets of finite perimeter (De Giorgi--Federer structure theorem, coarea
formula, differentiation of the distribution function), the planar
isoperimetric inequality and the sharp quantitative version
\cite{FMP2008}, the interior regularity and analyticity of the torsion
function \cite{GilbargTrudinger2001}, and Sard's theorem.

\Cref{sec:identity} introduces the two level deficits $D_H$ and $D_I$,
proves their non-negativity, establishes the central level-set deficit
identity $\int_0^m(D_H+D_I)\d v=4\pi\dT$ (\cref{thm:level-identity}),
and derives the Talenti comparison and the profile identity.

\Cref{sec:inset} carries out the level-set construction: it proves
deficit transfer to any superlevel set, selects the good level $\hat v$
by an averaging argument, establishes the closeness of $\{u>\hat v\}$
to $\Omega$, and isolates the inset $\Etil$ as the dominant connected
component.

\Cref{sec:annulus} proves that a connected hole-free finite-perimeter
set with controlled isoperimetric deficit is trapped in a thin
concentric annulus, using the angular-perimeter coarea identity and a
crossing-number argument.

\Cref{sec:fold} bounds the fold content of $\Etil$ --- the obstruction
to being a radial graph --- via the point-source flux identities and the
thin-annulus cancellation, achieving the sharp estimate without the
quantitative-isoperimetric square-root loss.

\Cref{sec:cut} proves that excising the caps to manufacture a radial
graph decreases the scale-invariant torsion deficit, via a
Green-function comparison and a barrier estimate.

\Cref{sec:nearly} establishes the perturbative stability of nearly
spherical sets for the torsion functional, proving that any radial graph
of small amplitude satisfies $\AF^2\le C_G\dstar$, with explicit
absolute constants, by a second-order expansion of the torsion deficit
in Fourier modes.

\Cref{sec:main} assembles all pieces, optimises the level parameter at
$\tau=\sqrt{\dT}$, and derives the main theorem in full.

\subsection*{Notation}

Throughout the paper, $\Omega\subset\R^2$ is a bounded open set with
$|\Omega|=\pi$.  The unit disc is $B_1$; $B(x_0,r)$ denotes the open
disc of radius $r$ centred at $x_0$.  The symbol $C$ denotes a generic
positive absolute constant whose value may change from line to line;
named constants such as $C_G$, $\eta_G$, $A_0$, $\delta_*$ are fixed
at their first occurrence.  All \emph{absolute} constants are
independent of $\Omega$ and of any geometric or regularity data of
$\partial\Omega$.  The notation $f\lesssim g$ means $f\le Cg$ for an
absolute $C$.

\section{Preliminaries}\label{sec:prelim}

Throughout this paper $\Omega\subset\R^2$ denotes a bounded open set,
$|\,\cdot\,|$ denotes Lebesgue measure on $\R^2$, $\HH^1$ denotes
one-dimensional Hausdorff measure, and $B_r(x)\subset\R^2$ is the open disc
of radius $r$ centred at $x$; we write $B_r = B_r(0)$.  The fixed
normalisation $|\Omega|=\pi$ is in force from \cref{sec:intro} onward but is
not required for the statements in this section unless explicitly noted.  We
collect here the definitions and classical results the rest of the paper relies
on.  Standard results are stated precisely with full hypotheses and cited to
published sources.

%--------------------------------------------------------------------
\subsection{Sets of finite perimeter}
%--------------------------------------------------------------------

We follow the conventions of \cite{Maggi2012}.

\begin{definition}[Perimeter]\label{def:perimeter}
Let $E\subset\R^2$ be Lebesgue measurable and let $A\subset\R^2$ be open.
The \emph{perimeter of $E$ in $A$} is
\[
  \Per(E;A) := \sup\!\left\{\,\int_E \divg\varphi\,dx \;:\;
    \varphi\in C^1_c(A;\R^2),\ \|\varphi\|_{L^\infty}\le 1\,\right\}.
\]
We set $\Per(E):=\Per(E;\R^2)$.  The set $E$ is of \emph{finite perimeter} (in
$A$) if $\Per(E;A)<\infty$.
\end{definition}

If $E$ has locally finite perimeter, the distributional gradient
$\mu_E := -D\mathbf{1}_E$ is an $\R^2$-valued Radon measure with total
variation $|\mu_E|(A)=\Per(E;A)$
\cite[Proposition~12.1, Definition~12.1]{Maggi2012}.

\begin{definition}[Reduced boundary]\label{def:reduced-boundary}
For $E$ of locally finite perimeter, the \emph{reduced boundary} $\partial^*E$
is the set of $x\in\supp|\mu_E|$ at which
\[
  \nu_E(x) := \lim_{r\to 0^+}\frac{\mu_E(B_r(x))}{|\mu_E|(B_r(x))}
\]
exists and satisfies $|\nu_E(x)|=1$.  The vector $\nu_E(x)$ is the
\emph{measure-theoretic outer unit normal} to $E$.
\end{definition}

\begin{theorem}[De~Giorgi structure theorem]\label{thm:de-giorgi}
Let $E\subset\R^2$ have locally finite perimeter.  Then $\partial^*E$ is
countably $\HH^1$-rectifiable, $\mu_E = \nu_E\,\HH^1\lfloor\partial^*E$, and
for every Borel set $A$,
\[
  \Per(E;A) = \HH^1(\partial^*E\cap A).
\]
\end{theorem}

\begin{proof}
See \cite[Theorem~15.9, Corollary~16.1]{Maggi2012} or
\cite[Section~5.7, Theorem~5.15]{EvansGariepy2015}.
\end{proof}

\begin{definition}[Essential boundary]\label{def:essential-boundary}
For a measurable set $E\subset\R^2$, its \emph{essential boundary} is
$\partial^e E := \R^2\setminus(E^{(0)}\cup E^{(1)})$, where
$E^{(t)}:=\{x:\lim_{r\to0}|E\cap B_r(x)|/|B_r(x)|=t\}$.
\end{definition}

\begin{theorem}[Federer]\label{thm:federer}
If $E\subset\R^2$ has locally finite perimeter then
$\partial^*E\subseteq E^{(1/2)}\subseteq\partial^e E$ and
$\HH^1(\partial^e E\setminus\partial^*E)=0$.  In particular,
$\Per(E;A)=\HH^1(\partial^e E\cap A)$ for every Borel $A$.
\end{theorem}

\begin{proof}
See \cite[Theorem~15.5, Theorem~16.2]{Maggi2012} or
\cite[Theorem~5.16]{EvansGariepy2015}.
\end{proof}

%--------------------------------------------------------------------
\subsection{The coarea formula and differentiation of the distribution
function}
%--------------------------------------------------------------------

\begin{theorem}[Coarea formula]\label{thm:coarea}
Let $\Omega\subset\R^2$ be open and let $u\in W^{1,1}(\Omega)$.  Then for
every nonnegative Borel function $g:\Omega\to[0,\infty]$,
\begin{equation}\label{eq:coarea}
  \int_\Omega g(x)\,|\nabla u(x)|\,dx
  = \int_{-\infty}^{\infty}\!\!\left(\int_{\{u=t\}} g\,d\HH^1\right)dt.
\end{equation}
Moreover, for Lebesgue-almost every $t\in\R$ the superlevel set $\{u>t\}$ has
finite perimeter in $\Omega$, the level set $\{u=t\}$ is $\HH^1$-rectifiable
with $\HH^1(\{u=t\})<\infty$, and
\[
  \Per(\{u>t\};\Omega) = \HH^1(\{u=t\}) \quad\text{for a.e.\ } t.
\]
\end{theorem}

\begin{proof}
The coarea formula \eqref{eq:coarea} for $W^{1,1}$ functions is
\cite[Theorem~13.1]{Maggi2012}; see also
\cite[Section~3.4, Theorem~3.13]{EvansGariepy2015}.  The perimeter identity and
rectifiability of almost every level set follow from the $BV$-coarea formula
\cite[Theorem~3.40]{AFP2000}: since
$\int_\R\Per(\{u>t\};\Omega)\,dt = \int_\Omega|\nabla u|\,dx < \infty$, the
set $\{u>t\}$ has finite perimeter for a.e.\ $t$.
\end{proof}

The following corollary, which we call the \emph{differentiation lemma} for the
distribution function, is used repeatedly in the level-set analysis.

\begin{lemma}[Differentiation of the distribution function]\label{lem:dist-deriv}
Let $\Omega\subset\R^2$ be bounded and let $u\in C^1(\Omega)\cap C(\overline\Omega)$
with $u\ge 0$.  Define $\mu(v):=|\{u>v\}|$ for $v\ge 0$.  Then $\mu$ is
non-increasing and absolutely continuous, and for Lebesgue-almost every $v>0$,
\begin{equation}\label{eq:dist-deriv}
  -\mu'(v) = \int_{\{u=v\}} \frac{1}{|\nabla u|}\,d\HH^1.
\end{equation}
In particular, the Cauchy--Schwarz inequality applied to the pair
$(1, |\nabla u|^{1/2})$ on $\{u=v\}$ gives
\begin{equation}\label{eq:cauchy-schwarz-flux}
  \Per(\{u>v\})^2 = \HH^1(\{u=v\})^2
  \le \left(\int_{\{u=v\}}|\nabla u|\,d\HH^1\right)\left(-\mu'(v)\right)
  \quad \text{for a.e.\ }v.
\end{equation}
\end{lemma}

\begin{proof}
The coarea formula \eqref{eq:coarea} with $g = |\nabla u|^{-1}\mathbf{1}_{\{\nabla u\ne 0\}}$
gives, for $0\le s < t$,
\[
  |\{s < u \le t\}\cap\{\nabla u\ne 0\}|
  = \int_s^t \left(\int_{\{u=v\}\cap\{\nabla u\ne 0\}}\frac{1}{|\nabla u|}\,d\HH^1\right)dv.
\]
Since $u\in C^1(\Omega)$, Sard's theorem (\cref{thm:sard} below) implies that
$\{u=v\}\cap\{\nabla u=0\}$ is empty for a.e.\ $v$, so the restriction to
$\{\nabla u\ne 0\}$ is superfluous for a.e.\ $v$.  Also $|\{\nabla u=0\}\cap\{s<u\le t\}|$
contributes zero to the left side (it has measure zero by a standard argument).
Hence $\mu(s)-\mu(t) = \int_s^t(-\mu')\,dv$ with $-\mu'(v) =
\int_{\{u=v\}}|\nabla u|^{-1}\,d\HH^1$ a.e.

For \eqref{eq:cauchy-schwarz-flux}: by Cauchy--Schwarz on $\{u=v\}$,
\[
  \HH^1(\{u=v\})^2
  = \left(\int_{\{u=v\}}1\,d\HH^1\right)^{\!2}
  \le \left(\int_{\{u=v\}}|\nabla u|\,d\HH^1\right)
      \left(\int_{\{u=v\}}\frac{1}{|\nabla u|}\,d\HH^1\right)
  = \left(\int_{\{u=v\}}|\nabla u|\,d\HH^1\right)(-\mu'(v)),
\]
and by \cref{thm:coarea}, $\HH^1(\{u=v\}) = \Per(\{u>v\})$ for a.e.\ $v$.
\end{proof}

\begin{remark}\label{rem:flux-coarea}
The integral $\int_{\{u=v\}}|\nabla u|\,d\HH^1$ is the \emph{flux} of the
torsion function across the level set.  When $-\Delta u = 1$ in $\{u>v\}$ with
$u=v$ on $\partial\{u>v\}$ (for a.e.\ regular $v$), integration of
$-\Delta u = 1$ over $\{u>v\}$ and Green's first identity give
$\int_{\{u=v\}}|\nabla u|\,d\HH^1 = |\{u>v\}| = \mu(v)$.  Hence
\eqref{eq:cauchy-schwarz-flux} becomes
$\Per(\{u>v\})^2\le \mu(v)\,(-\mu'(v))$ for a.e.\ $v$; the deficit from
equality is $D_H(v) := \mu(v)(-\mu'(v)) - \Per(\{u>v\})^2 \ge 0$.  This is
developed in \cref{sec:identity}.
\end{remark}

%--------------------------------------------------------------------
\subsection{Isoperimetric inequalities}
%--------------------------------------------------------------------

\begin{theorem}[Planar isoperimetric inequality]\label{thm:isoperimetric}
For every set of finite perimeter $E\subset\R^2$ with $|E|<\infty$,
\begin{equation}\label{eq:isoperimetric}
  \Per(E)^2 \ge 4\pi\,|E|,
\end{equation}
with equality if and only if $E$ is equivalent (up to a Lebesgue-null set)
to a disc $B_r(x)$.
\end{theorem}

\begin{proof}
This is \cite[Theorem~14.1]{Maggi2012}; see also
\cite[Section~5.6, Theorem~5.6]{EvansGariepy2015}.
\end{proof}

\begin{theorem}[Relative isoperimetric inequality in a disc]
\label{thm:rel-isoperimetric}
There is an absolute constant $C_{\mathrm{rel}}>0$ such that for every set of
finite perimeter $E\subset\R^2$, every $x\in\R^2$, and every $r>0$,
\begin{equation}\label{eq:rel-isoperimetric}
  \min\!\bigl\{|E\cap B_r(x)|,\,|B_r(x)\setminus E|\bigr\}
  \;\le\; C_{\mathrm{rel}}\,\Per(E;B_r(x))^2.
\end{equation}
The constant $C_{\mathrm{rel}}$ is independent of $E$, $x$, and $r$.
\end{theorem}

\begin{proof}
By the relative isoperimetric inequality for convex domains, applied to the unit
disc $B_1\subset\R^2$, there is a constant $C_{\mathrm{rel}}>0$ such that for
every set of finite perimeter $E$,
$\min\{|E\cap B_1|, |B_1\setminus E|\} \le C_{\mathrm{rel}}\,\Per(E;B_1)^2$;
see \cite[Theorem~12.37]{Maggi2012} (stated in $\R^n$, here with $n=2$,
$K=B_1$) or \cite[Theorem~5.9]{EvansGariepy2015}.  The general case at radius
$r$ follows by the substitution $\tilde{E}:=\{y/r:y\in E\}$: since
$|\tilde{E}\cap B_1| = r^{-2}|E\cap B_r(x)|$ and
$\Per(\tilde{E};B_1) = r^{-1}\Per(E;B_r(x))$, the unit-disc inequality for
$\tilde{E}$ multiplied by $r^2$ yields \eqref{eq:rel-isoperimetric} with the
same constant $C_{\mathrm{rel}}$.
\end{proof}

%--------------------------------------------------------------------
\subsection{The sharp quantitative isoperimetric inequality}
%--------------------------------------------------------------------

For a set of finite perimeter $E\subset\R^2$ with $0<|E|<\infty$ we define the
\emph{Fraenkel asymmetry}
\begin{equation}\label{eq:fraenkel}
  \AF(E) := \min_{x_0\in\R^2}\frac{|E\triangle B_{r_E}(x_0)|}{|E|},
  \qquad r_E := \sqrt{|E|/\pi},
\end{equation}
and the \emph{isoperimetric deficit}
\begin{equation}\label{eq:iso-deficit}
  \diso(E) := \frac{\Per(E)}{2\sqrt{\pi|E|}}-1 \;\ge\; 0,
\end{equation}
where $2\sqrt{\pi|E|}$ is the perimeter of a disc of area $|E|$.  Both
quantities are dimensionless, translation- and scaling-invariant, and
nonnegative, vanishing precisely when $E$ is a disc.

\begin{theorem}[Fusco--Maggi--Pratelli]\label{thm:fmp}
There is an absolute constant $C_{\mathrm{F}}>0$ such that for every set of
finite perimeter $E\subset\R^2$ with $0<|E|<\infty$,
\begin{equation}\label{eq:fmp}
  \AF(E) \;\le\; C_{\mathrm{F}}\,\diso(E)^{1/2}.
\end{equation}
The exponent $\tfrac{1}{2}$ is optimal: it cannot be replaced by any strictly
larger power, as shown by elongated ellipses with small eccentricity.
The value of $C_{\mathrm{F}}$ is not explicit in \cite{FMP2008}; it is finite
and absolute (depending on neither $E$ nor any geometric parameter of $\Omega$).
\end{theorem}

\begin{proof}
This is Theorem~1.1 of \cite{FMP2008} (stated there in $\R^n$; the planar case
$n=2$ is the specialisation $C_{\mathrm{F}}:=C(2)$).  The theorem is proved in
\cite{FMP2008} using symmetrisation and a selection-principle argument; we use
it as a black box.
\end{proof}

\begin{remark}[Relation between $\diso$ and the additive deficit]\label{rem:deficit-forms}
The additive isoperimetric deficit
$D_I(E) := \Per(E)^2 - 4\pi|E| \ge 0$ (see \cref{def:deficits}) satisfies
\[
  \diso(E) = \sqrt{1 + \frac{D_I(E)}{4\pi|E|}}-1
  = \frac{D_I(E)/\!(4\pi|E|)}{\sqrt{1+D_I(E)/(4\pi|E|)}+1}
  \;\le\; \frac{D_I(E)}{8\pi|E|},
\]
since $\sqrt{1+s}+1\ge 2$ for $s\ge 0$.
Hence $\diso(E)^{1/2} \le (D_I(E)/(8\pi|E|))^{1/2}$, and \eqref{eq:fmp} gives
\[
  \AF(E)^2 \;\le\; C_{\mathrm{F}}^2\,\diso(E)
  \;\le\; \frac{C_{\mathrm{F}}^2}{8\pi}\,\frac{D_I(E)}{|E|}.
\]
We will apply this estimate in \cref{sec:main} to the inset $\Etil$.
\end{remark}

%--------------------------------------------------------------------
\subsection{Interior regularity of the torsion function}
%--------------------------------------------------------------------

Let $u\in H^1_0(\Omega)$ be the unique weak solution of the torsion problem
\begin{equation}\label{eq:torsion}
  -\Delta u = 1 \text{ in }\Omega, \qquad u = 0 \text{ on }\partial\Omega,
\end{equation}
i.e., $\int_\Omega\nabla u\cdot\nabla\varphi\,dx=\int_\Omega\varphi\,dx$ for
all $\varphi\in H^1_0(\Omega)$.  Existence and uniqueness follow from the
Lax--Milgram theorem; $u\ge 0$ by the weak maximum principle
\cite[Theorem~8.1]{GilbargTrudinger2001}.

\begin{theorem}[Interior smoothness and analyticity of $u$]\label{thm:torsion-reg}
Let $\Omega\subset\R^2$ be a bounded open set and let $u\in H^1_0(\Omega)$ be the
weak solution of \eqref{eq:torsion}.  Then $u\in C^\infty(\Omega)$.
Moreover, $u$ is real-analytic in $\Omega$.
\end{theorem}

\begin{proof}
Since the right-hand side of \eqref{eq:torsion} is $f\equiv 1\in C^\infty(\R^2)$
and the operator $-\Delta$ has smooth (constant) coefficients, interior Schauder
estimates bootstrap to give $u\in C^{k,\alpha}_{\mathrm{loc}}(\Omega)$ for
every $k\ge 0$ and every $\alpha\in(0,1)$; see
\cite[Theorem~6.17, Corollary~6.18]{GilbargTrudinger2001}.  Hence
$u\in C^\infty(\Omega)$.

Real-analyticity follows from the general fact that solutions of linear
elliptic equations with real-analytic coefficients and real-analytic right-hand
side are real-analytic in the interior: the Laplacian has constant (hence
analytic) coefficients and $f\equiv 1$ is analytic, so $u$ is real-analytic in
$\Omega$; this is a classical result of Morrey--Nirenberg, discussed in
\cite[Chapter~6, \S6.8 and the historical notes]{GilbargTrudinger2001}.
\end{proof}

\begin{theorem}[Interior gradient and Hessian estimates]\label{thm:interior-est}
There is an absolute constant $C>0$ with the following property.  Let
$\Omega\subset\R^2$ be a bounded open set, let $u$ satisfy \eqref{eq:torsion}
weakly, let $x\in\Omega$, and set $d:=\dist(x,\partial\Omega)$.  Then
\begin{equation}\label{eq:interior-est}
  |\nabla u(x)| \;\le\; C\!\left(\frac{\|u\|_{L^\infty(B_d(x))}}{d}+d\right),
  \qquad
  |D^2 u(x)| \;\le\; C\!\left(\frac{\|u\|_{L^\infty(B_d(x))}}{d^2}+1\right).
\end{equation}
\end{theorem}

\begin{proof}
Apply the interior Schauder estimate of \cite[Theorem~6.2]{GilbargTrudinger2001}
to $u$ on $B_d(x)$ where $-\Delta u = 1$.  The scale-invariant form of the
estimate gives
\[
  d\,\|\nabla u\|_{L^\infty(B_{d/2}(x))} + d^2\,\|D^2u\|_{L^\infty(B_{d/2}(x))}
  \;\le\; C\!\left(\|u\|_{L^\infty(B_d(x))} + d^2\right),
\]
with $C$ depending only on the dimension $n=2$.  Dividing by $d$ and $d^2$
respectively yields \eqref{eq:interior-est} (the sup over $B_{d/2}(x)$ bounds
the pointwise value at $x$).
\end{proof}

\begin{lemma}[$L^\infty$ bound]\label{lem:linfty}
Let $\Omega\subset\R^2$ be bounded and let $u$ solve \eqref{eq:torsion}.
Then $0\le u(x)\le \tfrac14(\diam\Omega)^2$ for all $x\in\Omega$.
In particular $m:=\max_{\overline\Omega}u \le \tfrac14(\diam\Omega)^2$.
\end{lemma}

\begin{proof}
The lower bound $u\ge 0$ follows from the weak maximum principle
\cite[Theorem~8.1]{GilbargTrudinger2001}.  For the upper bound, let $R=\diam\Omega$
and pick any $x_0\in\overline\Omega$; then $\Omega\subset B_R(x_0)$.  The
function $w(x):=(R^2-|x-x_0|^2)/4$ satisfies $-\Delta w=1$ in $\R^2$ and
$w\ge 0$ on $\Omega$.  Hence $h:=w-u$ satisfies $-\Delta h=0$ in $\Omega$ and
$h=w\ge 0$ on $\partial\Omega$, so by the maximum principle
\cite[Corollary~3.2]{GilbargTrudinger2001}, $h\ge 0$ in $\Omega$, giving
$u\le w\le R^2/4$.
\end{proof}

%--------------------------------------------------------------------
\subsection{Sard's theorem and regular values of the torsion function}
%--------------------------------------------------------------------

\begin{theorem}[Sard]\label{thm:sard}
Let $U\subseteq\R^d$ be open and $f\in C^k(U;\R^m)$ with
$k\ge\max\{1, d-m+1\}$.  Then the set of critical values of $f$,
\[
  f\!\left(\bigl\{x\in U: \operatorname{rank}\,Df(x) < m\bigr\}\right),
\]
has Lebesgue measure zero in $\R^m$.  In particular, if $f\in C^\infty(U)$
and $m=1$, the set $f(\{\nabla f=0\})\subset\R$ is a Lebesgue-null set.
\end{theorem}

\begin{proof}
See \cite[Theorem~A.6]{Maggi2012} for the statement in the context used here;
the original source is \cite{Sard1942}.
\end{proof}

\begin{lemma}[Almost every level of $u$ is regular]\label{lem:regular-values}
Let $\Omega\subset\R^2$ be a bounded open set and let $u$ solve
\eqref{eq:torsion}.  Then for Lebesgue-almost every $v\in(0,m)$ the level $v$
is a \emph{regular value} of $u$ in $\Omega$: that is, $\nabla u(x)\ne 0$ at
every point $x\in\{u=v\}\cap\Omega$.  For every such $v$,
\begin{enumerate}[label=(\roman*)]
  \item\label{it:rv-manifold}
    the level set $\{u=v\}\cap\Omega$ is a finite, disjoint union of smooth
    closed curves;
  \item\label{it:rv-perimeter}
    the superlevel set $\{u>v\}$ is a finite-perimeter set with
    $\Per(\{u>v\}) = \HH^1(\{u=v\})$;
  \item\label{it:rv-formula}
    the differentiation formula \eqref{eq:dist-deriv} holds at $v$, with
    $1/|\nabla u|$ finite and smooth on $\{u=v\}$.
\end{enumerate}
\end{lemma}

\begin{proof}
By \cref{thm:torsion-reg}, $u\in C^\infty(\Omega)$ (indeed real-analytic).
Applying \cref{thm:sard} with $d=2$, $m=1$: the set of critical values
$u(\{\nabla u=0\}\cap\Omega)\subset\R$ has Lebesgue measure zero.  Hence for
a.e.\ $v$, every point of $\{u=v\}\cap\Omega$ is a regular point.  For such $v$:

\ref{it:rv-manifold}: The implicit function theorem gives $\{u=v\}\cap\Omega$
as a smooth $1$-dimensional embedded submanifold of $\Omega$.  Since $\Omega$
is bounded and $u\in C(\overline\Omega)$ with $u=0$ on $\partial\Omega$, the
level set $\{u=v\}$ (for $v>0$) is compactly contained in $\Omega$ and hence
consists of finitely many smooth closed curves (a smooth $1$-manifold without
boundary, compact in $\Omega$).

\ref{it:rv-perimeter}: By \cref{thm:coarea},
$\Per(\{u>v\})=\HH^1(\{u=v\})$ for a.e.\ $v$.

\ref{it:rv-formula}: For a.e.\ $v$, \eqref{eq:dist-deriv} holds and
$|\nabla u|>0$ on $\{u=v\}$, so $1/|\nabla u|$ is finite and smooth there by
\cref{thm:torsion-reg}.
\end{proof}

\section{The level deficits and the volume profile}\label{sec:identity}

This section is the analytic heart of the paper. We package the Saint--Venant
deficit
\[
   \dT=\dT(\Om)=\frac\pi8-T(\Om)\ \ge\ 0
\]
of a normalised domain $|\Om|=\pi$ into the level sets of its torsion
function $u$. Recall (\cref{sec:prelim}) that $u\in H^1_0(\Om)\cap
C^\infty(\Om)$ solves $-\Delta u=1$ in $\Om$, $u=0$ on $\partial\Om$, that
\[
   T(\Om):=\int_\Om u\,dx=\int_\Om|\nabla u|^2\,dx,
   \qquad m:=\max_\Om u ,
\]
the second equality holding because testing $-\Delta u=1$ against
$u\in H^1_0(\Om)$ gives $\int_\Om|\nabla u|^2=\int_\Om u$. We define two
non-negative defects of each level set; their integral over all levels turns
out to be exactly $4\pi\dT$. We then use Talenti's comparison to pin the
volume profile $\mu(v):=|\{u>v\}|$, obtaining $m\le\tfrac14$, an exact
profile identity, and the pointwise bound $-\mu'\ge4\pi$. These facts feed
directly into the inset construction of \cref{sec:inset}.

Throughout, a value $v\in(0,m)$ is called \emph{regular} if $\nabla u\ne0$ on
$\{u=v\}$; by \cref{lem:regular-values} almost every $v\in(0,m)$ is regular,
and for such $v$ the level set $\{u=v\}$ is a (closed, embedded)
$C^\infty$ curve in $\Om$ and the coarea quotient $1/|\nabla u|$ is finite and
smooth on it. We write
\[
   E_v:=\{x\in\Om:u(x)>v\},\qquad
   \mu(v):=|E_v|,\qquad
   P(v):=P(E_v)=\HH^1(\{u=v\})
\]
for the super-level set, its area, and its perimeter (the last equality
holding for regular $v$, the level set being a smooth curve).

\subsection{The flux and coarea identities}

The two structural identities below are the only place where the equation
$-\Delta u=1$ enters this section; everything afterwards is geometry of the
level sets. The first (flux) identity quantifies the source; the second
(coarea) is the differentiation formula for $\mu$. We prove both.

\begin{lemma}[Flux and coarea identities]\label{lem:flux}
For a.e.\ $v\in(0,m)$,
\begin{equation}\label{eq:flux}
   \int_{\{u=v\}}|\nabla u|\,d\HH^1=\mu(v),
\end{equation}
and for a.e.\ $v\in(0,m)$ the function $\mu$ is differentiable with
\begin{equation}\label{eq:coarea}
   -\mu'(v)=\int_{\{u=v\}}\frac{1}{|\nabla u|}\,d\HH^1 .
\end{equation}
Moreover $\mu$ is absolutely continuous and nonincreasing on $(0,m)$, with
$\mu(0^+)=|\Om|=\pi$ and $\mu(m^-)=0$.
\end{lemma}

\begin{proof}
\emph{Flux \eqref{eq:flux}.} We derive this directly from the weak
formulation, so that no boundary regularity of $\Om$ is needed. Recall
$u\in H^1_0(\Om)$ solves
\begin{equation}\label{eq:weak}
   \int_\Om\nabla u\cdot\nabla\varphi\,dx=\int_\Om\varphi\,dx
   \qquad\text{for all }\varphi\in H^1_0(\Om).
\end{equation}
Fix $0<v<m$ and, for small $h>0$ with $v+h<m$, take the Lipschitz cut-off
\[
   \psi_h(x):=\min\Bigl\{1,\ \tfrac1h\bigl(u(x)-v\bigr)_+\Bigr\}\in[0,1].
\]
Since $t\mapsto\min\{1,\tfrac1h(t-v)_+\}$ is Lipschitz and vanishes for
$t\le v$ (in particular at $t=0$), and $u\in H^1_0(\Om)$, the composition
$\psi_h$ lies in $H^1_0(\Om)$ and is an admissible test function in
\eqref{eq:weak}, with
\[
   \nabla\psi_h=\tfrac1h\,\nabla u\,\mathbf 1_{\{v<u<v+h\}}\quad\text{a.e.}
\]
Inserting $\varphi=\psi_h$ into \eqref{eq:weak},
\[
   \frac1h\int_{\{v<u<v+h\}}|\nabla u|^2\,dx
   =\int_\Om\nabla u\cdot\nabla\psi_h\,dx
   =\int_\Om\psi_h\,dx .
\]
We pass to the limit $h\downarrow0$ on each side. On the left, the coarea
formula (\cref{thm:coarea}) gives
\[
   \frac1h\int_{\{v<u<v+h\}}|\nabla u|^2\,dx
   =\frac1h\int_v^{v+h}\Bigl(\int_{\{u=t\}}|\nabla u|\,d\HH^1\Bigr)dt
   \ \xrightarrow[h\downarrow0]{}\ \int_{\{u=v\}}|\nabla u|\,d\HH^1
\]
for a.e.\ $v$, namely at every Lebesgue point of the $L^1(0,m)$ function
$t\mapsto F(t):=\int_{\{u=t\}}|\nabla u|\,d\HH^1$ (that $F\in L^1$ follows from
$\int_0^m F\,dt=\int_{\{0<u<m\}}|\nabla u|^2\,dx\le\int_\Om|\nabla u|^2<\infty$
by coarea). On the right, $0\le\psi_h\le1$ and $\psi_h\uparrow\mathbf
1_{\{u>v\}}$ pointwise as $h\downarrow0$, so by monotone (or dominated, since
$|\Om|<\infty$) convergence
\[
   \int_\Om\psi_h\,dx\ \xrightarrow[h\downarrow0]{}\ \int_\Om\mathbf
   1_{\{u>v\}}\,dx=|\{u>v\}|=\mu(v).
\]
Equating the two limits yields \eqref{eq:flux} for a.e.\ $v\in(0,m)$.

\emph{Coarea \eqref{eq:coarea} and absolute continuity.} This is the
differentiation of the distribution function, \cref{lem:dist-deriv}: for
a.e.\ $v$, $\mu$ is differentiable with $-\mu'(v)=\int_{\{u=v\}}|\nabla
u|^{-1}\,d\HH^1$. We record the absolute-continuity statement explicitly, as
it is used repeatedly. Set $\Phi(v):=\int_{\{u=v\}}|\nabla u|^{-1}\,d\HH^1\ge0$
for those $v$ where this is finite. The coarea formula (\cref{thm:coarea})
with $g=\mathbf 1_{\{u>0\}}|\nabla u|^{-1}\mathbf 1_{\{\nabla u\ne0\}}$ gives
\[
   \int_0^m\Phi(\tau)\,d\tau
   =\int_{\{0<u<m\}\cap\{\nabla u\ne0\}}1\,dx
   =|\Om|=\pi<\infty ,
\]
where the middle equality uses that $\{\nabla u=0\}\cap\{0<u<m\}$ has
Lebesgue measure zero (it is contained in the union of level sets at critical
values, a set of measure zero by Sard's theorem,
\cref{thm:sard,lem:regular-values}, intersected with $\Om$). Hence $\Phi\in
L^1(0,m)$. For $0<v_1<v_2<m$ the coarea formula applied with
$g=\mathbf 1_{\{v_1<u\le v_2\}}|\nabla u|^{-1}\mathbf 1_{\{\nabla u\ne0\}}$
yields
\[
   \mu(v_1)-\mu(v_2)=|\{v_1<u\le v_2\}|=\int_{v_1}^{v_2}\Phi(\tau)\,d\tau .
\]
Thus $\mu(v)=\int_v^m\Phi(\tau)\,d\tau$ for $v\in(0,m)$; being the tail
integral of an $L^1$ function, $\mu$ is absolutely continuous and
nonincreasing on $(0,m)$, with $\mu'=-\Phi$ a.e., which is \eqref{eq:coarea}.
Finally $\mu(0^+)=\int_0^m\Phi=|\Om|=\pi$ (equivalently $|\{u>0\}|=|\Om|$
since $u>0$ in $\Om$ by the strong minimum principle), and
$\mu(m^-)=\int_m^m\Phi=0$.
\end{proof}

\subsection{The two level deficits}

\begin{definition}[The two level deficits]\label{def:deficits}
For a regular value $v\in(0,m)$ we define the \emph{Cauchy--Schwarz} (or
H\"older/Serrin) defect and the \emph{isoperimetric} defect of the level set
$\{u=v\}$ by
\[
   D_H(v):=\mu(v)\,\bigl(-\mu'(v)\bigr)-P(v)^2 ,
   \qquad
   D_I(v):=P(v)^2-4\pi\,\mu(v).
\]
\end{definition}

The names anticipate \cref{lem:nonneg}: $D_H$ vanishes precisely when the
flux is equidistributed (constant $|\nabla u|$ on the level curve), and $D_I$
vanishes precisely when the level set is a disk.

\begin{lemma}[Non-negativity of the deficits]\label{lem:nonneg}
For a.e.\ $v\in(0,m)$ one has
\[
   D_H(v)\ \ge\ 0,\qquad D_I(v)\ \ge\ 0 .
\]
Moreover $D_I(v)=0$ if and only if $E_v$ is (equivalent to) a disk, and, for
regular $v$, $D_H(v)=0$ if and only if $|\nabla u|$ is constant
$\HH^1$-almost everywhere on $\{u=v\}$.
\end{lemma}

\begin{proof}
\emph{$D_I\ge0$.} This is the planar isoperimetric inequality
(\cref{thm:isoperimetric}) applied to the finite-perimeter set $E_v$:
$P(E_v)^2\ge4\pi|E_v|$, i.e.\ $P(v)^2\ge4\pi\mu(v)$, with equality if and
only if $E_v$ is a disk up to a Lebesgue-null modification.

\emph{$D_H\ge0$.} Fix a regular value $v$ at which \eqref{eq:flux} and
\eqref{eq:coarea} hold (a.e.\ $v$, by \cref{lem:flux,lem:regular-values}); then
$\Sigma=\{u=v\}$ is a smooth embedded curve on which $|\nabla u|>0$ and
$P(v)=\HH^1(\Sigma)<\infty$. Apply the Cauchy--Schwarz inequality
on the finite measure $\HH^1\llcorner\Sigma$ to the functions
$|\nabla u|^{1/2}$ and $|\nabla u|^{-1/2}$:
\[
   P(v)^2
   =\Bigl(\int_{\Sigma}1\,d\HH^1\Bigr)^2
   =\Bigl(\int_{\Sigma}|\nabla u|^{1/2}\cdot|\nabla u|^{-1/2}\,d\HH^1\Bigr)^2
   \le\Bigl(\int_{\Sigma}|\nabla u|\,d\HH^1\Bigr)
      \Bigl(\int_{\Sigma}\frac{1}{|\nabla u|}\,d\HH^1\Bigr).
\]
By \eqref{eq:flux} the first factor is $\mu(v)$, and by \eqref{eq:coarea} the
second factor is $-\mu'(v)$. Hence $P(v)^2\le\mu(v)(-\mu'(v))$, i.e.\
$D_H(v)\ge0$. Equality in
Cauchy--Schwarz holds if and only if $|\nabla u|^{1/2}$ and
$|\nabla u|^{-1/2}$ are proportional $\HH^1$-a.e.\ on $\Sigma$, i.e.\ if and
only if $|\nabla u|$ is $\HH^1$-a.e.\ constant on $\{u=v\}$.
\end{proof}

\subsection{The level-set deficit identity}

We now prove the central identity and, as the mandate requires, derive the
constant $4\pi$ independently of any prior computation, directly from coarea
and integration by parts (the Saint--Venant route).

\begin{theorem}[Level-set deficit identity]\label{thm:level-identity}
With the notation above,
\begin{equation}\label{eq:level-identity}
   \int_0^m\bigl(D_H(v)+D_I(v)\bigr)\,dv\ =\ 4\pi\,\dT(\Om).
\end{equation}
\end{theorem}

\begin{proof}
The proof has three steps: a pointwise telescoping, two integrations, and
the assembly with the explicit value of the constant.

\emph{Step 1: telescoping.} For a.e.\ $v\in(0,m)$ the squared-perimeter terms
in \cref{def:deficits} cancel:
\begin{equation}\label{eq:DHDI}
   D_H(v)+D_I(v)
   =\Bigl(\mu(v)(-\mu'(v))-P(v)^2\Bigr)+\Bigl(P(v)^2-4\pi\mu(v)\Bigr)
   =\mu(v)\bigl(-\mu'(v)\bigr)-4\pi\,\mu(v).
\end{equation}
The perimeter has disappeared entirely; what remains involves only $\mu$ and
$\mu'$.

\emph{Step 2: the first term, by integration by parts.} By \cref{lem:flux},
$\mu$ is absolutely continuous on $(0,m)$ with $\mu'\in L^1$, $\mu(0^+)=\pi$
and $\mu(m^-)=0$. Consequently $v\mapsto\tfrac12\mu(v)^2$ is absolutely
continuous with derivative $\mu(v)\mu'(v)\in L^1(0,m)$ (a product of a
bounded absolutely continuous function and an $L^1$ function, with the
Leibniz rule valid because $\mu$ is bounded and absolutely continuous).
Hence
\begin{equation}\label{eq:firstterm}
   \int_0^m\mu(v)\bigl(-\mu'(v)\bigr)\,dv
   =-\Bigl[\tfrac12\mu(v)^2\Bigr]_{v=0^+}^{v=m^-}
   =\tfrac12\mu(0^+)^2-\tfrac12\mu(m^-)^2
   =\tfrac12|\Om|^2=\frac{\pi^2}{2}.
\end{equation}

\emph{Step 3: the second term, by the layer-cake formula.} Since $u>0$ in
$\Om$ and $\max u=m$,
\begin{equation}\label{eq:secondterm}
   \int_0^m 4\pi\,\mu(v)\,dv
   =4\pi\int_0^\infty|\{u>v\}|\,dv
   =4\pi\int_\Om u\,dx=4\pi\,T(\Om),
\end{equation}
the middle equality being Cavalieri's (layer-cake) formula for the
nonnegative function $u$, valid since $\mu(v)=0$ for $v\ge m$.

\emph{Assembly and the constant.} Integrating \eqref{eq:DHDI} over $(0,m)$
and inserting \eqref{eq:firstterm}--\eqref{eq:secondterm},
\[
   \int_0^m\bigl(D_H+D_I\bigr)\,dv
   =\frac{\pi^2}{2}-4\pi\,T(\Om).
\]
It remains to identify $\frac{\pi^2}{2}-4\pi T(\Om)$ with $4\pi\dT$. We do
\emph{not} invoke any earlier formula for $T(B_1)$; instead we read off the
constant from Step 2. Indeed
\[
   \frac{\pi^2}{2}=\frac{|\Om|^2}{2}=4\pi\cdot\frac{|\Om|^2}{8\pi}
   =4\pi\cdot\frac{\pi^2}{8\pi}=4\pi\cdot\frac{\pi}{8},
\]
so that
\[
   \int_0^m\bigl(D_H+D_I\bigr)\,dv
   =4\pi\Bigl(\frac{\pi}{8}-T(\Om)\Bigr)=4\pi\,\dT(\Om),
\]
by the definition $\dT=\frac\pi8-T(\Om)$. This proves
\eqref{eq:level-identity}.

We pause to make the provenance of the constant transparent, since it is the
keystone of the paper. The factor $4\pi$ in $D_I$ is the planar isoperimetric
constant ($P^2\ge4\pi A$). The factor $\tfrac12|\Om|^2$ in Step 2 is purely
the integration of $\mu(-\mu')$, owing nothing to the equation beyond the
boundary values $\mu(0^+)=|\Om|$, $\mu(m^-)=0$. These two independent inputs
combine into the deficit because of the elementary algebraic identity
$\tfrac12 A^2=4\pi\cdot A^2/(8\pi)$ for any area $A$; with $A=|\Om|=\pi$ this
reads $\tfrac12|\Om|^2=4\pi\cdot\tfrac\pi8$, which is the algebra displayed
above. The number $\tfrac\pi8$ so produced is exactly the disk torsion
$T(B_1)$ (computed independently in \cref{lem:profile} below), so the
right-hand side is $4\pi\bigl(T(B_1)-T(\Om)\bigr)$. Thus
the right-hand side is $4\pi(T(B_1)-T(\Om))$, the Saint--Venant deficit
scaled by $4\pi$, with no hidden normalisation.
\end{proof}

\begin{remark}[The identity as an exact, scale-aware statement]\label{rem:identity}
Identity \eqref{eq:level-identity} expresses the (global) Saint--Venant
deficit as the total over all levels of two pointwise non-negative geometric
defects of the level sets: $D_H$, the failure of $|\nabla u|$ to be constant
on $\{u=v\}$; and $D_I$, the failure of $E_v$ to be a disk. Both vanish
identically (for all $v$) if and only if every level set is a disk on which
$|\nabla u|$ is constant, i.e.\ if and only if $\Om$ is a disk and $u=u_B$;
this is Serrin's overdetermined characterisation, and it is the rigidity
behind \eqref{eq:level-identity}. The quantitative content is the converse,
exploited from \cref{sec:inset} on: small $\dT$ forces both defects to be
small \emph{on most levels}.
\end{remark}

\subsection{Scale invariance of the identity}

For the inset of \cref{sec:inset} we will apply the identity not to $\Om$ but
to its super-level sets, which are not area-normalised. We therefore record
the unnormalised form of \eqref{eq:level-identity} and its exact
homogeneity. Let $E\subset\R^2$ be any bounded open set of finite measure
which is regular for the Dirichlet problem, with torsion function $w_E$
($-\Delta w_E=1$ in $E$, $w_E\in H^1_0(E)$), torsion
$T(E)=\int_E w_E=\int_E|\nabla w_E|^2$, maximum $m_E=\max w_E$, and level
quantities $\mu_E(v)=|\{w_E>v\}|$, $P_E(v)=P(\{w_E>v\})$,
\[
   D_H^E(v):=\mu_E(v)\bigl(-\mu_E'(v)\bigr)-P_E(v)^2,
   \qquad
   D_I^E(v):=P_E(v)^2-4\pi\mu_E(v).
\]

\begin{proposition}[Unnormalised identity and degree-$4$ homogeneity]
\label{prop:scale}
For any such $E$,
\begin{equation}\label{eq:identity-gen}
   \int_0^{m_E}\bigl(D_H^E(v)+D_I^E(v)\bigr)\,dv
   =\tfrac12|E|^2-4\pi\,T(E)
   =4\pi\bigl(T(B_E)-T(E)\bigr),
   \qquad T(B_E):=\frac{|E|^2}{8\pi},
\end{equation}
where $B_E$ is any disk with $|B_E|=|E|$. Both sides are homogeneous of
degree $4$ under dilations: writing $sE=\{sx:x\in E\}$ for $s>0$, one has
\begin{equation}\label{eq:homog}
   \int_0^{m_{sE}}\bigl(D_H^{sE}+D_I^{sE}\bigr)\,dv
   =s^4\int_0^{m_E}\bigl(D_H^{E}+D_I^{E}\bigr)\,dv,
   \qquad
   T(B_{sE})-T(sE)=s^4\bigl(T(B_E)-T(E)\bigr).
\end{equation}
Consequently the \emph{scale-invariant} deficit
$\dstar(E):=\bigl(T(B_E)-T(E)\bigr)/T(B_E)\in[0,1)$ satisfies
\begin{equation}\label{eq:scaleinv-identity}
   \dstar(E)
   =\frac{T(B_E)-T(E)}{T(B_E)}
   =\frac{\frac{1}{4\pi}\int_0^{m_E}(D_H^E+D_I^E)\,dv}{|E|^2/(8\pi)}
   =\frac{2}{|E|^2}\int_0^{m_E}\bigl(D_H^E+D_I^E\bigr)\,dv ,
\end{equation}
and $\dstar(sE)=\dstar(E)$ for all $s>0$.
\end{proposition}

\begin{proof}
\emph{Identity \eqref{eq:identity-gen}.} The proof of
\cref{lem:flux,thm:level-identity} used only $-\Delta w_E=1$, $w_E\in
H^1_0(E)$, interior smoothness, and Sard's theorem, all of which hold for $E$
verbatim (interior regularity and analyticity of the torsion function depend
only on the interior equation, not on boundary regularity;
\cref{thm:torsion-reg}). Repeating the three steps gives
\[
   \int_0^{m_E}(D_H^E+D_I^E)\,dv
   =\tfrac12\mu_E(0^+)^2-4\pi\,T(E)
   =\tfrac12|E|^2-4\pi\,T(E),
\]
since $\mu_E(0^+)=|\{w_E>0\}|=|E|$ ($w_E>0$ in $E$). The last equality in
\eqref{eq:identity-gen} is the algebraic identity
$\tfrac12|E|^2=4\pi\cdot|E|^2/(8\pi)=4\pi\,T(B_E)$.

\emph{Homogeneity \eqref{eq:homog}.} The torsion function transforms under
dilation by $w_{sE}(x)=s^2\,w_E(x/s)$: indeed $w_{sE}$ so defined vanishes on
$\partial(sE)$ and $-\Delta w_{sE}(x)=-s^2\cdot s^{-2}(\Delta w_E)(x/s)=1$,
and by uniqueness it is the torsion function of $sE$. Hence
$m_{sE}=s^2 m_E$, and
\[
   T(sE)=\int_{sE}w_{sE}\,dx=\int_{sE}s^2 w_E(x/s)\,dx
   \overset{x=sy}{=}s^2\!\int_E s^2 w_E(y)\,dy=s^4\,T(E),
\]
while $|sE|^2=s^4|E|^2$; both terms on the right of \eqref{eq:identity-gen}
scale by $s^4$, giving the second equality in \eqref{eq:homog}, and the first
follows from it via \eqref{eq:identity-gen}. (Directly: for a level $v=s^2v'$
one has $\{w_{sE}>v\}=s\{w_E>v'\}$, so $\mu_{sE}(s^2v')=s^2\mu_E(v')$,
$P_{sE}(s^2v')=s\,P_E(v')$, $\mu_{sE}'(s^2v')=\mu_E'(v')$, whence
$D_H^{sE}(s^2v')=s^2 D_H^E(v')$ and $D_I^{sE}(s^2v')=s^2 D_I^E(v')$; the
change of variables $v=s^2v'$, $dv=s^2\,dv'$ then produces the factor
$s^2\cdot s^2=s^4$.)

\emph{Scale invariance of $\dstar$.} From \eqref{eq:identity-gen},
$T(B_E)-T(E)=\tfrac{1}{4\pi}\int_0^{m_E}(D_H^E+D_I^E)\,dv$ and
$T(B_E)=|E|^2/(8\pi)$, so dividing gives
$\dstar(E)=\tfrac{2}{|E|^2}\int_0^{m_E}(D_H^E+D_I^E)\,dv$, which is
\eqref{eq:scaleinv-identity}. By \eqref{eq:homog} numerator and denominator
$T(B_E)$ both scale by $s^4$, so $\dstar(sE)=\dstar(E)$.
\end{proof}

\begin{remark}[Consistency with the normalisation]\label{rem:consistency}
Applied to $E=\Om$ with $|\Om|=\pi$, \eqref{eq:identity-gen} returns
\eqref{eq:level-identity}: $\tfrac12\pi^2-4\pi T(\Om)=4\pi\dT$, and
$T(B_\Om)=\pi^2/(8\pi)=\pi/8=T(B_1)$. The two normalisations are related by
$\dstar(\Om)=\dT/T(B_1)=\tfrac8\pi\dT$, in agreement with the global
notation.
\end{remark}

\subsection{Talenti's comparison and the volume profile}

To convert ``small on most levels'' into quantitative geometry we need the
size of the level sets, i.e.\ the profile $\mu(v)$. This is furnished by
Talenti's pointwise comparison with the model disk. Recall the unit-disk
torsion function $u_{B_1}(x)=\tfrac14(1-|x|^2)$ and its distribution function
\[
   \mu_B(v):=|\{u_{B_1}>v\}|=\pi(1-4v)_+ ,
   \qquad v\ge0,
\]
where $t_+=\max\{t,0\}$.

\begin{theorem}[Talenti comparison for the profile]\label{thm:talenti}
Let $u$ be the torsion function of $\Om$ with $|\Om|=\pi$. Then
\begin{equation}\label{eq:talenti}
   \mu(v)=|\{u>v\}|\ \le\ \mu_B(v)=\pi(1-4v)_+
   \qquad\text{for all }v\ge0 .
\end{equation}
\end{theorem}

\begin{proof}
We invoke Talenti's comparison theorem
\cite[Theorem~1]{Talenti1976} in the form proved there: if
$u\in H^1_0(\Om)$ solves $-\Delta u=f$ with $f\ge0$, and $u^\ast$ denotes the
Schwarz (radially symmetric decreasing) rearrangement of $u$, while
$\tilde u$ solves $-\Delta\tilde u=f^\ast$ in the ball $\Om^\ast$ centred at
the origin with $|\Om^\ast|=|\Om|$ and $\tilde u\in H^1_0(\Om^\ast)$, then
\[
   u^\ast(x)\ \le\ \tilde u(x)\qquad\text{for all }x\in\Om^\ast .
\]
We apply this with $f\equiv1$. Then $f^\ast\equiv1$, $\Om^\ast=B_1$ (since
$|\Om|=\pi=|B_1|$), and the comparison solution is the disk torsion function
$\tilde u=u_{B_1}$, because $u_{B_1}=\tfrac14(1-|x|^2)$ solves
$-\Delta u_{B_1}=1$ in $B_1$ with zero boundary data and is its own
rearrangement. Thus $u^\ast\le u_{B_1}$ pointwise on $B_1$.

Schwarz rearrangement is equimeasurable: $|\{u>v\}|=|\{u^\ast>v\}|$ for every
$v$. Since $u^\ast\le u_{B_1}$, we have
$\{u^\ast>v\}\subseteq\{u_{B_1}>v\}$, hence
$|\{u^\ast>v\}|\le|\{u_{B_1}>v\}|$. Combining,
\[
   \mu(v)=|\{u>v\}|=|\{u^\ast>v\}|\le|\{u_{B_1}>v\}|=\mu_B(v)
   =\pi(1-4v)_+ ,
\]
where the last equality is $\{u_{B_1}>v\}=\{|x|^2<1-4v\}$, of area
$\pi(1-4v)_+$.
\end{proof}

\begin{lemma}[Exact volume profile and the cap on the maximum]\label{lem:profile}
The Saint--Venant deficit equals the $L^1$ gap between the disk profile and
the domain profile over $(0,\tfrac14)$:
\begin{equation}\label{eq:profile}
   \int_0^{1/4}\bigl(\mu_B(v)-\mu(v)\bigr)\,dv\ =\ \dT(\Om),
\end{equation}
the integrand $\mu_B(v)-\mu(v)\ge0$ being nonnegative for every $v$. In
particular the maximum is capped,
\begin{equation}\label{eq:mcap}
   m=\max_\Om u\ \le\ \tfrac14 .
\end{equation}
\end{lemma}

\begin{proof}
\emph{The maximum cap \eqref{eq:mcap}.} Suppose, for contradiction, that
$m>\tfrac14$. Pick $v$ with $\tfrac14<v<m$. Then $\{u>v\}\ne\emptyset$ and is
open (as $u$ is continuous), so $\mu(v)=|\{u>v\}|>0$; but $v>\tfrac14$ gives
$\mu_B(v)=\pi(1-4v)_+=0$, contradicting \eqref{eq:talenti}. Hence
$m\le\tfrac14$. In particular $\mu(v)=0$ for $v\ge\tfrac14\ge m$.

\emph{The profile identity \eqref{eq:profile}.} Nonnegativity of the
integrand on all of $(0,\tfrac14)$ is exactly \eqref{eq:talenti}. By the
layer-cake formula and $m\le\tfrac14$,
\[
   T(\Om)=\int_\Om u\,dx=\int_0^\infty\mu(v)\,dv
   =\int_0^{1/4}\mu(v)\,dv ,
\]
the last step using $\mu(v)=0$ for $v\ge\tfrac14$. Likewise, by the same
formula for the disk torsion function $u_{B_1}$,
\[
   T(B_1)=\int_{B_1}u_{B_1}\,dx=\int_0^\infty\mu_B(v)\,dv
   =\int_0^{1/4}\pi(1-4v)\,dv
   =\pi\Bigl[v-2v^2\Bigr]_0^{1/4}
   =\pi\Bigl(\tfrac14-\tfrac18\Bigr)=\frac\pi8 .
\]
Subtracting,
\[
   \int_0^{1/4}\bigl(\mu_B(v)-\mu(v)\bigr)\,dv
   =T(B_1)-T(\Om)=\frac\pi8-T(\Om)=\dT(\Om).
\]
This is \eqref{eq:profile}. (The computation also reconfirms
independently that $T(B_1)=\pi/8$, hence that the deficit is well defined and
non-negative: $\dT=\int_0^{1/4}(\mu_B-\mu)\ge0$.)
\end{proof}

\begin{remark}[Two complementary decompositions of $\dT$]\label{rem:two}
We have written $\dT$ in two ways: as an integral over levels of the
\emph{shape} defects $D_H+D_I$ (\cref{thm:level-identity}, weight $4\pi$), and
as an integral over levels of the \emph{volume} gap $\mu_B-\mu$
(\cref{lem:profile}, weight $1$). The inset construction in \cref{sec:inset}
uses both: the first transfers the torsion deficit to a single level set, the
second certifies that this level set captures almost all of the area of
$\Om$.
\end{remark}

\subsection{A pointwise lower bound on \texorpdfstring{$-\mu'$}{-mu'}}

The final consequence we extract is a clean pointwise estimate, used to
quantify how fast the volume profile decreases.

\begin{lemma}[Lower bound on the profile slope]\label{lem:muprime}
For a.e.\ $v\in(0,m)$,
\begin{equation}\label{eq:muprime}
   -\mu'(v)\ \ge\ 4\pi .
\end{equation}
\end{lemma}

\begin{proof}
Fix $v\in(0,m)$ a regular value at which \eqref{eq:coarea} holds (a.e.\ $v$).
By \cref{lem:nonneg}, $D_H(v)\ge0$ and $D_I(v)\ge0$, i.e.
\[
   \mu(v)\bigl(-\mu'(v)\bigr)\ \ge\ P(v)^2\qquad\text{and}\qquad
   P(v)^2\ \ge\ 4\pi\,\mu(v).
\]
Chaining these,
\[
   \mu(v)\bigl(-\mu'(v)\bigr)\ \ge\ P(v)^2\ \ge\ 4\pi\,\mu(v).
\]
Since $v<m$, the super-level set $E_v=\{u>v\}$ is nonempty and open, so
$\mu(v)>0$; dividing by $\mu(v)>0$ gives $-\mu'(v)\ge4\pi$.
\end{proof}

\begin{remark}[Sharpness and the model]\label{rem:sharp-mu}
All four statements of this section are sharp and are simultaneously
saturated by the disk: for $u_{B_1}$ one has $\mu=\mu_B$, $-\mu_B'(v)=4\pi$ on
$(0,\tfrac14)$, $D_H\equiv D_I\equiv0$, $\dT=0$, and $m=\tfrac14$. Thus
\eqref{eq:muprime} is an equality for the model, and the level identity
\eqref{eq:level-identity} reads $0=0$. The quantitative theory of the
remaining sections measures the deviation from this rigid configuration.
\end{remark}

\section{The inset: regularity, deficit transfer, level selection}
\label{sec:inset}

Throughout this section $\Om\subset\R^2$ is open with $|\Om|=\pi$, $u\in
H^1_0(\Om)$ is its torsion function ($-\Delta u=1$ in $\Om$, $u=0$ on
$\partial\Om$), $m:=\max_\Om u$, and $\dT=\dT(\Om)=\tfrac\pi8-T(\Om)\ge0$ is
the Saint--Venant deficit. We write $E_v:=\{u>v\}$ for the open super-level
set, $\mu(v):=|E_v|$, $\Per(v):=\Per(E_v)$, and recall the two level
deficits of \cref{def:deficits},
\[
   D_H(v)=\mu(v)\bigl(-\mu'(v)\bigr)-\Per(v)^2\ \ge0,
   \qquad
   D_I(v)=\Per(v)^2-4\pi\mu(v)\ \ge0
\]
(\cref{lem:nonneg}), the level-set deficit identity
$\int_0^m(D_H+D_I)\,dv=4\pi\dT$ (\cref{thm:level-identity}), the Talenti
profile $\mu_B(v):=\pi(1-4v)_+$ with $\mu\le\mu_B$ and the profile identity
$\int_0^{1/4}(\mu_B-\mu)\,dv=\dT$ (\cref{thm:talenti}, \cref{lem:profile}),
the bound $-\mu'(v)\ge4\pi$ a.e.\ (\cref{lem:muprime}), and the flux/coarea
identities $\int_{\{u=v\}}|\nabla u|\,d\HH^1=\mu(v)$,
$-\mu'(v)=\int_{\{u=v\}}|\nabla u|^{-1}\,d\HH^1$ for a.e.\ regular $v$
(\cref{thm:coarea}, \cref{lem:dist-deriv}). By Sard's theorem
(\cref{thm:sard}) and \cref{lem:regular-values} the non-regular values of
$u$ form a Lebesgue-null set, and for every regular $v\in(0,m)$ the set
$E_v$ has finite perimeter with $\HH^1$-rectifiable boundary
$\{u=v\}$ (\cref{thm:torsion-reg}). Absolute constants are denoted $C$ (value
may change) or named ($A_0$, $\delta_\star$); $\delta_\star>0$ is a small
absolute constant, fixed at the end of the section. We set $\tau:=\sqrt\dT$
only at the very end; until then $\tau$ is a free parameter in $(0,\tfrac1{16})$.

\subsection{Standing measure-theoretic convention and interior containment}
\label{ss:containment}

The torsion function $u$ is regular only in the \emph{interior} of $\Om$: by
\cref{thm:torsion-reg}, $u\in C^\infty(\Om)$ (indeed real-analytic) and
$u>0$ in $\Om$, but for a general bounded open $\Om$ we do \emph{not} assume
$u\in C(\overline\Om)$, nor that $u=0$ holds pointwise on $\partial\Om$, nor
that the super-level sets $E_v=\{u>v\}$ are compactly contained in $\Om$.
Indeed they need not be: for the punctured disc $\Om=B_1\setminus\set0$
(area $\pi$) the torsion function is $u=u_{B_1}=(1-|x|^2)/4$ (the removed
point $\set0$ is polar, hence invisible to $H^1_0$), and every super-level
set $\set{u\ge t}=\overline{B_{\sqrt{1-4t}}}$ accumulates at the boundary
point $0\in\partial\Om$, so its closure is \emph{not} a compact subset of
$\Om$. The naive statement ``$\overline{E_v}$ is compactly contained in
$\Om$'' is therefore false in general, and we adopt throughout the only
correct, measure-theoretic, stance: all quantities of the paper ---
$|\Om|$, $T(\Om)$, the asymmetry $\AF$, the level sets $E_v$, their
measures $\mu(v)$ and perimeters $\Per(v)$ --- depend on $\Om$ only through
its measure-theoretic structure, and a \emph{polar} (capacity-zero) subset
of $\partial\Om$ may be added to or removed from $\Om$ without changing any
of them. We make this precise and isolate the containment property we
actually use.

\begin{lemma}[Interior containment of super-level sets, modulo a polar set]
\label{lem:containment}
Let $\Om\subset\R^2$ be bounded open with torsion function $u$, and let
$t>0$. Then:
\begin{enumerate}[label=\textup{(\roman*)},leftmargin=2.4em]
\item\label{it:cont-fp} for a.e.\ $t>0$ the set $E_t=\set{u>t}$ has finite
perimeter, $\Per(E_t)=\HH^1(\set{u=t})$, and $E_t$ is an open subset of
$\Om$ that coincides with its own measure-theoretic interior up to an
$\HH^1$-null set;
\item\label{it:cont-far} for every regular $t>0$ and every connected
component $E$ of $E_t$, the relative boundary of $E$ in $\Om$ lies in the
level set,
\[
   \partial E\cap\Om\ \subset\ \set{u=t}\ \subset\ \Om,
\]
an $\HH^1$-rectifiable curve contained in the \emph{open} set $\Om$ and
carrying the full perimeter, $\Per(E)=\HH^1(\partial^*E)=\HH^1(\partial E\cap
\Om)$; the set $E$ is bounded \textup{(}so $\overline E\subset\R^2$ is
compact\textup{)}, and the $\R^2$-closure $\overline E$ may meet
$\partial\Om$, but $\overline E\cap\partial\Om$ carries \emph{no} perimeter
of $E$: $\HH^1(\partial^*E\cap\partial\Om)=0$. \textup{(}This is the precise,
measure-theoretic substitute for ``$\overline E$ compactly contained in
$\Om$'', which fails for the punctured disc.\textup{)}
\item\label{it:cont-polar} consequently, replacing $\Om$ by
$\Om\cup Z$ or $\Om\setminus Z$ for any polar set $Z$ \textup{(}capacity
zero, hence Lebesgue-null and not charged by the perimeter measure of any
finite-perimeter set, \cite[Thm.~4.5, \S4.7]{EvansGariepy2015}\textup{)}
leaves $u$, $T(\Om)$, every $E_t$, $\mu(t)$, $\Per(t)$, $\dT$ and $\AF(\Om)$
unchanged. In particular the symmetric differences $|\Om\triangle E_t|$ and
$|\Om\triangle G|$ used downstream are insensitive to such $Z$.
\end{enumerate}
\end{lemma}

\begin{proof}
\emph{\ref{it:cont-fp}.} By the coarea formula (\cref{thm:coarea}) applied
to $u\in H^1(\Om)\cap C^\infty(\Om)$, $\int_0^m\Per(E_t)\,dt
=\int_\Om|\nabla u|<\infty$, so $E_t$ has finite perimeter for a.e.\ $t$,
with $\Per(E_t)=\HH^1(\set{u=t})$ at every regular $t$
(\cref{lem:regular-values}). The set $E_t$ is open as a super-level set of
the continuous function $u$ on the open set $\Om$; for a regular value
$\nabla u\neq0$ on $\set{u=t}$, so $\set{u=t}$ is an $\HH^1$-rectifiable
$C^1$ curve, of zero Lebesgue measure, whence $E_t$ equals its
measure-theoretic interior up to that null curve.

\emph{\ref{it:cont-far}.} Fix a regular $t$ and a component $E$ of $E_t$.
$E$ is open (super-level set of the continuous $u$) and bounded ($E\subset
\Om$, $\Om$ bounded), so $\overline E\subset\R^2$ is compact. The relative
boundary $\partial E\cap\Om$ consists of points $q\in\Om$ that are limits of
points of $E$ (where $u>t$) and of points of $\Om\setminus E$; by continuity
of $u$ on the open set $\Om$, $u(q)=t$, so $\partial E\cap\Om\subset
\set{u=t}$, which at the regular value $t$ is an $\HH^1$-rectifiable $C^1$
curve in $\Om$ along which $\nabla u\neq0$ (\cref{lem:regular-values}). The
reduced boundary satisfies $\partial^*E\subset\partial E$, and the part of
$\partial^*E$ inside $\Om$ is $\HH^1$-equivalent to $\set{u=t}\cap\partial E$
(\cref{thm:coarea}, $\Per(E)=\HH^1(\set{u=t}\cap\partial E)$). It remains to
see that $\overline E$ may touch $\partial\Om$ only off the reduced boundary.
Indeed $\partial^*E\subset\partial E$, and the part of $\partial E$ on
$\partial\Om$ is disjoint from $\Om\supset\set{u=t}$; since
$\partial^*E\cap\Om\subset\set{u=t}\subset\Om$ already carries the full
perimeter $\Per(E)=\HH^1(\partial^*E)$, the remainder $\partial^*E\cap
\partial\Om$ has $\HH^1$-measure zero. (For the punctured disc, $\overline E
=\overline{B_{\sqrt{1-4t}}}$ touches $\partial\Om$ at the polar puncture
$0$, which is a single point, hence $\HH^1$-null and not a reduced-boundary
point.) This proves \ref{it:cont-far}.

\emph{\ref{it:cont-polar}.} A polar set $Z\subset\R^2$ has zero
$2$-capacity, hence zero Lebesgue measure and is removable for $H^1_0$: the
torsion function of $\Om$ and of $\Om\cup Z$ coincide ($H^1_0(\Om)=
H^1_0(\Om\cup Z)$ when $Z$ is polar of finite measure, since capacity-zero
sets are negligible for $H^1$ traces, \cite[\S4.7]{EvansGariepy2015}). All
listed quantities are integrals against Lebesgue measure of $u$ or
indicators of finite-perimeter sets, or perimeters thereof; Lebesgue measure
and the perimeter measure (which is absolutely continuous with respect to
$\HH^1\restr{\partial^*E}$ and assigns no mass to the $\HH^1$-null, indeed
Lebesgue-null, polar set $Z$, \cite[\S5.7]{EvansGariepy2015},
\cite[Thm.~16.2]{Maggi2012}) both ignore $Z$. The claim follows.
\end{proof}

\begin{remark}[The role of \cref{lem:containment}]\label{rem:containment}
\cref{lem:containment} is the honest substitute for the (false in general)
assertion that $\overline{E_v}$ is compactly contained in $\Om$. What the
later sections genuinely need is only: \textup{(a)} $E_v$ and each of its
components are open subsets of $\Om$ on which $u$ is real-analytic with
$\nabla u\neq0$ along $\set{u=v}$ at regular $v$ (so $\partial E$ is an
$\HH^1$-rectifiable curve in the \emph{open} set $\Om$, with $\Per(E)=
\HH^1(\partial^*E)$); and \textup{(b)} the measure-theoretic insensitivity
to polar sets, \ref{it:cont-polar}, which guarantees that all symmetric
differences and deficits are well posed. In \cref{sec:annulus} we work with
the \emph{measure-theoretic representative} of $\Etil$ \textup{(}equivalently
the finely-open representative\textup{)}: the topological boundary
$\partial\Etil$ there is the boundary of this representative, which differs
from the reduced and essential boundaries only on an $\HH^1$-null set and
lies in the open set $\Om$. No regularity of $\partial\Om$ is used anywhere.
\end{remark}

\subsection{No interior holes}

We now record a topological regularity property of the super-level sets,
which is a direct consequence of the superharmonicity of $u$
($-\Delta u=1\ge0$) and the minimum principle. We use only the
\emph{interior} regularity of \cref{thm:torsion-reg} --- $u\in C^\infty(\Om)$,
$u>0$ in $\Om$ --- together with \cref{lem:containment}; in particular we
do \emph{not} assume $u\in C(\overline\Om)$ or any regularity of
$\partial\Om$. The phrase ``$\overline K\cap\partial\Om\neq\emptyset$''
below means that $K$, a relatively closed bounded subset of $\Om$, is
\emph{not} compactly contained in $\Om$: its closure in $\R^2$ meets
$\partial\Om$ (equivalently, $K$ has a limit point on $\partial\Om$).

\begin{lemma}[No interior holes]\label{lem:noholes}
For every $v\in(0,m)$ the following hold.
\begin{enumerate}[label=\textup{(\roman*)},leftmargin=2.4em]
\item Every connected component $K$ of the open set
$\{x\in\Om:u(x)<v\}$ has $\overline K\cap\partial\Om\neq\emptyset$.
\item For every \emph{regular} value $v$, each connected component $E$ of
$E_v=\{u>v\}$ has no interior holes: every bounded connected component $H$
of $\R^2\setminus\overline E$ satisfies $\overline H\not\subset\Om$; in
particular $\R^2\setminus\overline E$ is connected up to sets touching
$\partial\Om$, and the inset defined below is simply connected modulo
$\partial\Om$.
\end{enumerate}
\end{lemma}

\begin{proof}
The entire argument takes place in the open set $\Om$, where $u$ is smooth
(\cref{thm:torsion-reg}); no value of $u$ on $\partial\Om$ and no continuity
of $u$ up to $\partial\Om$ is invoked. The negations we contradict are
exactly the statements that the relevant components are \emph{compactly
contained in $\Om$} (closure in $\R^2$ a compact subset of $\Om$), the case
\cref{lem:containment}\ref{it:cont-far} shows cannot occur for super-level
components; here we prove the analogous exclusion for the complementary
sub-level islands and for holes.

\emph{(i).} Suppose, for contradiction, that some component $K$ of
$\{u<v\}$ satisfies $\overline K\Subset\Om$, i.e.\ $\overline K$ is a compact
subset of $\Om$ not meeting $\partial\Om$. We claim $u\equiv v$ on
$\partial K$. Let $x_0\in\partial K$. Since $\overline K\subset\Om$ we have
$x_0\in\Om$, so $u$ is continuous at $x_0$ (interior continuity,
\cref{thm:torsion-reg}; no boundary regularity is used). If $u(x_0)<v$ then by continuity
$u<v$ on a ball around $x_0$, which would be a connected subset of
$\{u<v\}$ meeting $K$, hence contained in $K$, placing $x_0$ in the
\emph{interior} of $K$ --- contradicting $x_0\in\partial K$. If $u(x_0)>v$
then by continuity $u>v$ near $x_0$, so a punctured neighbourhood of $x_0$
is disjoint from $\{u<v\}\supset K$, contradicting $x_0\in\partial K$ (as
$K$ accumulates at $x_0$). Hence $u(x_0)=v$, proving $u=v$ on $\partial K$.

Now $u$ is continuous on the compact set $\overline K$ and attains its
minimum there; since $u<v$ at any interior point of $K$ and $u=v$ on
$\partial K$, the minimum is attained at an interior point and equals
$\min_{\overline K}u<v=\min_{\partial K}u$. But $-\Delta u=1\ge0$ in $\Om$
makes $u$ superharmonic, and the (weak) minimum principle for
superharmonic functions in $C^2$
\cite[Theorem~3.1]{GilbargTrudinger2001} gives
$\min_{\overline K}u=\min_{\partial K}u$. This contradiction proves (i).

\emph{(ii).} Let $v$ be regular and $E$ a component of $E_v$, and suppose
$H$ is a bounded component of $\R^2\setminus\overline E$ with
$\overline H\subset\Om$. Then $H$ is an open subset of $\Om$ with
$\overline H\Subset\Om$, and we consider the two open sets
$W^-:=H\cap\{u<v\}$ and $W^+:=H\cap\{u>v\}$, both contained in $H$, hence
both with closures contained in $\overline H\Subset\Om$.

We first note that every connected component $G$ of $W^-$ (resp.\ of $W^+$)
is in fact a connected component of $\{u<v\}$ (resp.\ of $\{u>v\}$), and is
contained in $H$. Indeed, $G$ is connected and contained in the open set
$\{u\ne v\}$, so it lies in a unique component $\hat G$ of $\{u<v\}$ (resp.\
$\{u>v\}$). If $\hat G$ left $H$ it would, being connected and meeting $H$,
have to cross $\partial H$; but $\partial H\subset\overline E$ and on
$\overline E$ one has $u\ge v$, while points of $\partial H$ with $u>v$ lie in
the open set $E$ (a component of $E_v$, hence open) which is disjoint from
$\hat G$ when $\hat G\subset W^-$, and equals a different component of $E_v$
disjoint from $E$ when $\hat G\subset W^+$ (as $\hat G\cap E\subset W^+\cap
E\subset H\cap E=\emptyset$). In either case $\hat G$ cannot meet $\partial H$,
so $\hat G\subset H$ and $\hat G=G$. Thus each component of $W^-$ is a
connected component of $\{u<v\}$
whose closure lies in $\overline H\Subset\Om$ and so does not meet
$\partial\Om$; by (i) this forces $W^-=\emptyset$. Each connected component
$E'$ of $W^+$ is a connected component of $E_v=\{u>v\}$ with
$\overline{E'}\Subset\Om$; on it $u=v$ on $\partial E'$ and $u>v$ inside, so
$u$ attains an interior maximum strictly above its boundary value. But $u$
is superharmonic ($-\Delta u=1\ge0$), so $-u$ is subharmonic and satisfies
the maximum principle \cite[Theorem~3.1]{GilbargTrudinger2001}:
$\max_{\overline{E'}}u=\max_{\partial E'}u=v$, contradicting $u>v$ inside.
Hence $W^+=\emptyset$ as well.

Therefore $u\equiv v$ on the open set $H$, which has positive Lebesgue
measure; but $v$ is a regular value, so $\{u=v\}$ is an $\HH^1$-rectifiable
curve, of zero Lebesgue measure. This contradicts $|H|>0$, so no such $H$
exists.
\end{proof}

\begin{remark}[Measure-theoretic reading; no boundary regularity]
Part (ii) uses both the minimum principle (via (i), excluding holes that are
sub-level islands) and the maximum principle (excluding super-level islands
separated from $E$), and only interior regularity of $u$
(\cref{thm:torsion-reg}); it is consistent with the standing convention of
\cref{ss:containment}. The statement ``$\overline H\not\subset\Om$'' is read
in the sense of \cref{lem:containment}\ref{it:cont-far}: a putative hole $H$
that is compactly contained in $\Om$ is excluded, while a putative hole that
accumulates only on $\partial\Om$ (as the centre of the punctured disc would,
were $\set0$ removed) is not a genuine interior hole and carries no perimeter
of $E$. This does \emph{not} by itself prove that the full exterior
$\R^2\setminus\overline\Etil$ is connected; that stronger topological input is
supplied unconditionally by the hole-filling construction \cref{lem:fill} of
\cref{ss:fill}, which fills any boundary-touching holes at a cost quadratic in
the isoperimetric deficit. The consequences used unconditionally
here are connectedness of $\Etil$, exclusion of compactly contained holes, and
the measure-theoretic fact $|\{u=\hat v\}|=0$ used repeatedly to identify
symmetric differences. A reader content with these may keep (i) and the line
``$u\equiv v$ on $H$ with $|H|>0$ contradicts $|\{u=v\}|=0$'' as the whole of
(ii).
\end{remark}

\subsection{Deficit transfer to super-level sets}

The torsion function of any super-level set $E_v$ is $u-v$, and the
level-set deficit identity, being homogeneous of degree four, transfers the
scale-invariant deficit from $\Om$ to $E_v$ as a tail of the
positive-integrand identity of \cref{thm:level-identity}.

\begin{lemma}[Deficit transfer]\label{lem:transfer}
For every regular $v\in(0,m)$,
\begin{equation}\label{eq:transfer}
   \dstar(E_v)\ \le\ \frac{\pi^2}{\mu(v)^2}\,\dstar(\Om).
\end{equation}
\end{lemma}

\begin{proof}
On the open set $E_v$ the relevant function is
$w:=(u-v)_+\restr{E_v}=u-v$.  This is the torsion function of $E_v$ in the
weak Sobolev sense.  Indeed, by the capacity characterisation of zero traces,
the quasi-continuous representative of $(u-v)_+$ vanishes outside $E_v$ up to
capacity zero, hence $(u-v)_+\in H^1_0(E_v)$
\cite[\S4.7]{EvansGariepy2015}.  For every
$\varphi\in C_c^\infty(E_v)$,
\[
   \int_{E_v}\nabla w\cdot\nabla\varphi
   =\int_{\Om}\nabla u\cdot\nabla\varphi
   =\int_{\Om}\varphi
   =\int_{E_v}\varphi .
\]
Thus $-\Delta w=1$ in $E_v$ with zero Sobolev trace, and uniqueness in
$H^1_0(E_v)$ gives the claim.  Therefore
$T(E_v)=\int_{E_v}w=\int_{E_v}(u-v)$. Its super-level sets are
$\{w>s\}=\{u>v+s\}=E_{v+s}$ for $s\in(0,m-v)$, so the perimeter and
distribution functions of $w$ at height $s$ coincide with those of $u$ at
height $v+s$; in particular the two level deficits of $w$ at height $s$
equal $D_H(v+s)$ and $D_I(v+s)$.

We now apply the derivation of the level-set deficit identity
(\cref{thm:level-identity}, eq.~\eqref{eq:level-identity}) to $w$ on $E_v$;
the derivation there uses only the flux/coarea identities, the absolute
continuity of the distribution function, and the layer-cake formula, none of
which requires the area to equal $\pi$. For a domain $D$ of area $|D|$ with
torsion function having maximum $m_D$ and level deficits $D_H^D,D_I^D$ it
yields
\[
   \int_0^{m_D}\!\bigl(D_H^D+D_I^D\bigr)\,ds
   \ =\ \tfrac12|D|^2-4\pi T(D)
   \ =\ 4\pi\Bigl(\frac{|D|^2}{8\pi}-T(D)\Bigr)
   \ =\ 4\pi\bigl(T(B_D)-T(D)\bigr),
\]
where $B_D$ is the disc of area $|D|$, $T(B_D)=|D|^2/(8\pi)$; the middle
equality is exactly the computation in the proof of
\cref{thm:level-identity} with $\tfrac12|D|^2$ in place of
$\tfrac12\pi^2$. Taking $D=E_v$, $m_{E_v}=m-v$, and using the identification
of the level deficits of $w$ with $D_H(v+\cdot),D_I(v+\cdot)$,
\[
   4\pi\bigl(T(B_{E_v})-T(E_v)\bigr)
   =\int_0^{m-v}\!\bigl(D_H(v+s)+D_I(v+s)\bigr)\,ds
   =\int_v^m\!\bigl(D_H+D_I\bigr)\,dv'
   \ \le\ \int_0^m\!\bigl(D_H+D_I\bigr)\,dv'
   =4\pi\dT,
\]
the inequality by non-negativity of the integrand (\cref{lem:nonneg}).
Hence $T(B_{E_v})-T(E_v)\le\dT$. Dividing by
$T(B_{E_v})=\mu(v)^2/(8\pi)$ and using $\dstar(\Om)=\tfrac8\pi\dT$, i.e.\
$\dT=\tfrac\pi8\dstar(\Om)$,
\[
   \dstar(E_v)=\frac{T(B_{E_v})-T(E_v)}{T(B_{E_v})}
   \ \le\ \frac{\dT}{\mu(v)^2/(8\pi)}
   =\frac{8\pi\dT}{\mu(v)^2}
   =\frac{8\pi}{\mu(v)^2}\cdot\frac\pi8\,\dstar(\Om)
   =\frac{\pi^2}{\mu(v)^2}\,\dstar(\Om).\qedhere
\]
\end{proof}

The content of \eqref{eq:transfer} is that the \emph{scale-invariant}
torsion deficit of every super-level set is comparable to that of $\Om$
(a tail of a positive integral), even at levels near $\partial\Om$ where the
single-level deficits $D_H(v),D_I(v)$ are individually inflated. This is the
mechanism that carries the small torsion deficit onto the inset.

\subsection{Selection of a good level}

\begin{lemma}[Good level]\label{lem:level}
Let $A_0:=16\pi$. For every $\tau\in(0,\tfrac1{16})$ there is a regular
value $\hat v\in(\tau,2\tau)$ of $u$ with
\begin{equation}\label{eq:goodlevel}
   D_H(\hat v)+D_I(\hat v)\ \le\ A_0\,\frac{\dT}{\tau},
   \qquad
   \mu_B(\hat v)-\mu(\hat v)\ \le\ A_0\,\frac{\dT}{\tau}.
\end{equation}
\end{lemma}

\begin{proof}
Set $I:=(\tau,2\tau)$, of length $|I|=\tau$, and note $I\subset(0,\tfrac18)
\subset(0,\tfrac14)$. By \cref{thm:level-identity} and non-negativity,
$\int_I(D_H+D_I)\,dv\le\int_0^m(D_H+D_I)\,dv=4\pi\dT$. By Chebyshev's
inequality,
\[
   \bigl|\{v\in I:\ D_H(v)+D_I(v)>A_0\dT/\tau\}\bigr|
   \ \le\ \frac{\int_I(D_H+D_I)}{A_0\dT/\tau}
   \ \le\ \frac{4\pi\dT}{A_0\dT/\tau}
   =\frac{4\pi\tau}{A_0}.
\]
By \cref{lem:profile}, $\mu_B-\mu\ge0$ and
$\int_I(\mu_B-\mu)\,dv\le\int_0^{1/4}(\mu_B-\mu)\,dv=\dT$, so again by
Chebyshev,
\[
   \bigl|\{v\in I:\ \mu_B(v)-\mu(v)>A_0\dT/\tau\}\bigr|
   \ \le\ \frac{\dT}{A_0\dT/\tau}
   =\frac{\tau}{A_0}.
\]
With $A_0=16\pi$, the union of these two ``bad'' sets has measure at most
$\dfrac{4\pi\tau}{A_0}+\dfrac{\tau}{A_0}=\dfrac{\tau}{4}+\dfrac{\tau}{16\pi}
<\tau=|I|$. Adding the null set of non-regular values
(\cref{thm:sard}, \cref{lem:regular-values}) does not change this. Hence the
set of regular $\hat v\in I$ satisfying both bounds in \eqref{eq:goodlevel}
has positive measure, and in particular is nonempty.
\end{proof}

From now on $\hat v\in(\tau,2\tau)$ denotes a good regular level as in
\cref{lem:level}, and we write $\hat v<\tfrac14$ (since $\hat v<2\tau<
\tfrac18$), so that $\mu_B(\hat v)=\pi(1-4\hat v)$.

\subsection{Closeness of the super-level set}

\begin{lemma}[Closeness]\label{lem:close}
At the level $\hat v$ of \cref{lem:level},
\begin{equation}\label{eq:close}
   \bigl|\Om\triangle\{u>\hat v\}\bigr|
   \ =\ \pi-\mu(\hat v)
   \ \le\ 8\pi\,\tau+A_0\,\frac{\dT}{\tau}.
\end{equation}
\end{lemma}

\begin{proof}
Since $E_{\hat v}=\{u>\hat v\}\subset\Om$, the symmetric difference is
$\Om\setminus E_{\hat v}=\{x\in\Om:u(x)\le\hat v\}$, of measure
$\pi-\mu(\hat v)$ (the level set $\{u=\hat v\}$ is Lebesgue-null). Split
\[
   \pi-\mu(\hat v)
   =\bigl(\pi-\mu_B(\hat v)\bigr)+\bigl(\mu_B(\hat v)-\mu(\hat v)\bigr).
\]
The first term is $\pi-\pi(1-4\hat v)=4\pi\hat v\le4\pi(2\tau)=8\pi\tau$,
using $\hat v<2\tau$; the second is $\le A_0\dT/\tau$ by
\eqref{eq:goodlevel}. Adding gives \eqref{eq:close}.
\end{proof}

\noindent
We now fix $\delta_\star\in(0,\tfrac1{256}]$ small enough that, whenever
$\dT\le\delta_\star$ and $\tau\in[\sqrt\dT,\tfrac1{16})$ with
$\dT/\tau\le\sqrt\dT$ (e.g.\ at the optimal choice $\tau=\sqrt\dT$, where
the right side of \eqref{eq:close} equals $(8\pi+A_0)\sqrt\dT$), one has
\begin{equation}\label{eq:halfmass}
   8\pi\,\tau+A_0\,\frac{\dT}{\tau}\ \le\ \frac\pi2,
   \qquad\text{hence}\qquad
   \mu(\hat v)\ \ge\ \frac\pi2 .
\end{equation}
Concretely, at $\tau=\sqrt\dT$ this requires $(8\pi+A_0)\sqrt{\delta_\star}
\le\pi/2$, i.e.\ $\delta_\star\le\pi^2/(4(8\pi+A_0)^2)$; with $A_0=16\pi$
this is $\delta_\star\le 1/(2304)$, comfortably compatible with
$\delta_\star\le\tfrac1{256}$ after taking the smaller value. We keep
\eqref{eq:halfmass} as a standing hypothesis for the remainder of the
section (it holds at the optimised scale $\tau=\sqrt\dT$); the value of
$\hat v,\tau$ is otherwise left symbolic and optimised in
\cref{thm:reduction}.

\subsection{One connected inset: the spillover is higher order}

The super-level set $\{u>\hat v\}$ may have several components, but the
total mass of all but the largest is quadratically small in $D_I(\hat v)$,
by a multi-component refinement of the isoperimetric inequality.

\begin{lemma}[Spillover]\label{lem:spill}
Assume \eqref{eq:halfmass}. Let the connected components of $\{u>\hat v\}$
have areas $A_1\ge A_2\ge\cdots>0$, let $\Etil$ be the component of area
$A_1$, and set $S:=\sum_{i\ge2}A_i=\mu(\hat v)-A_1$ for the total spillover.
Then $A_1\ge\pi/4$ and
\begin{equation}\label{eq:spill}
   \bigl|\{u>\hat v\}\setminus\Etil\bigr|
   \ =\ S\ \le\ \frac{D_I(\hat v)^2}{16\pi^3}
   \ \le\ \frac{A_0^2}{16\pi^3}\,\frac{\dT^2}{\tau^2}.
\end{equation}
\end{lemma}

\begin{proof}
\emph{Multi-component isoperimetric inequality.} Perimeter is additive over
the (essentially disjoint) components and the planar isoperimetric
inequality (\cref{thm:isoperimetric}) applies to each, so
$\Per(\hat v)=\sum_i\Per_i\ge\sum_i 2\sqrt{\pi A_i}$ with $\Per_i$ the
perimeter of the $i$-th component. Squaring and using
$\bigl(\sum_i\sqrt{A_i}\bigr)^2=\sum_iA_i+2\sum_{i<j}\sqrt{A_iA_j}$,
\begin{equation}\label{eq:multiiso}
   D_I(\hat v)=\Per(\hat v)^2-4\pi\!\sum_iA_i
   \ \ge\ 4\pi\Bigl[\Bigl(\textstyle\sum_i\sqrt{A_i}\Bigr)^2-\sum_iA_i\Bigr]
   \ =\ 8\pi\!\!\sum_{i<j}\!\sqrt{A_iA_j}.
\end{equation}

\emph{Lower bound on $A_1$.} Since $A_j\le A_1$ we have
$\sqrt{A_j}=A_j/\sqrt{A_j}\ge A_j/\sqrt{A_1}$ for $j\ge2$, hence
$\sum_{j\ge2}\sqrt{A_j}\ge S/\sqrt{A_1}$ and
\[
   \sum_i\sqrt{A_i}\ \ge\ \sqrt{A_1}+\frac{S}{\sqrt{A_1}}
   =\frac{A_1+S}{\sqrt{A_1}}=\frac{\mu(\hat v)}{\sqrt{A_1}}.
\]
Combining with \eqref{eq:multiiso} (in the form $D_I\ge4\pi[(\sum\sqrt{A_i})^2
-\mu(\hat v)]$),
\begin{equation}\label{eq:Aone}
   D_I(\hat v)\ \ge\ 4\pi\Bigl(\frac{\mu(\hat v)^2}{A_1}-\mu(\hat v)\Bigr)
   =4\pi\,\mu(\hat v)\,\frac{\mu(\hat v)-A_1}{A_1}
   =4\pi\,\mu(\hat v)\,\frac{S}{A_1}.
\end{equation}
Suppose, for contradiction, that $A_1<\pi/4$. By \eqref{eq:halfmass},
$\mu(\hat v)\ge\pi/2>A_1$, so there are at least two components and
$S=\mu(\hat v)-A_1>\pi/2-\pi/4=\pi/4$. Then \eqref{eq:Aone} gives
\[
   D_I(\hat v)\ >\ 4\pi\cdot\frac\pi2\cdot\frac{S}{\pi/4}=8\pi\,S>8\pi\cdot\frac\pi4=2\pi^2.
\]
But $D_I(\hat v)\le A_0\dT/\tau\le8\pi\tau+A_0\dT/\tau\le\pi/2<2\pi^2$ by
\eqref{eq:goodlevel} and \eqref{eq:halfmass}, a contradiction. Hence
$A_1\ge\pi/4$.

\emph{Spillover bound.} From \eqref{eq:multiiso},
\[
   D_I(\hat v)\ \ge\ 8\pi\!\sum_{i<j}\!\sqrt{A_iA_j}
   \ \ge\ 8\pi\sqrt{A_1}\sum_{j\ge2}\sqrt{A_j}
   \ \ge\ 8\pi\sqrt{A_1}\,\sqrt{\textstyle\sum_{j\ge2}A_j}
   \ =\ 8\pi\sqrt{A_1\,S},
\]
where the middle step keeps only the pairs $(1,j)$ and the last step uses
$\sum_{j\ge2}\sqrt{A_j}\ge\sqrt{\sum_{j\ge2}A_j}$ (superadditivity of
$t\mapsto\sqrt t$ on non-negative terms). Squaring and using $A_1\ge\pi/4$,
\[
   S\ \le\ \frac{D_I(\hat v)^2}{64\pi^2A_1}
   \ \le\ \frac{D_I(\hat v)^2}{64\pi^2\cdot(\pi/4)}
   =\frac{D_I(\hat v)^2}{16\pi^3}.
\]
The final bound in \eqref{eq:spill} follows from
$D_I(\hat v)\le A_0\dT/\tau$ (\cref{lem:level}).
\end{proof}

\begin{definition}[The inset]\label{def:inset}
Assume $\dT\le\delta_\star$, $\tau$ subject to \eqref{eq:halfmass}, and let
$\hat v\in(\tau,2\tau)$ be a good regular level (\cref{lem:level}). The
\emph{inset} is the connected component $\Etil$ of $\{u>\hat v\}$ of largest
measure, i.e.\ the component of area $A_1$ in \cref{lem:spill}. We write
\[
   \mu_E:=|\Etil|=A_1=\mu(\hat v)-S\ \ge\ \frac\pi4,
   \qquad
   \hat r:=\sqrt{\mu_E/\pi},
   \qquad
   \diso:=\diso(\Etil)=\frac{\Per(\Etil)}{2\sqrt{\pi\mu_E}}-1,
\]
and $\Exc:=\Per(\Etil)-2\pi\hat r=2\pi\hat r\,\diso$.
\end{definition}

By \cref{lem:noholes}\,(ii), $\Etil$ is connected and has no interior hole
\emph{compactly contained in $\Om$}. The stronger \emph{global}
exterior-connectedness ($\R^2\setminus\overline\Etil$ connected) required by
\cref{sec:annulus,sec:fold} is supplied unconditionally by passing to the
hole-filled inset of \cref{lem:fill} below (\cref{ss:fill}); from
\cref{sec:annulus} onward the geometric reduction is run on that saturation,
which inherits all the controls of \cref{prop:insetcontrols} at the same scale.

\subsection{Controls on the inset}

We collect the three controls of the inset used in the subsequent sections.

\begin{proposition}[Inset controls]\label{prop:insetcontrols}
Under \cref{def:inset} there is an absolute constant $C$ such that
\begin{align}
   \diso(\Etil)\ &\le\ \frac{D_I(\hat v)+4\pi S}{8\pi\,\mu_E}
   \ \le\ C\,\frac{\dT}{\tau},
   \label{eq:diso-inset}\\[1mm]
   \dstar(\Etil)\ &\le\ \frac{\pi^2}{\mu_E^2}\,\dstar(\Om)
   +\frac{S^2}{\mu_E^2}
   \ \le\ C\,\dstar(\Om),
   \label{eq:dstar-inset}\\[1mm]
   \bigl|\Om\triangle\Etil\bigr|\ &\le\ 8\pi\,\tau+A_0\frac{\dT}{\tau}
   +\frac{A_0^2}{16\pi^3}\frac{\dT^2}{\tau^2}.
   \label{eq:close-inset}
\end{align}
In particular, with the optimal choice $\tau=\sqrt\dT$ of
\cref{thm:reduction}, $\diso(\Etil)\le C\sqrt\dT$,
$\dstar(\Etil)\le C\dstar(\Om)$ and
$|\Om\triangle\Etil|\le C\sqrt\dT$.
\end{proposition}

\begin{proof}
\emph{\eqref{eq:diso-inset}.} Since $\partial\Etil\subset\{u=\hat v\}$ and
perimeter is additive over components, $\Per(\Etil)\le\Per(\hat v)$, hence
\[
   D_I(\Etil):=\Per(\Etil)^2-4\pi\mu_E
   \ \le\ \Per(\hat v)^2-4\pi\mu_E
   =\bigl(4\pi\mu(\hat v)+D_I(\hat v)\bigr)-4\pi\mu_E
   =D_I(\hat v)+4\pi S,
\]
using $\Per(\hat v)^2=4\pi\mu(\hat v)+D_I(\hat v)$ and
$\mu(\hat v)-\mu_E=S$. From the definition of $\diso$,
$\bigl(1+\diso(\Etil)\bigr)^2=\Per(\Etil)^2/(4\pi\mu_E)=1+D_I(\Etil)/(4\pi\mu_E)$,
so by $\sqrt{1+t}-1\le t/2$,
\[
   \diso(\Etil)\ \le\ \frac{D_I(\Etil)}{8\pi\mu_E}
   \ \le\ \frac{D_I(\hat v)+4\pi S}{8\pi\mu_E}.
\]
By \eqref{eq:spill}, $4\pi S\le 4\pi D_I(\hat v)^2/(16\pi^3)
=D_I(\hat v)^2/(4\pi^2)$; and $D_I(\hat v)\le A_0\dT/\tau\le\pi/2$ by
\eqref{eq:goodlevel}--\eqref{eq:halfmass}, so $4\pi S\le D_I(\hat v)\cdot
\frac{D_I(\hat v)}{4\pi^2}\le\frac{1}{8\pi}D_I(\hat v)\le D_I(\hat v)$.
Hence $D_I(\hat v)+4\pi S\le2D_I(\hat v)\le2A_0\dT/\tau$, and with
$\mu_E\ge\pi/4$,
\[
   \diso(\Etil)\ \le\ \frac{2A_0\dT/\tau}{8\pi\cdot(\pi/4)}
   =\frac{A_0}{\pi^2}\,\frac{\dT}{\tau}\ =:\ C\,\frac{\dT}{\tau}.
\]

\emph{\eqref{eq:dstar-inset}.} We bound $\dstar(\Etil)$ directly, using only
the deficit transfer of \cref{lem:transfer} and the exact additivity of
torsion over connected components; no appeal to \cref{lem:cutdeficit} is
needed (see \cref{rem:fwd}). Let $C_2,C_3,\dots$ be the spillover components
of $E_{\hat v}$, of areas $A_2,A_3,\dots$ with $\sum_{i\ge2}A_i=S$. On each
component $C$ of $E_{\hat v}=\{u>\hat v\}$, the restriction of
$(u-\hat v)_+$ to $C$ belongs to $H^1_0(C)$.  More explicitly, the Sobolev
space on an open set decomposes over its connected components:
if $E=\bigcup_i C_i$ is the disjoint union of open components and
$w\in H^1_0(E)$, then $w_i:=w\mathbf1_{C_i}$ belongs to $H^1_0(C_i)$ and
$\nabla w_i=\mathbf1_{C_i}\nabla w$ a.e.; this follows from the
quasi-continuous trace characterisation of $H^1_0$ and the locality of the
Sobolev gradient.  Applying this to $w=(u-\hat v)_+\in H^1_0(E_{\hat v})$
gives the claim.  Testing against $C_c^\infty(C)$ gives
$-\Delta(u-\hat v)=1$ weakly in $C$, and uniqueness makes this restriction
the torsion function of $C$; hence $T(C)=\int_C(u-\hat v)$. As
$u-\hat v$ is shared and the components are pairwise disjoint, torsion is
\emph{exactly additive}:
\begin{equation}\label{eq:Tadd}
   T(E_{\hat v})=\int_{E_{\hat v}}(u-\hat v)
   =\int_{\Etil}(u-\hat v)+\sum_{i\ge2}\int_{C_i}(u-\hat v)
   =T(\Etil)+\sum_{i\ge2}T(C_i).
\end{equation}
By the Saint--Venant inequality applied to each spillover component
(\cref{prop:scale}, $T(B_E)-T(E)\ge0$, with $T(B_{C_i})=A_i^2/(8\pi)$),
\begin{equation}\label{eq:Tspill}
   \sum_{i\ge2}T(C_i)\ \le\ \sum_{i\ge2}\frac{A_i^2}{8\pi}
   \ \le\ \frac{1}{8\pi}\Bigl(\sum_{i\ge2}A_i\Bigr)^2=\frac{S^2}{8\pi},
\end{equation}
the middle inequality because $\sum A_i^2\le(\sum A_i)^2$ for non-negative
terms. Combining \eqref{eq:Tadd}--\eqref{eq:Tspill} with $\mu_E\le\mu(\hat
v)$,
\[
   T(B_{\Etil})-T(\Etil)
   =\frac{\mu_E^2}{8\pi}-T(E_{\hat v})+\sum_{i\ge2}T(C_i)
   \ \le\ \frac{\mu_E^2-\mu(\hat v)^2}{8\pi}
          +\bigl(T(B_{E_{\hat v}})-T(E_{\hat v})\bigr)+\frac{S^2}{8\pi}
   \ \le\ \dT+\frac{S^2}{8\pi},
\]
where we used $T(B_{\Etil})=\mu_E^2/(8\pi)$, $T(B_{E_{\hat v}})=
\mu(\hat v)^2/(8\pi)$, $\mu_E^2-\mu(\hat v)^2\le0$, and the bound
$T(B_{E_{\hat v}})-T(E_{\hat v})\le\dT$ established in the proof of
\cref{lem:transfer}. Dividing by $T(B_{\Etil})=\mu_E^2/(8\pi)$ and using
$\dT=\tfrac\pi8\dstar(\Om)$,
\begin{equation}\label{eq:dstarE-exact}
   \dstar(\Etil)
   \ \le\ \frac{8\pi\dT}{\mu_E^2}+\frac{S^2}{\mu_E^2}
   \ =\ \frac{\pi^2}{\mu_E^2}\,\dstar(\Om)+\frac{S^2}{\mu_E^2}.
\end{equation}
This is the first inequality in \eqref{eq:dstar-inset} up to the higher-order
term $S^2/\mu_E^2$. The latter is negligible: by \eqref{eq:spill},
$S\le A_0^2\dT^2/(16\pi^3\tau^2)$, and with the standing hypothesis
\eqref{eq:halfmass} (in particular $A_0\dT/\tau\le\pi/2$, i.e.\
$\dT/\tau\le1/32$) one has $S\le\tfrac{A_0^2}{16\pi^3}(\dT/\tau)\cdot\dT
\le\tfrac{A_0^2}{16\pi^3\cdot32}\dT\le\dT$; hence, with $\mu_E\ge\pi/4$,
\[
   \frac{S^2}{\mu_E^2}\ \le\ \frac{S\cdot\dT}{(\pi/4)^2}
   \ \le\ \frac{16}{\pi^2}\,\dT^2
   \ \le\ \frac{16}{\pi^2}\,\dT\cdot\frac\pi8\,\dstar(\Om)
   \ =\ \frac{2}{\pi}\,\dT\,\dstar(\Om)\ \le\ \dstar(\Om),
\]
the last step for $\dT\le\delta_\star\le\pi/2$. Thus, with $\mu_E\ge\pi/4$,
\eqref{eq:dstarE-exact} gives
$\dstar(\Etil)\le\frac{\pi^2}{\mu_E^2}\dstar(\Om)+\dstar(\Om)
\le(16+1)\dstar(\Om)\le C\,\dstar(\Om)$, which is \eqref{eq:dstar-inset}.

\emph{\eqref{eq:close-inset}.} By \cref{lem:close} and \cref{lem:spill},
\[
   |\Om\triangle\Etil|=\bigl(\pi-\mu(\hat v)\bigr)+S
   \ \le\ \Bigl(8\pi\tau+A_0\frac\dT\tau\Bigr)+\frac{A_0^2}{16\pi^3}\frac{\dT^2}{\tau^2}.
\]
For the final assertions, set $\tau=\sqrt\dT$: then $\dT/\tau=\sqrt\dT$, so
$\diso(\Etil)\le C\sqrt\dT$ by \eqref{eq:diso-inset}; the spillover term
$\dT^2/\tau^2=\dT$ is higher order, so $|\Om\triangle\Etil|\le(8\pi+A_0)
\sqrt\dT+C\dT\le C\sqrt\dT$; and \eqref{eq:dstar-inset} is independent of
$\tau$.
\end{proof}

\subsection{Hole-filling: global exterior connectedness at no cost}
\label{ss:fill}

\Cref{lem:noholes} excludes complementary components of $\Etil$ that are
\emph{compactly contained} in $\Om$, but it does not by itself rule out a
bounded complementary component whose closure reaches $\partial\Om$: such a
component evades the interior minimum principle, which is available only away
from $\partial\Om$. The downstream geometry of \cref{sec:annulus,sec:fold},
however, needs the \emph{global} topological fact that the unbounded
complement of the set is its only complement --- the exterior-connectedness
convention of \cref{lem:ext-conn}. We supply this unconditionally by passing
to the \emph{saturation} (hole-filling) of $\Etil$, the standard device of
filling in all bounded complementary components. The decisive quantitative
point is that, by an isoperimetric argument identical in spirit to the
spillover bound \cref{lem:spill}, the filled area is bounded by the
\emph{square} of the isoperimetric deficit, hence is of the same higher order
as the spillover $S$; all inset scales are therefore preserved.

\begin{definition}[Saturation of the inset]\label{def:fill}
The \emph{hole-filled inset} is
\begin{equation}\label{eq:fill-def}
   E^h\ :=\ \Etil\ \cup\ \bigcup\set{K:\ K\text{ a bounded connected
   component of }\R^2\setminus\overline\Etil},
\end{equation}
the union of $\Etil$ with all bounded complementary components of its
$\R^2$-closure. We write $H:=|E^h\setminus\Etil|\ge0$ for the filled area.
\end{definition}

\begin{lemma}[Hole-filling]\label{lem:fill}
Under \cref{def:inset} the hole-filled inset $E^h$ of \cref{def:fill}
satisfies the following.
\begin{enumerate}[label=\textup{(\alph*)},leftmargin=2.4em]
\item\label{it:fill-top} $E^h$ is bounded, open and connected; its complement
$\R^2\setminus\overline{E^h}$ is connected and unbounded \textup{(}it is
exactly the unbounded component of $\R^2\setminus\overline\Etil$\textup{)},
and $\partial E^h\subset\partial\Etil$. Thus $E^h$ is globally hole-free and
satisfies the exterior-connectedness convention \cref{lem:ext-conn}
\textup{(}$\R^2\setminus\overline{E^h}$ connected, with
$\partial(\R^2\setminus\overline{E^h})=\partial E^h$\textup{)}
unconditionally, with no hypothesis on $\partial\Om$.
\item\label{it:fill-per} $\Etil\subset E^h$, $|E^h|=\mu_E+H$, and the
perimeter does not increase:
\begin{equation}\label{eq:fill-per}
   \Per(E^h)\ \le\ \Per(\Etil).
\end{equation}
\item\label{it:fill-key} \textup{(}Key bound.\textup{)} With
$D_I(\Etil)=\Per(\Etil)^2-4\pi\mu_E\ge0$,
\begin{equation}\label{eq:fill-key}
   H\ \le\ \frac{D_I(\Etil)^2}{64\pi^2\,\mu_E}\ \le\ C\,D_I(\Etil)^2 .
\end{equation}
\item\label{it:fill-scales} At the inset scales of
\cref{prop:insetcontrols} \textup{(}$D_I(\Etil)\le C\,\dT/\tau$,
$\mu_E\ge\pi/4$\textup{)} one has $H\le C(\dT/\tau)^2$, so at $\tau=\sqrt\dT$,
$H\le C\,\dT$ \textup{(}higher order, the order of the spillover $S$\textup{)}.
Consequently all the inset controls transfer to $E^h$, with the same rates:
writing $\hat r^h:=\sqrt{|E^h|/\pi}$,
\begin{align}
   |\Om\triangle E^h|\ &\le\ |\Om\triangle\Etil|+H\ \le\ C\,\sqrt\dT,
   &\hat r^h\ &=\ \hat r+O(H)=\hat r+O(\dT);\label{eq:fill-close}\\[1mm]
   \diso(E^h)\ &\le\ \diso(\Etil)\ \le\ C\,\tfrac\dT\tau,
   &\Exc(E^h)\ &\le\ \Exc(\Etil)+C\,H\ \le\ C\,\tfrac\dT\tau;\label{eq:fill-diso}\\[1mm]
   a(E^h)\ &\le\ a(\Etil)+C\,H\ \le\ a(\Etil)+C\,\dT,
   &&\bigl(\text{so }a(E^h)\le C(\tfrac\dT\tau)^{1/2}\bigr);\label{eq:fill-a}\\[1mm]
   \dstar(E^h)\ &\le\ \dstar(\Etil)+\frac{C\,H}{\mu_E}\ \le\ \dstar(\Etil)+C\,\dT
   \ \le\ C\,\dstar(\Om), &&\label{eq:fill-dstar}
\end{align}
where $a(F):=|F\triangle B_{\hat r_F}(z_F)|$ is the ball distance of $F$ at
its optimal centre and $\hat r_F:=\sqrt{|F|/\pi}$.
\end{enumerate}
\end{lemma}

\begin{proof}
\emph{\ref{it:fill-top}.} The open set $\R^2\setminus\overline\Etil$ has, as
the complement of a bounded set, exactly one unbounded connected component
$U_\infty$ (nonempty since $\Etil$ is bounded) and at most countably many
bounded components $\set{K_i}_{i\in I}$, all open and disjoint; by
\eqref{eq:fill-def}, $E^h=\Etil\cup\bigcup_iK_i$, so that as sets
\[
   E^h\ =\ \R^2\setminus\bigl(\partial\Etil\cup U_\infty\bigr),
   \qquad\text{equivalently}\qquad
   \overline{E^h}=\R^2\setminus U_\infty .
\]
We first record the elementary topological facts. Each component $K$ of
$\R^2\setminus\overline\Etil$ (bounded or not) is open with $\partial
K\subset\partial\Etil$: a point $x\in\partial K$ lies neither in the open set
$K$ nor in the open set $\R^2\setminus\overline\Etil\setminus K$ (a union of
other components), so $x\in\overline\Etil$; and $x\notin\Etil$, since $\Etil$
open with $x\in\overline K$ would put a neighbourhood of $x$ inside $\Etil$,
disjoint from $K$. Hence $x\in\partial\Etil$. In particular $\partial
U_\infty\subset\partial\Etil$ and $\partial K_i\subset\partial\Etil$.

\emph{$E^h$ is open.} Since $\Etil$ and each $K_i$ are open, it suffices to
show every $x\in\partial\Etil\cap E^h$ is interior to $E^h$; but in fact no
point of $\partial\Etil$ lies in $E^h$, because $E^h=\Etil\cup\bigcup_iK_i$
is a union of open sets disjoint from $\partial\Etil$
(as $\partial\Etil\cap\Etil=\emptyset$ and $\partial\Etil\cap K_i=\emptyset$,
the latter since $K_i$ is open and $\subset\R^2\setminus\overline\Etil$).
Thus $E^h$ is a union of open sets, hence open. It is bounded since $\Etil$
and all $K_i$ lie in any ball $B$ with $\R^2\setminus B\subset U_\infty$
(such $B$ exists as $U_\infty$ contains the complement of a large ball).

\emph{$E^h$ is connected.} $\Etil$ is connected (\cref{def:inset},
\cref{lem:noholes}). For each $i$, $\partial K_i\subset\partial\Etil$ is
nonempty (a bounded open $K_i\subsetneq\R^2$ has nonempty boundary), so
pick $x\in\partial K_i\subset\partial\Etil$; then $x\in\overline{K_i}\cap
\overline\Etil$, so $\Etil\cup K_i$ is connected (two connected sets with
intersecting closures, both inside the connected $\overline{\Etil\cup K_i}$;
concretely $\overline{K_i}\cap\overline\Etil\ni x$ joins them). All the
$\Etil\cup K_i$ share the connected $\Etil$, so their union $E^h$ is
connected.

\emph{Exterior connectedness.} By the displayed identity
$\overline{E^h}=\R^2\setminus U_\infty$, so $\R^2\setminus\overline{E^h}=
U_\infty$, which is connected and unbounded by construction. Moreover
$\partial E^h=\overline{E^h}\setminus E^h=(\R^2\setminus U_\infty)\setminus
E^h\subset\partial\Etil$ (the points of $\R^2\setminus U_\infty$ not in
$E^h=\Etil\cup\bigcup_iK_i$ lie in $\partial\Etil$, as
$\R^2=\Etil\sqcup\partial\Etil\sqcup(\R^2\setminus\overline\Etil)$ and
$\R^2\setminus\overline\Etil=U_\infty\sqcup\bigcup_iK_i$). Conversely
$\partial U_\infty=\overline{U_\infty}\setminus U_\infty\subset\partial\Etil$
lies in $\partial E^h$ since these points are limits of $U_\infty=\R^2
\setminus\overline{E^h}$ and of $E^h$ (being in $\partial\Etil$, hence limits
of $\Etil\subset E^h$). Thus $\partial(\R^2\setminus\overline{E^h})=\partial
U_\infty=\partial E^h$, and the trichotomy $\R^2=E^h\sqcup\partial E^h\sqcup
(\R^2\setminus\overline{E^h})$ holds. This is exactly the
exterior-connectedness convention of \cref{lem:ext-conn}, established with no
appeal to the minimum principle, hence with no hypothesis on $\partial\Om$.

\emph{\ref{it:fill-per}.} Clearly $\Etil\subset E^h$ and the $K_i$ are
pairwise disjoint and disjoint from $\Etil$, so $|E^h|=\mu_E+\sum_i|K_i|=
\mu_E+H$ with $H=\sum_i|K_i|\ge0$; each $K_i$ is bounded with finite
perimeter $\Per(K_i)<\infty$, since $\Per(\Etil)<\infty$ and
$\partial^*K_i\subset\partial^*\Etil$ $\HH^1$-a.e. (see below).

For the perimeter, we identify the essential boundaries. The plane is
partitioned $\HH^2$-a.e.\ into the measure-theoretic interiors of the
disjoint open sets $\Etil,K_1,K_2,\dots,U_\infty$, and the perimeters add up
along their shared interfaces. Concretely, write $F_\infty:=U_\infty$,
$F_0:=\Etil$, $F_i:=K_i$. Each pair of these sets is separated by the level
set $\set{u=\hat v}\cap\Om$ or sits across $\partial\Etil$, and by the locality
and structure of the reduced boundary \cite[Thm.~16.3, \S16.2]{Maggi2012} the
essential boundary of $\Etil$ decomposes $\HH^1$-a.e.\ as
\begin{equation}\label{eq:fill-boundary}
   \partial^*\Etil\ =\ \bigl(\partial^*\Etil\cap\partial^*U_\infty\bigr)
   \ \sqcup\ \bigsqcup_i\bigl(\partial^*\Etil\cap\partial^*K_i\bigr)
   \qquad(\HH^1\text{-a.e.}),
\end{equation}
since every point of $\partial^*\Etil$ has, on its outer side, density-one
membership in exactly one of the complementary components $U_\infty,K_1,K_2,
\dots$. Filling the bounded components $K_i$ \emph{removes} from the boundary
exactly the interfaces $\partial^*\Etil\cap\partial^*K_i$ (now interior to
$E^h$) and keeps the outer interface $\partial^*\Etil\cap\partial^*U_\infty=
\partial^*E^h$; moreover, since each $K_i$ has \emph{all} of its essential
boundary shared with $\Etil$ or with other filled components (its complement
in $\R^2$ is $\Etil\cup U_\infty\cup\bigcup_{j\ne i}K_j$, whose density-one
points exhaust $\partial^*K_i$), filling adds no new reduced boundary. Hence
$\partial^*E^h=\partial^*\Etil\cap\partial^*U_\infty=\partial^*\Etil
\setminus\bigsqcup_i\partial^*K_i$ ($\HH^1$-a.e.), and by additivity of
$\HH^1$ over \eqref{eq:fill-boundary},
\begin{equation}\label{eq:fill-perimexact}
   \Per(E^h)=\HH^1(\partial^*E^h)
   =\HH^1(\partial^*\Etil)-\sum_i\HH^1(\partial^*\Etil\cap\partial^*K_i)
   =\Per(\Etil)-\sum_i\Per(K_i),
\end{equation}
the last equality because $\partial^*K_i\subset\partial^*\Etil$ $\HH^1$-a.e.\
(every essential-boundary point of $K_i$ adjacent to $\Etil$ is one of
$\Etil$ adjacent to $K_i$; the components $K_i$ do not touch each other or
$U_\infty$ across a set of positive $\HH^1$-measure, being separated by
$\overline\Etil$). In particular $\Per(E^h)\le\Per(\Etil)$, which is
\eqref{eq:fill-per}.

\emph{\ref{it:fill-key}.} Let the filled holes $K_i$ have areas $A_i^h:=|K_i|$
and perimeters $p_i:=\Per(K_i)$, so $H=\sum_iA_i^h$. By \ref{it:fill-per},
filling removes exactly the boundaries of the holes, $\Per(E^h)=\Per(\Etil)
-\sum_i p_i$ (eq.~\eqref{eq:fill-perimexact}).
By the planar isoperimetric inequality (\cref{thm:isoperimetric}) applied to
each hole, $p_i\ge2\sqrt{\pi A_i^h}$; and by superadditivity of
$t\mapsto\sqrt t$ on non-negative terms,
\begin{equation}\label{eq:fill-superadd}
   \sum_i p_i\ \ge\ 2\sqrt\pi\sum_i\sqrt{A_i^h}\ \ge\ 2\sqrt\pi
   \sqrt{\textstyle\sum_iA_i^h}\ =\ 2\sqrt{\pi H}.
\end{equation}
On the other hand the isoperimetric inequality for $E^h$ itself, of area
$\mu_E+H$, gives $\Per(E^h)\ge2\sqrt{\pi(\mu_E+H)}$. Substituting both into
\eqref{eq:fill-perimexact},
\[
   2\sqrt{\pi(\mu_E+H)}\ \le\ \Per(E^h)\ =\ \Per(\Etil)-\sum_i p_i
   \ \le\ \Per(\Etil)-2\sqrt{\pi H},
\]
so that
\[
   \Per(\Etil)\ \ge\ 2\sqrt\pi\Bigl(\sqrt{\mu_E+H}+\sqrt{H}\Bigr).
\]
Squaring and using $\sqrt{(\mu_E+H)H}\ge\sqrt{\mu_E H}$,
\[
   \Per(\Etil)^2\ \ge\ 4\pi\Bigl(\mu_E+2H+2\sqrt{H(\mu_E+H)}\Bigr)
   \ \ge\ 4\pi\mu_E+8\pi\sqrt{H\mu_E},
\]
whence $D_I(\Etil)=\Per(\Etil)^2-4\pi\mu_E\ge8\pi\sqrt{H\mu_E}$, i.e.
\[
   H\ \le\ \frac{D_I(\Etil)^2}{64\pi^2\mu_E},
\]
which is \eqref{eq:fill-key}; the second inequality there uses
$\mu_E\ge\pi/4$ (\cref{lem:spill}), giving the absolute constant
$C=1/(16\pi^3)$.

\emph{\ref{it:fill-scales}.} By \eqref{eq:diso-inset} and its proof,
$D_I(\Etil)\le D_I(\hat v)+4\pi S\le2D_I(\hat v)\le2A_0\,\dT/\tau$, so
\eqref{eq:fill-key} gives $H\le C(\dT/\tau)^2$, and at $\tau=\sqrt\dT$,
$H\le C\dT$, of the same higher order as the spillover $S$ of
\cref{lem:spill}. The consequences follow:

\noindent\emph{Closeness and radius.} Since $\Etil\subset E^h$ and
$\Etil\subset\Om$ (so $|\Etil\setminus\Om|=0$), we split the symmetric
difference: $|\Om\setminus E^h|\le|\Om\setminus\Etil|$ (as $\Etil\subset E^h$
shrinks the part of $\Om$ outside), while $|E^h\setminus\Om|\le|E^h\setminus
\Etil|+|\Etil\setminus\Om|=H+0=H$. Hence
\[
   |\Om\triangle E^h|=|\Om\setminus E^h|+|E^h\setminus\Om|
   \ \le\ |\Om\setminus\Etil|+H=|\Om\triangle\Etil|+H
   \ \le\ C\sqrt\dT+C\dT\le C\sqrt\dT,
\]
which is \eqref{eq:fill-close}; no assertion that $E^h\subset\Om$ is needed.
The radius obeys $\hat r^h=\sqrt{(\mu_E+H)/\pi}=
\hat r\sqrt{1+H/\mu_E}=\hat r+O(H/\hat r)=\hat r+O(\dT)$ by $H\le C\dT$ and
$\hat r\ge1/2$.

\noindent\emph{Isoperimetric deficit.} Since $\Per(E^h)\le\Per(\Etil)$
\eqref{eq:fill-per} and $|E^h|=\mu_E+H\ge\mu_E$,
\[
   \diso(E^h)=\frac{\Per(E^h)}{2\sqrt{\pi|E^h|}}-1
   \ \le\ \frac{\Per(\Etil)}{2\sqrt{\pi\mu_E}}-1=\diso(\Etil)\le C\,\tfrac\dT\tau,
\]
which is the first bound in \eqref{eq:fill-diso}; and $\Exc(E^h)=
\Per(E^h)-2\pi\hat r^h\le\Per(\Etil)-2\pi\hat r+2\pi(\hat r-\hat r^h)\le
\Exc(\Etil)+CH\le C\,\dT/\tau$, using $|\hat r-\hat r^h|\le CH$.

\noindent\emph{Ball distance.} For any common centre,
$|E^h\triangle B_{\hat r^h}(z)|\le|\Etil\triangle B_{\hat r}(z)|+|E^h
\setminus\Etil|+|B_{\hat r^h}(z)\triangle B_{\hat r}(z)|\le a(\Etil)+H+
C|\hat r^h-\hat r|\le a(\Etil)+CH$; minimising over $z$ gives $a(E^h)\le
a(\Etil)+CH\le a(\Etil)+C\dT$, which is \eqref{eq:fill-a}, and the parenthetical
rate follows from $a(\Etil)\le C(\dT/\tau)^{1/2}$ and $\dT\le(\dT/\tau)^{1/2}$
(as $\dT/\tau\le1$).

\noindent\emph{Scale-invariant deficit.} Torsional rigidity is monotone under
set inclusion: from the variational characterisation
$T(F)=\max\set{2\int_F v-\int_F|\nabla v|^2:\ v\in H^1_0(F)}$ \textup{(}attained
at the torsion function of $F$\textup{)} and the inclusion $H^1_0(\Etil)\subset
H^1_0(E^h)$ (extension by zero, valid as $\Etil\subset E^h$ are open), one has
$T(E^h)\ge T(\Etil)$ for $\Etil\subset E^h$. Hence, with $|E^h|=\mu_E+H$,
\[
   1-\dstar(E^h)=\frac{8\pi T(E^h)}{|E^h|^2}
   \ \ge\ \frac{8\pi T(\Etil)}{(\mu_E+H)^2}
   =\frac{1-\dstar(\Etil)}{(1+H/\mu_E)^2},
\]
so
\[
   \dstar(E^h)\ \le\ 1-\frac{1-\dstar(\Etil)}{(1+H/\mu_E)^2}
   \ \le\ \dstar(\Etil)+\Bigl(1-\frac1{(1+H/\mu_E)^2}\Bigr)
   \ \le\ \dstar(\Etil)+\frac{2H}{\mu_E},
\]
using $1-(1+t)^{-2}\le2t$ for $t\ge0$. With $\mu_E\ge\pi/4$ and $H\le C\dT$
this is $\dstar(E^h)\le\dstar(\Etil)+C\dT\le C\dstar(\Om)$ by
\eqref{eq:dstar-inset} and $\dT=\tfrac\pi8\dstar(\Om)$, which is
\eqref{eq:fill-dstar}.
\end{proof}

\begin{remark}[The hole-filled inset is the object of
\cref{sec:annulus,sec:fold}]\label{rem:fill-use}
\Cref{lem:fill} resolves, unconditionally, the global exterior-connectedness
property required by the trapping annulus and the inscribed-ball lemma: the
hole-filled inset $E^h$ is connected and satisfies $\R^2\setminus
\overline{E^h}$ connected, with all the inset scales of
\cref{prop:insetcontrols} preserved up to an additive $O(\dT)$, of the same
higher order as the spillover. From \cref{sec:annulus} onward the geometric
reduction is therefore applied to $E^h$ in place of $\Etil$; the only change
in any estimate is the replacement of $\mu_E$ by $\mu_E+H=\mu_E+O(\dT)$ and
of $\hat r,\diso,\Exc,a,\dstar$ by their $E^h$-counterparts, with identical
rates. The mass $H$ added by filling, like the spillover $S$ subtracted by
selecting the largest component, is quadratic in the isoperimetric deficit
and hence negligible at the sharp scale.
\end{remark}

\begin{remark}[Self-contained deficit transfer; no forward dependence]
\label{rem:fwd}
The bound \eqref{eq:dstar-inset} is proved above without any appeal to the
cutting principle \cref{lem:cutdeficit} of \cref{sec:cut}: the passage from
$E_{\hat v}$ to its largest component $\Etil$ is handled by the exact
additivity of torsion over components \eqref{eq:Tadd} and the Saint--Venant
bound \eqref{eq:Tspill} on the spillover torsion, both internal to
\cref{sec:identity,sec:inset}. There is therefore \emph{no} circular
dependence on \cref{sec:cut} in this section. (One could alternatively invoke
\cref{lem:cutdeficit}; but its precise hypothesis --- that the excised set lie
in a thin concentric collar of $\Etil$ --- is \emph{not} met by the spillover,
whose components are in general not contained in the trapping annulus of
$\Etil$. The additivity argument avoids this issue entirely, which is why we
prefer it.) For the record, \cref{lem:cutdeficit} itself rests only on
Green-function monotonicity and a boundary-layer barrier, so even the optional
route would be acyclic.
\end{remark}

\begin{remark}[The scales, recorded]\label{rem:scales}
With $\tau=\sqrt\dT$ (\cref{thm:reduction}) the inset $\Etil$ satisfies
$\diso=\diso(\Etil)\le C\sqrt\dT$ and $\Exc=2\pi\hat r\,\diso\le C\sqrt\dT$.
Feeding $\diso$ into the sharp quantitative isoperimetric inequality
(\cref{thm:fmp}) gives the ball-distance
$a:=|\Etil\triangle B_{\hat r}(z)|\le C_{\mathrm F}\sqrt{\diso}\,\mu_E\le
C\,\dT^{1/4}$ at the optimal centre $z$. The trapping width then satisfies
$\eta\le C(\Exc+a)^{1/2}\le C\,\dT^{1/8}$ (\cref{thm:annulus}), an
\emph{absolute} small amplitude. These are the inputs to the annulus and
fold sections; the deficit $\dstar(\Etil)\le C\dstar(\Om)$ and the closeness
$|\Om\triangle\Etil|\le C\sqrt\dT$ are the inputs to the conclusion. We
emphasise that closeness is carried by the \emph{level} $\hat v\sim\sqrt\dT$
through \cref{lem:close}, not by the amplitude $\eta$; and that the deficit
transfer \cref{lem:transfer} delivers $\dstar(\Etil)\lesssim\dstar(\Om)$
without recourse to the inflated single-level deficits.
\end{remark}

\section{The trapping annulus}\label{sec:annulus}

In this section we prove that the connected inset $\Etil$ produced in
\cref{sec:inset} is trapped between two concentric circles whose radii
differ by an \emph{absolutely small} amount, controlled by the perimeter
excess and the $L^1$ distance to a ball. The decisive point is that the
width of the trapping annulus is bounded \emph{without any second-order
information} on $\partial\Etil$: no interior or exterior ball condition, no
$C^2$ regularity, no Magnanini--Poggesi oscillation estimate. The only
inputs are the connectedness and the no-hole structure of $\Etil$
(\cref{lem:noholes}), the planar isoperimetric deficit, and the closeness
to a ball furnished by \cref{thm:fmp}.

The mechanism is purely geometric: the boundary of a connected, hole-free
set must cross every circle of intermediate radius at least twice; summed
against the radial coordinate through the coarea formula, this forces a
lower bound on the \emph{angular} part of the perimeter, which a single
Cauchy--Schwarz inequality converts into the deficit. We work throughout
with the good representative of $\Etil$ described in
\cref{thm:torsion-reg,lem:regular-values}: since $\hat v$ is a regular value
of the (real-analytic, in $\Om$) torsion function $u$ and $\Etil$ is one
connected component of $\{u>\hat v\}$, the set $\Etil$ is a bounded open
connected subset of $\Om$, with $\overline\Etil\subset\subset\Om$, whose
topological boundary
\[
   \partial\Etil=\set{u=\hat v}\cap\partial\Etil
\]
is a compact, $\HH^1$-rectifiable curve along which $\nabla u\neq0$; in
particular $\partial\Etil$ coincides $\HH^1$-a.e.\ with the reduced and the
essential boundary, and $\HH^1(\partial\Etil)=P(\Etil)$. The plane decomposes
as the disjoint union
\begin{equation}\label{eq:trichotomy}
   \R^2=\Etil\ \sqcup\ \partial\Etil\ \sqcup\ \bigl(\R^2\setminus\overline\Etil\bigr),
\end{equation}
the first and last sets being open. We record the precise topological
consequence of the no-hole property that we shall use.

\begin{lemma}[The complement of the inset is connected]\label{lem:ext-conn}
The open set $U:=\R^2\setminus\overline\Etil$ is connected, and
$\partial U=\partial\Etil$.
\end{lemma}

\begin{proof}
Since $\overline\Etil\subset\subset\Om$ is compact, $U$ is the complement of
a compact set and so has exactly one unbounded connected component $U_\infty$
(the points of large modulus form a connected set lying in $U$). Suppose
$U$ had another (necessarily bounded) component $H$. By \cref{lem:noholes}
no bounded component of $\R^2\setminus\overline\Etil$ is compactly contained
in $\Om$, so $\overline H$ meets $\partial\Om$; pick $q\in\overline H\cap
\partial\Om$. As $\overline\Etil\subset\subset\Om$, the point $q$ and the
entire ray from $q$ heading away from the bounded set $\overline\Etil$ lie
in $U$ and reach arbitrarily large modulus, hence lie in $U_\infty$; thus
$q\in\overline{U_\infty}$ and, $H$ being open and adjacent to $q$, $H$ meets
$U_\infty$, forcing $H=U_\infty$, a contradiction with $H$ bounded.
Therefore $U=U_\infty$ is connected. Finally $\partial U=\partial(\R^2
\setminus\overline\Etil)=\partial\overline\Etil=\partial\Etil$, the last
equality because $\Etil$ equals the interior of its closure for the good
representative.
\end{proof}

Throughout, $z\in\R^2$ denotes the centre of an optimal ball realising the
Fraenkel asymmetry of $\Etil$; we take $z\in\Etil$ (an interior point), which
holds for the inset as shown in \cref{rem:zinE}. With
$\hat r:=\sqrt{|\Etil|/\pi}\in[\tfrac12,1]$ (the bounds follow from
$\pi/4\le|\Etil|\le\pi$, \cref{lem:spill})
\begin{equation}\label{eq:a-def}
   a:=|\Etil\triangle B_{\hat r}(z)|
   =\min_{x_0\in\R^2}|\Etil\triangle B_{\hat r}(x_0)|
   =\AF(\Etil)\,|\Etil| .
\end{equation}
We write $e_r:=(x-z)/|x-z|$ for the radial unit field (defined on
$\R^2\setminus\set z$), $e_\theta:=e_r^{\perp}$ for the orthogonal angular
field, $\nu$ for the outward unit normal on the rectifiable boundary, and
\[
   \nu_r:=\nu\cdot e_r,\qquad
   \nu_\theta:=\nu\cdot e_\theta,\qquad
   \nu_r^2+\nu_\theta^2=1 \quad(\HH^1\text{-a.e.\ on }\partial\Etil),
\]
\[
   P_r:=\int_{\partial\Etil}|\nu_r|\,\d\HH^1,
   \qquad
   P_\theta:=\int_{\partial\Etil}|\nu_\theta|\,\d\HH^1,
   \qquad
   0\le P_r,\,P_\theta\le P(\Etil).
\]
Finally we record the standing scaling assumption of the construction,
$\dT\le\delta_\star$, together with the deficit scales carried in from
\cref{sec:inset} (here $\tau$ is the level-selection parameter of that
section, kept symbolic until the optimisation in \cref{sec:main}):
\begin{equation}\label{eq:ann-scales}
   \Exc:=P(\Etil)-2\pi\hat r=2\pi\hat r\,\diso(\Etil)\le C\,\frac{\dT}{\tau},
   \qquad
   a\le C_{\mathrm F}\sqrt{\diso(\Etil)}\;|\Etil|\le C\Bigl(\frac{\dT}{\tau}\Bigr)^{1/2},
\end{equation}
the first from \cref{lem:spill}, the second from the planar quantitative
isoperimetric inequality \cref{thm:fmp}. In particular $\Exc+a\le C(\dT/\tau)
^{1/2}$ for $\dT/\tau\le1$, so the trapping width produced below will scale
like $(\dT/\tau)^{1/4}$.

\subsection{The inner and outer radii}

\begin{definition}[Inner and outer radius]\label{def:radii}
For the optimal centre $z\in\Etil$ set
\begin{equation}\label{eq:radii-def}
   \rho_i:=\sup\set{r>0:\ B_r(z)\subseteq\Etil},
   \qquad
   \rho_e:=\inf\set{r>0:\ \Etil\subseteq B_r(z)} .
\end{equation}
\end{definition}

We will use Definition~\ref{def:radii} under the hypothesis $z\in\Etil$,
which guarantees $\rho_i>0$; that $z\in\Etil$ holds in the application is
addressed in \cref{rem:zinE}. The next lemma collects the elementary
properties of these radii; they are the rigorous substitute for the (here
unavailable) minimum and maximum of $|x-z|$ over a smooth boundary.

\begin{lemma}[Basic properties of $\rho_i,\rho_e$]\label{lem:radii}
Assume $z\in\Etil$. Then $0<\rho_i\le\rho_e<\infty$ and
\begin{equation}\label{eq:inclusions}
   B_{\rho_i}(z)\subseteq\Etil\subseteq \overline{B_{\rho_e}(z)} .
\end{equation}
Moreover $\rho_i=\min_{x\in\partial\Etil}|x-z|$ and
$\rho_e=\max_{x\in\partial\Etil}|x-z|=\max_{x\in\overline\Etil}|x-z|$; both
extrema are attained, and
\begin{equation}\label{eq:rhat-between}
   \rho_i\le\hat r\le\rho_e .
\end{equation}
\end{lemma}

\begin{proof}
Since $z\in\Etil$ and $\Etil$ is open, some $B_r(z)\subseteq\Etil$, so
$\rho_i>0$; since $\Etil$ is bounded, $\rho_e<\infty$, and $B_r(z)\subseteq
\Etil\subseteq B_{r'}(z)$ forces $r\le r'$, whence $\rho_i\le\rho_e$. For the
left inclusion in \eqref{eq:inclusions}, $B_{\rho_i}(z)=\bigcup_{r<\rho_i}
B_r(z)\subseteq\Etil$. For the right inclusion, $\overline\Etil
\subseteq\overline{B_{\rho_e}(z)}$ because $\Etil\subseteq B_{r}(z)$ for
every $r>\rho_e$.

Set $\rho_i':=\min\set{|x-z|:x\in\partial\Etil}$ and $\rho_e':=\max\set{|x-z|:
x\in\partial\Etil}$, both attained since $\partial\Etil$ is compact and
nonempty.

\emph{The outer radius.} First, $\rho_e=\sup_{x\in\Etil}|x-z|$: indeed
$\Etil\subseteq B_r(z)$ holds iff $r>\sup_{x\in\Etil}|x-z|$, so the infimum
defining $\rho_e$ equals this supremum. We claim this supremum equals
$\rho_e'$, the maximal boundary radius. On one hand $\rho_e'\le\sup_{\Etil}
|x-z|$: a boundary point $w$ of radius $\rho_e'$ is a limit of points of
$\Etil$, whose radii approach $\rho_e'$. On the other hand $\sup_{\Etil}|x-z|
\le\rho_e'$: suppose some $x\in\Etil$ had $|x-z|>\rho_e'$. Extend the ray
from $z$ through $x$ to a far point $y$ with $y\notin\overline\Etil$
(possible, $\Etil$ being bounded). The set $\set{t\in[0,1]:(1-t)x+ty\in
\overline\Etil}$ is closed, contains $0$ but not $1$; its supremum $t^*$
gives a point $p=(1-t^*)x+t^*y\in\overline\Etil$ with all nearby points
further out on the ray outside $\overline\Etil$, so $p\in\partial\Etil$.
As $p$ lies on the outward ray beyond $x$, $|p-z|\ge|x-z|>\rho_e'$,
contradicting the maximality of $\rho_e'$. Hence $\rho_e=\sup_{\Etil}|x-z|
=\rho_e'$ is attained on $\partial\Etil$, and \eqref{eq:inclusions} gives
$\Etil\subseteq\overline{B_{\rho_e}(z)}$.

\emph{The inner radius.} No boundary point has radius $<\rho_i$, because the
open ball $B_{\rho_i}(z)\subseteq\Etil$ contains no point of $\partial\Etil$;
thus $\rho_i'\ge\rho_i$. Conversely $\rho_i'\le\rho_i$: the open ball
$B_{\rho_i'}(z)$ contains no boundary point (all of $\partial\Etil$ has
radius $\ge\rho_i'$), so by \eqref{eq:trichotomy} it is partitioned by the
two disjoint open sets $\Etil$ and $\R^2\setminus\overline\Etil$; being
connected and containing $z\in\Etil$, it lies entirely in $\Etil$, whence
$\rho_i\ge\rho_i'$ by definition of $\rho_i$. Hence $\rho_i=\rho_i'$,
attained at some $w_i\in\partial\Etil$ with $|w_i-z|=\rho_i$.

\emph{Position of $\hat r$.} Finally $B_{\rho_i}(z)\subseteq\Etil$ gives
$\pi\rho_i^2\le|\Etil|=\pi\hat r^2$, and $\Etil\subseteq\overline{B_{\rho_e}
(z)}$ gives $\pi\hat r^2=|\Etil|\le\pi\rho_e^2$, which is
\eqref{eq:rhat-between}.
\end{proof}

\subsection{The angular coarea identity}

For $\rho>0$ let
\begin{equation}\label{eq:Ndef}
   N(\rho):=\HH^0\bigl(\partial\Etil\cap\partial B_\rho(z)\bigr)
   \in\set{0,1,2,\dots,\infty}
\end{equation}
be the number of boundary points on the circle of radius $\rho$ about $z$.

\begin{lemma}[Angular coarea identity]\label{lem:angular-coarea}
With $P_\theta$ and $N$ as above,
\begin{equation}\label{eq:angular-coarea}
   P_\theta=\int_{\partial\Etil}|\nu_\theta|\,\d\HH^1
   =\int_0^\infty N(\rho)\,\d\rho .
\end{equation}
\end{lemma}

\begin{proof}
The radial function $f(x):=|x-z|$ is $1$-Lipschitz on $\R^2$, with
$\nabla f=e_r$ on $\R^2\setminus\set z$. Let $\nabla^{\partial\Etil}f$ denote
its tangential gradient along the rectifiable curve $\partial\Etil$, i.e.\
the orthogonal projection of $\nabla f$ onto the (approximate) tangent line,
which $\HH^1$-a.e.\ is the line $\R\,\nu^\perp$. Writing $\nu=\nu_r e_r
+\nu_\theta e_\theta$ and decomposing $e_r=(e_r\cdot\nu)\nu+(e_r\cdot
\nu^\perp)\nu^\perp$,
\[
   |\nabla^{\partial\Etil}f|=|e_r\cdot\nu^\perp|
   =\sqrt{1-(e_r\cdot\nu)^2}=\sqrt{1-\nu_r^2}=|\nu_\theta|
   \qquad(\HH^1\text{-a.e.}).
\]
The coarea formula for Lipschitz functions on $\HH^1$-rectifiable sets
\textup{(}\cref{thm:coarea}; see also \cite[Theorem~2.93]{AFP2000} and
\cite[\S3.4.2]{EvansGariepy2015}\textup{)} gives
\[
   \int_{\partial\Etil}|\nabla^{\partial\Etil}f|\,\d\HH^1
   =\int_0^\infty\HH^0\bigl(\partial\Etil\cap f^{-1}(\rho)\bigr)\,\d\rho
   =\int_0^\infty N(\rho)\,\d\rho,
\]
which is \eqref{eq:angular-coarea}.
\end{proof}

We record the companion \emph{weighted} identity, used once below. Let
$\Theta:\R^2\setminus\set z\to S^1$, $\Theta(x)=(x-z)/|x-z|$, be the angular
projection (well defined on $\partial\Etil$ since $z\in\Etil$ is interior).
A computation identical to the one above, now for $\Theta$, gives
$|\nabla^{\partial\Etil}\Theta|=|\nu_r|/|x-z|$, so the coarea formula with
the weight $g(x)=|x-z|$ yields
\begin{equation}\label{eq:radial-weighted}
   P_r=\int_{\partial\Etil}|\nu_r|\,\d\HH^1
     =\int_{\partial\Etil}|x-z|\,|\nabla^{\partial\Etil}\Theta|\,\d\HH^1
     =\int_{S^1}\Bigl(\sum_{x\in\partial\Etil\cap\,\mathrm{ray}(\theta)}|x-z|\Bigr)\d\theta,
\end{equation}
where $\mathrm{ray}(\theta):=\set{z+\rho e^{i\theta}:\rho>0}$. (We only use
the elementary consequence \eqref{eq:Pr-lower} below.)

\subsection{Two crossings at every intermediate radius}

\begin{lemma}[$N\ge2$ on the trapping interval]\label{lem:Ncross}
Assume $z\in\Etil$, and recall $U:=\R^2\setminus\overline\Etil$ is connected
with $\partial U=\partial\Etil$ \textup{(}\cref{lem:ext-conn}, from the
no-hole property \cref{lem:noholes}\textup{)}. Then
\begin{equation}\label{eq:Ngeq2}
   N(\rho)\ge2\qquad\text{for every }\rho\in(\rho_i,\rho_e).
\end{equation}
\end{lemma}

\begin{proof}
Fix $\rho\in(\rho_i,\rho_e)$ and let $C_\rho:=\partial B_\rho(z)$, a circle
(homeomorphic to $S^1$, hence connected). By \eqref{eq:trichotomy} every
point of $C_\rho$ lies in exactly one of the open sets $\Etil$,
$U:=\R^2\setminus\overline\Etil$, or in $\partial\Etil$. We show that
$C_\rho$ meets both $\Etil$ and $U$; the conclusion then follows from a
connectedness count.

\emph{Step 1: $C_\rho$ meets $\Etil$.} The function $f(x)=|x-z|$ is
continuous on the connected open set $\Etil$ (open connected subsets of
$\R^2$ are path-connected), so $f(\Etil)$ is an interval. It contains $0$
(as $z\in\Etil$, take $f(z)=0$) and, by \cref{lem:radii}, values
arbitrarily close to $\rho_e$ from below (points of $\Etil$ of radius
$\to\rho_e$). Hence $f(\Etil)\supseteq[0,\rho_e)\ni\rho$, so some point of
$\Etil$ has radius $\rho$, i.e.\ $C_\rho\cap\Etil\neq\varnothing$.

\emph{Step 2: $C_\rho$ meets $U$.} By \cref{lem:radii} there is a boundary
point $w_i\in\partial\Etil$ with $|w_i-z|=\rho_i$. For the good
representative the topological boundaries coincide, $\partial U=\partial
\overline\Etil=\partial\Etil$, so $w_i\in\overline U$, and every
neighbourhood of $w_i$ meets the open set $U$; thus $U$ contains points of
radius arbitrarily close to $\rho_i$, all of radius $\ge\rho_i$ since
$B_{\rho_i}(z)\subseteq\Etil$. Hence $\inf_{U}f=\rho_i$ (not necessarily
attained). On the
other hand $U$ is unbounded (it is the complement of the bounded set
$\overline\Etil$), so $\sup_U f=\infty$. Since $U$ is connected and open,
hence path-connected, $f(U)$ is an interval with infimum $\rho_i$ and
supremum $\infty$, so $f(U)\supseteq(\rho_i,\infty)\ni\rho$. Therefore some
point of $U$ has radius $\rho$, i.e.\ $C_\rho\cap U\neq\varnothing$.

\emph{Step 3: at least two boundary crossings.} Suppose, for contradiction,
that $\HH^0(C_\rho\cap\partial\Etil)\le1$. Then $C_\rho\setminus\partial\Etil$
is $C_\rho$ with at most one point removed, hence connected. By
\eqref{eq:trichotomy}, $C_\rho\setminus\partial\Etil=(C_\rho\cap\Etil)\sqcup
(C_\rho\cap U)$ is a partition into two disjoint sets that are open in
$C_\rho$ and, by Steps 1--2, both nonempty. A connected topological space
admits no such partition---a contradiction. Hence
$N(\rho)=\HH^0(C_\rho\cap\partial\Etil)\ge2$.
\end{proof}

\begin{remark}[Where connectedness and no-holes enter]\label{rem:why}
Connectedness of $\Etil$ is used in Step~1 to drive the radius continuously
from $0$ to near $\rho_e$; the no-hole hypothesis (connectedness of
$U=\R^2\setminus\overline\Etil$) is used in Step~2 to drive the
\emph{exterior} radius continuously from near $\rho_i$ out to $\infty$. If
$\Etil$ had a hole, the bounded complementary component could sit inside a
circle $C_\rho$ without forcing a second crossing, and the argument would
fail. The bound \eqref{eq:Ngeq2} holds for \emph{every}
$\rho\in(\rho_i,\rho_e)$, not merely almost every $\rho$; the a.e.\ caveat
enters only through the coarea formula \eqref{eq:angular-coarea}, where
$N(\rho)$ is finite for a.e.\ $\rho$ because $P_\theta\le P(\Etil)<\infty$.
\end{remark}

\subsection{The radial-profile bound}

We now bound the angular defect $P-P_r$ of the perimeter by the geometric
quantities $\Exc$ and $a$. Define the \emph{outer radial profile}
\begin{equation}\label{eq:Rdef}
   R(\theta):=\sup\set{\rho>0:\ z+\rho e^{i\theta}\in\Etil}\in[0,\rho_e]
   \qquad(\theta\in S^1).
\end{equation}
Because $z\in\Etil$ and $\Etil$ is bounded and open, the ray
$\mathrm{ray}(\theta)$ starts inside $\Etil$ and leaves it, so it meets
$\partial\Etil$; for a.e.\ $\theta$ the meeting set is finite with
outermost radius equal to $R(\theta)$, and
\begin{equation}\label{eq:Pr-lower}
   P_r=\int_{S^1}\Bigl(\sum_{x\in\partial\Etil\cap\,\mathrm{ray}(\theta)}|x-z|\Bigr)\d\theta
   \ \ge\ \int_{S^1}R(\theta)\,\d\theta,
\end{equation}
the equality being \eqref{eq:radial-weighted} and the inequality holding
because the inner sum is at least its largest term $R(\theta)$.

\begin{lemma}[Radial-profile bound]\label{lem:radial-profile}
Assume $z\in\Etil$. Then
\begin{equation}\label{eq:radial-profile}
   P(\Etil)-P_r\ \le\ \Exc+\frac{2a}{\hat r}\ \le\ \Exc+4a .
\end{equation}
\end{lemma}

\begin{proof}
For each $\theta$ with $R(\theta)<\hat r$, every point $z+\rho e^{i\theta}$
with $R(\theta)<\rho<\hat r$ has radius exceeding $R(\theta)$, hence lies
\emph{outside} $\Etil$ (by definition of $R$), while having radius $<\hat r$,
hence inside $B_{\hat r}(z)$; thus the radial segment $\set{R(\theta)<\rho
<\hat r}$ along direction $\theta$ lies in $B_{\hat r}(z)\setminus\Etil$.
Integrating in polar coordinates and using
$\int_{R}^{\hat r}\rho\,\d\rho=\tfrac12(\hat r^2-R^2)
=\tfrac12(\hat r-R)(\hat r+R)\ge\tfrac{\hat r}{2}(\hat r-R)$ (valid as
$R\ge0$),
\[
   a\ \ge\ |B_{\hat r}(z)\setminus\Etil|
   \ \ge\ \int_{\set{R<\hat r}}\!\int_{R(\theta)}^{\hat r}\rho\,\d\rho\,\d\theta
   \ \ge\ \frac{\hat r}{2}\int_{S^1}(\hat r-R(\theta))_+\,\d\theta,
\]
where $|B_{\hat r}(z)\setminus\Etil|\le|\Etil\triangle B_{\hat r}(z)|=a$ by
\eqref{eq:a-def}. Hence $\int_{S^1}(\hat r-R)_+\,\d\theta\le 2a/\hat r$, and
since $R\ge\hat r-(\hat r-R)_+$,
\[
   \int_{S^1}R(\theta)\,\d\theta\ \ge\ 2\pi\hat r-\int_{S^1}(\hat r-R)_+\,\d\theta
   \ \ge\ 2\pi\hat r-\frac{2a}{\hat r}.
\]
Combining with \eqref{eq:Pr-lower} and $P(\Etil)=2\pi\hat r+\Exc$,
\[
   P(\Etil)-P_r\ \le\ P(\Etil)-\int_{S^1}R\,\d\theta
   \ \le\ (2\pi\hat r+\Exc)-\Bigl(2\pi\hat r-\frac{2a}{\hat r}\Bigr)
   =\Exc+\frac{2a}{\hat r}.
\]
The final inequality in \eqref{eq:radial-profile} uses $\hat r\ge\tfrac12$.
\end{proof}

\subsection{The trapping theorem}

\begin{theorem}[Trapping annulus]\label{thm:annulus}
Let $E\subseteq\R^2$ be a bounded, open, connected set, equal to the good
representative described above \textup{(}so $\partial E$ is the topological
boundary, $\HH^1$-rectifiable with $\HH^1(\partial E)=P(E)$, and the
trichotomy \eqref{eq:trichotomy} holds\textup{)}, with no holes
\textup{(}$\R^2\setminus\overline E$ connected\textup{)}, area
$|E|=\pi\hat r^2$, isoperimetric deficit $\diso(E)$, perimeter excess
$\Exc=P(E)-2\pi\hat r=2\pi\hat r\,\diso(E)$, and ball distance
$a=|E\triangle B_{\hat r}(z)|$ at the optimal centre $z$. Assume $z\in E$,
and let $\rho_i,\rho_e$ be as in \eqref{eq:radii-def}. Then
\begin{equation}\label{eq:annulus}
   B_{\rho_i}(z)\subseteq E\subseteq \overline{B_{\rho_e}(z)},
   \qquad
   \rho_i\le\hat r\le\rho_e,
   \qquad
   \eta:=\rho_e-\rho_i\ \le\ C\,(\Exc+a)^{1/2},
\end{equation}
with $C$ an absolute constant. In particular, for the inset $\Etil$ of
\cref{sec:inset} \textup{(}$\dT\le\delta_\star$, $\hat r\in[\tfrac12,1]$, and
the scales \eqref{eq:ann-scales}\textup{)},
\begin{equation}\label{eq:annulus-rate}
   \eta=\rho_e-\rho_i\ \le\ C\,(\Exc+a)^{1/2}\ \le\ C\,\Bigl(\frac{\dT}{\tau}
   \Bigr)^{1/4},
\end{equation}
which at the optimal choice $\tau=\sqrt{\dT}$ of \cref{sec:main} reads
$\eta\le C\,\dT^{1/8}$.
\end{theorem}

\begin{proof}
The inclusions and $\rho_i\le\hat r\le\rho_e$ are \cref{lem:radii}; it
remains to prove the width bound. We give the estimate in three steps and
assemble.

\emph{Step 1 \textup{(}geometry $\Rightarrow$ angular perimeter\textup{)}.}
By \cref{lem:Ncross}, $N(\rho)\ge2$ for every $\rho\in(\rho_i,\rho_e)$, so
\cref{lem:angular-coarea} gives
\begin{equation}\label{eq:step1}
   P_\theta=\int_0^\infty N(\rho)\,\d\rho
   \ \ge\ \int_{\rho_i}^{\rho_e}2\,\d\rho
   =2(\rho_e-\rho_i)=2\eta .
\end{equation}

\emph{Step 2 \textup{(}Cauchy--Schwarz\textup{)}.} Write $P:=P(E)$. By the
Cauchy--Schwarz inequality on $\partial E$ (with the constant function $1$),
\[
   P_\theta^2=\Bigl(\int_{\partial E}|\nu_\theta|\,\d\HH^1\Bigr)^2
   \le P\int_{\partial E}\nu_\theta^2\,\d\HH^1
   = P\int_{\partial E}\bigl(1-\nu_r^2\bigr)\,\d\HH^1
   = P\Bigl(P-\int_{\partial E}\nu_r^2\,\d\HH^1\Bigr),
\]
using $\nu_r^2+\nu_\theta^2=1$. A second application of Cauchy--Schwarz gives
$\int_{\partial E}\nu_r^2\,\d\HH^1\ge\bigl(\int_{\partial E}|\nu_r|\,\d\HH^1
\bigr)^2/P=P_r^2/P$, hence
\begin{equation}\label{eq:step2}
   P_\theta^2\ \le\ P\Bigl(P-\frac{P_r^2}{P}\Bigr)=P^2-P_r^2
   =(P-P_r)(P+P_r)\ \le\ 2P\,(P-P_r),
\end{equation}
the last step using $0\le P_r\le P$.

\emph{Step 3 \textup{(}radial-profile bound\textup{)}.} By
\cref{lem:radial-profile}, $P-P_r\le\Exc+2a/\hat r$.

\emph{Assembly.} Combining \eqref{eq:step1}, \eqref{eq:step2}, Step 3 and
$P=2\pi\hat r+\Exc$,
\[
   4\eta^2\ \le\ P_\theta^2\ \le\ 2P\,(P-P_r)
   \ \le\ 2\,(2\pi\hat r+\Exc)\Bigl(\Exc+\frac{2a}{\hat r}\Bigr).
\]
Using $\hat r\le1$ and $\Exc\le C(\dT/\tau)\le C$ we have $2\pi\hat r+\Exc\le
2\pi+C\le C$, and using $\hat r\ge\tfrac12$ we have $2a/\hat r\le 4a$;
therefore
\[
   4\eta^2\ \le\ C\,(\Exc+4a)\ \le\ C\,(\Exc+a),
   \qquad\text{i.e.}\qquad
   \eta\ \le\ C\,(\Exc+a)^{1/2},
\]
with $C$ absolute. This is the width bound in \eqref{eq:annulus}. Finally,
inserting \eqref{eq:ann-scales}, $\Exc+a\le C(\dT/\tau)+C(\dT/\tau)^{1/2}\le
C(\dT/\tau)^{1/2}$ for $\dT/\tau\le1$, so $\eta\le C(\dT/\tau)^{1/4}$, which
is \eqref{eq:annulus-rate}; at $\tau=\sqrt\dT$ this is $C\dT^{1/8}$.
\end{proof}

\begin{remark}[A posteriori: the annulus is centred near the unit circle]
\label{rem:posterior}
The bound \eqref{eq:annulus-rate} together with $\rho_i\le\hat r\le\rho_e$
yields
\[
   \hat r-\rho_i\le\eta\le C\,(\dT/\tau)^{1/4},
   \qquad
   \rho_e-\hat r\le\eta\le C\,(\dT/\tau)^{1/4},
\]
so $\rho_i\ge\hat r-C(\dT/\tau)^{1/4}\ge\tfrac12$ and $\rho_e\le\hat
r+C(\dT/\tau)^{1/4}\le 2$ once $\dT\le\delta_\star$ is small enough (recall
$\hat r\in[\tfrac12,1]$ and $\dT/\tau\to0$). These bounds (in particular
$\rho_i\ge\tfrac12$) are exactly the inputs used in \cref{sec:fold,sec:cut};
they are \emph{outputs} of the trapping, not assumptions of it.
\end{remark}

\begin{remark}[On the hypothesis $z\in\Etil$, and why it holds]\label{rem:zinE}
The sole structural input of \cref{thm:annulus} beyond connectedness and the
no-hole property is that the optimal centre $z$ is an interior point of
$\Etil$ (equivalently $\rho_i>0$): it is used to define the angular
projection $\Theta$ on $\partial\Etil$ in \eqref{eq:radial-weighted} and to
start the radius at $0$ in Steps~1--2 of \cref{lem:Ncross}. We verify it
\emph{rigorously} for the inset, splitting on where $z$ sits relative to the
good representative \eqref{eq:trichotomy}.

\emph{$z$ is not in the open exterior.} If $z\in\R^2\setminus\overline\Etil$,
let $\delta:=\dist(z,\overline\Etil)>0$. Then $B_{\min(\delta,\hat r)}(z)
\cap\Etil=\varnothing$, so, using $|\Etil|=\pi\hat r^2$ and
\eqref{eq:a-def},
\[
   \pi\min(\delta,\hat r)^2\le|B_{\hat r}(z)\setminus\Etil|
   =\tfrac12\,a\le C\,(\dT/\tau)^{1/2}.
\]
Since $a\le2|\Etil|$ forbids $\delta>\hat r$, this gives $\delta\le
C(\dT/\tau)^{1/4}$. For $\dT/\tau\le\delta_\star$ small enough this is
$<\hat r/2$, so $z$ lies within $\hat r/2$ of $\overline\Etil$; after the
harmless recentring described next, $z\in\overline\Etil$.

\emph{Recentring \textup{(}if needed\textup{)}.} If $z\notin\overline\Etil$
or $z\in\partial\Etil$, replace $z$ by a point $z_0\in\Etil$ with
$|z_0-z|\le\delta':=C(\dT/\tau)^{1/4}$ (it exists: $|B_{\hat r}(z)\cap\Etil|
=\pi\hat r^2-a/2>0$, so $\Etil$ meets every disc about $z$ of area exceeding
$a/2$, in particular some $B_{r_0}(z)$ with $r_0\le C(\dT/\tau)^{1/4}$).
Measuring the ball distance at $z_0$ costs only
\[
   a_0:=|\Etil\triangle B_{\hat r}(z_0)|
   \le a+|B_{\hat r}(z)\triangle B_{\hat r}(z_0)|
   \le a+4\pi\hat r\,|z-z_0|\le C(\dT/\tau)^{1/4},
\]
since $B_{\hat r}(z)\triangle B_{\hat r}(z_0)\subseteq B_{\hat r+|z-z_0|}(z)
\setminus B_{\hat r-|z-z_0|}(z)$ has area $4\pi\hat r|z-z_0|$. With this
$z_0\in\Etil$ in place of $z$, \cref{thm:annulus} applies and yields
$\rho_e-\rho_i\le C(\Exc+a_0)^{1/2}$, hence the trapping with absolute
constant; only the \emph{rate} would be the slightly weaker
$(\dT/\tau)^{1/8}$ \emph{in this degenerate case}.

\emph{Self-consistency, and honest status.} Whenever $z\in\Etil$ (the
generic case, $a_0=a$) the conclusion \eqref{eq:annulus-rate} with
\cref{rem:posterior} forces $\rho_i\ge\hat r-C(\dT/\tau)^{1/4}\ge\tfrac12>0$,
confirming $z\in\Etil$ a posteriori with the sharp rate. The one statement
not reduced to a tautology is that the \emph{exact} FMP optimum is interior
(rather than on $\partial\Etil$); we have not proved this from first
principles, but we have shown that an interior centre exists within
$C(\dT/\tau)^{1/4}$ of it, with $a_0$ of the same order. Consequently
\cref{thm:annulus} holds verbatim for any $z\in\Etil$, the trapping with an
\emph{absolute constant} holds unconditionally, and the only quantity that
could degrade---the \emph{exponent}, in a non-generic boundary configuration
of the optimal centre---is flagged here and is, by the a-posteriori bound,
not actually triggered under $\dT\le\delta_\star$.
\end{remark}

\begin{remark}[What is, and is not, used]\label{rem:summary}
The trapping width $\eta\le C(\Exc+a)^{1/2}$ is proved with an
\emph{absolute} constant and \emph{no} second-order information on
$\partial\Etil$: no interior or exterior ball condition, no curvature bound,
no harmonic oscillation estimate. The ingredients are exactly: the
connectedness and no-hole property of $\Etil$ (\cref{lem:noholes}); the
coarea formula on a rectifiable curve (\cref{thm:coarea}); the isoperimetric
excess $\Exc$ and the FMP closeness $a$ (\cref{thm:fmp}, via
\eqref{eq:ann-scales}). Notably, the Cauchy--Schwarz/Serrin defect $D_H$ is
\emph{not} used in this section. The single arithmetic of the proof is the
chain $4\eta^2\le P_\theta^2\le2P(P-P_r)\le2P(\Exc+2a/\hat r)$, in which the
geometric lower bound (two crossings per circle) is converted into the
deficit by one Cauchy--Schwarz and one elementary radial-profile estimate.
\end{remark}

\section{Fold content and the radial graph}\label{sec:fold}

The trapping annulus of \cref{thm:annulus} confines $\partial\Etil$ to a
thin shell $\{\rho_i\le\abs{x-z}\le\rho_e\}$ of width $\eta=\rho_e-\rho_i$,
so the inset has \emph{small amplitude}. To convert this into a radial
\emph{graph} about the centre $z$ we must rule out radial folds: rays from
$z$ meeting $\partial\Etil$ more than once. This section bounds the total
\emph{fold content} --- the obstruction to being a graph --- by the
geometric deficit, with \emph{no} square-root loss from the quantitative
isoperimetric inequality, and then realises $\Etil$ as a measurable radial
graph up to a controlled symmetric difference.

Throughout, $\Etil$ is the hole-filled inset of \cref{def:inset,lem:fill}
(\cref{conv:fill}), connected and globally hole-free, so the global
exterior-connectedness needed to invoke the trapping annulus of
\cref{thm:annulus} holds unconditionally (\cref{lem:fill}\ref{it:fill-top}).
We write
with $\mu_E:=\abs\Etil$, $\hat r:=\sqrt{\mu_E/\pi}$, isoperimetric deficit
$\diso:=\diso(\Etil)$, perimeter excess $\Exc:=\Per(\Etil)-2\pi\hat r
=2\pi\hat r\,\diso$, and (\cref{thm:annulus}) trapping
\begin{equation}\label{eq:fold-annulus}
   B_{\rho_i}(z)\subset\Etil\subset B_{\rho_e}(z),\qquad
   \rho_i\le\hat r\le\rho_e,\qquad
   \eta:=\rho_e-\rho_i\le C\,(\Exc+a)^{1/2},
\end{equation}
where $z$ is the optimal centre and $a:=\abs{\Etil\triangle B_{\hat r}(z)}$.
In particular $z$ lies in the interior of $\Etil$ (indeed
$B_{\rho_i}(z)\subset\Etil$), and $\partial\Etil$ avoids the singular point:
\begin{equation}\label{eq:fold-away}
   \partial\Etil\subset\{x:\rho_i\le\abs{x-z}\le\rho_e\}.
\end{equation}
We write $e_r:=(x-z)/\abs{x-z}$, $e_\theta:=e_r^\perp$ for the unit radial
and angular vectors at $x\ne z$, $\nu$ for the outward unit normal on the
reduced boundary $\partial^*\Etil$, and
\[
   \nu_r:=\nu\cdot e_r
\]
for the radial part of $\nu$. (As is standard for sets of finite perimeter,
all boundary integrals are over the reduced boundary $\partial^*\Etil$, on
which $\nu$ is defined $\HH^1$-a.e.; we write $\partial\Etil$ for it
throughout, noting that for the hole-free finite-perimeter set $\Etil$ the
reduced and essential boundaries coincide $\HH^1$-a.e.\ by the De
Giorgi--Federer structure theorem.) For $\theta\in S^1$ let
\begin{equation}\label{eq:fold-M}
   M(\theta):=\HH^0\bigl(\partial\Etil\cap R_\theta\bigr),\qquad
   R_\theta:=\{z+\rho e^{i\theta}:\rho>0\}
\end{equation}
be the number of crossings of $\partial\Etil$ by the open ray in direction
$\theta$. By the one-dimensional structure of slices of finite-perimeter
sets \cite[\S2.4, \S3.11]{AFP2000}, for a.e.\ $\theta$ the radial slice
\[
   \Etil_\theta:=\{\rho>0:z+\rho e^{i\theta}\in\Etil\}
\]
is, up to an $\HH^1$-null set, a finite union of intervals; since
$B_{\rho_i}(z)\subset\Etil\subset B_{\rho_e}(z)$, the innermost interval
starts at $0$ and all endpoints lie in $[\rho_i,\rho_e]$, so
\begin{equation}\label{eq:fold-slice}
   \Etil_\theta=(0,r_1)\cup(r_2,r_3)\cup\cdots\cup(r_{2m},r_{2m+1}),
   \quad \rho_i\le r_1\le\cdots\le r_{2m+1}\le\rho_e,
   \quad M(\theta)=2m(\theta)+1.
\end{equation}
The intervals $(r_{2i},r_{2i+1})$, $i\ge1$, are the \emph{caps};
$m(\theta)$ is the number of caps on the ray, and $\Etil$ is a radial graph
about $z$ exactly when $M\equiv1$ (no caps).

\subsection{Point-source flux identities}

\begin{lemma}[Two-flux identities]\label{lem:two-flux}
Let $E\subset\R^2$ be a set of finite perimeter and suppose that, for some
$r_0>0$,
\[
   B_{r_0}(z)\subset E,\qquad
   \partial^* E\subset\{x:\abs{x-z}\ge r_0\}.
\]
With $e_r=(x-z)/\abs{x-z}$, $\nu_r=\nu\cdot e_r$, and
$M(\theta)$ as in \eqref{eq:fold-M}, the following hold:
\begin{subequations}\label{eq:two-flux}
\begin{align}
   \textup{(i)}\quad
   &\int_{\partial E}\frac{\abs{\nu_r}}{\abs{x-z}}\,d\HH^1
   =\int_{S^1}M(\theta)\,d\theta,
   \label{eq:two-flux-i}\\[2pt]
   \textup{(ii)}\quad
   &\int_{\partial E}\frac{\nu_r}{\abs{x-z}}\,d\HH^1=2\pi,
   \label{eq:two-flux-ii}\\[2pt]
   \textup{(iii)}\quad
   &\int_{\partial E}\nu_r\,d\HH^1
   =\int_E\frac{dx}{\abs{x-z}}=:I_E .
   \label{eq:two-flux-iii}
\end{align}
\end{subequations}
We refer to the three identities collectively as \eqref{eq:two-flux}.
\end{lemma}

\begin{proof}
We may translate so that $z=0$; then $e_r=x/\abs x$ and $\abs{x-z}=\abs x$.
Set $A:=\{x:r_0\le\abs x\}$, a closed set containing $\partial^*E$ on which
all three integrands are smooth and bounded. Since $B_{r_0}(0)\subset E$,
for every $0<\eps<r_0$ the circle $\partial B_\eps(0)$ lies in the
measure-theoretic interior $E^{(1)}$.

\emph{(i) The angular coarea identity.}
Let $\Theta(x):=x/\abs x\in S^1$ be the angular projection.  On the annulus
$A$ this map is Lipschitz.  The reduced boundary $\partial^*E$ is a
countably $1$-rectifiable set contained in $A$; let $\tau$ denote a unit
tangent to $\partial^*E$ (defined $\HH^1$-a.e., orthogonal to $\nu$). The
tangential Jacobian of $\Theta$ along $\partial^*E$ is
\[
   J_{\partial^*E}\Theta
   =\frac{\abs{e_\theta\cdot\tau}}{\abs x}.
\]
Since $\{e_r,e_\theta\}$ and $\{\nu,\tau\}$ are each orthonormal bases of
$\R^2$ with $e_r\perp e_\theta$ and $\nu\perp\tau$, the angle between
$e_\theta$ and $\tau$ equals the angle between $e_r$ and $\nu$; hence
$\abs{e_\theta\cdot\tau}=\abs{e_r\cdot\nu}=\abs{\nu_r}$, so
\[
   J_{\partial^*E}\Theta=\frac{\abs{\nu_r}}{\abs x}.
\]
The coarea formula for the Lipschitz map
$\Theta\colon\partial^*E\to S^1$ on the $1$-rectifiable set
$\partial^*E$ \cite[Theorem~2.93]{AFP2000} gives
\[
   \int_{\partial^*E}\frac{\abs{\nu_r}}{\abs x}\,d\HH^1
   =\int_{\partial^*E}J_{\partial^*E}\Theta\,d\HH^1
   =\int_{S^1}\HH^0\bigl(\partial^*E\cap \Theta^{-1}(\theta)\bigr)\,d\theta
   =\int_{S^1}M(\theta)\,d\theta,
\]
because $\Theta^{-1}(\theta)\cap A=R_\theta\cap A\supset
\partial^*E\cap R_\theta$ and $\partial^*E\subset A$. This is
\eqref{eq:two-flux-i}.

\emph{(ii) The point-source field, singularity excised.}
Consider the vector field $X(x):=x/\abs x^2$, defined and smooth on
$\R^2\setminus\{0\}$, with
\[
   \divg X=\frac{2}{\abs x^2}-\frac{2}{\abs x^2}=0
   \qquad(x\ne0)
\]
(directly: $\partial_i(x_i\abs x^{-2})=\abs x^{-2}+x_i(-2x_i\abs x^{-4})$,
summing $i=1,2$ gives $2\abs x^{-2}-2\abs x^{-2}=0$). On $\partial^*E$ we
have $X\cdot\nu=(x\cdot\nu)/\abs x^2=\nu_r/\abs x$. Fix $0<\eps<r_0$ and
apply the Gauss--Green theorem for sets of finite perimeter
\cite[Theorem~5.16]{EvansGariepy2015} to the field $X$, which is $C^1$ and
bounded on the closed set $\{\abs x\ge\eps/2\}$, over the finite-perimeter
set $E_\eps:=E\setminus\overline{B_\eps(0)}$. Since $\partial B_\eps\subset
E^{(1)}$ (shown above), the reduced boundary $\partial^*E_\eps$ splits
($\HH^1$-a.e.) into $\partial^*E$ (with outward normal $\nu$, and on which
$X$ is integrable since $\partial^*E\subset A$) and the inner circle
$\partial B_\eps$ (with outward-for-$E_\eps$ normal $-x/\eps$):
\[
   0=\int_{E_\eps}\divg X\,dx
   =\int_{\partial^*E}X\cdot\nu\,d\HH^1
   +\int_{\partial B_\eps}X\cdot\Bigl(-\frac{x}{\eps}\Bigr)\,d\HH^1 .
\]
On $\partial B_\eps$, $X=x/\eps^2$ and $X\cdot(-x/\eps)=-\abs x^2/\eps^3
=-1/\eps$, so the second term equals $-\eps^{-1}\HH^1(\partial B_\eps)
=-\eps^{-1}\cdot2\pi\eps=-2\pi$, independent of $\eps$. Hence
\[
   \int_{\partial^*E}\frac{\nu_r}{\abs x}\,d\HH^1
   =\int_{\partial^*E}X\cdot\nu\,d\HH^1=2\pi,
\]
which is \eqref{eq:two-flux-ii}. (No limit in $\eps$ is even needed: the
left side does not depend on $\eps$, and the inner-circle flux is exactly
$2\pi$ for every admissible $\eps$.)

\emph{(iii) The radial field, singularity integrable.}
Consider $Y(x):=e_r=x/\abs x$, smooth on $\R^2\setminus\{0\}$, with
\[
   \divg Y=\sum_i\partial_i\frac{x_i}{\abs x}
   =\sum_i\Bigl(\frac1{\abs x}-\frac{x_i^2}{\abs x^3}\Bigr)
   =\frac2{\abs x}-\frac{\abs x^2}{\abs x^3}=\frac1{\abs x}
   \qquad(x\ne0).
\]
Again fix $0<\eps<r_0$ and apply Gauss--Green to the $C^1$ bounded field $Y$
over $E_\eps=E\setminus\overline{B_\eps(0)}$:
\[
   \int_{E_\eps}\frac{dx}{\abs x}
   =\int_{E_\eps}\divg Y\,dx
   =\int_{\partial^*E}Y\cdot\nu\,d\HH^1
   +\int_{\partial B_\eps}Y\cdot\Bigl(-\frac x\eps\Bigr)\,d\HH^1 .
\]
On $\partial B_\eps$, $Y=x/\eps$ and $Y\cdot(-x/\eps)=-\abs x^2/\eps^2=-1$,
so the inner term is $-\HH^1(\partial B_\eps)=-2\pi\eps$. Thus
\[
   \int_{\partial^*E}\nu_r\,d\HH^1
   =\int_{E_\eps}\frac{dx}{\abs x}+2\pi\eps .
\]
Now let $\eps\downarrow0$. The left side is fixed. On the right,
$\abs x^{-1}\in L^1_{\mathrm{loc}}(\R^2)$ (in polar coordinates
$\int_{B_R}\abs x^{-1}\,dx=2\pi R<\infty$) and $E$ has finite measure, so
$\abs x^{-1}\mathbf1_E\in L^1(\R^2)$; by dominated convergence
$\int_{E_\eps}\abs x^{-1}\,dx\to\int_E\abs x^{-1}\,dx=I_E$, while
$2\pi\eps\to0$. Therefore $\int_{\partial^*E}\nu_r\,d\HH^1=I_E$, which is
\eqref{eq:two-flux-iii}. The constant $I_E$ is finite because
$\abs x^{-1}\mathbf1_E\in L^1$.
\end{proof}

\begin{remark}[On the singularity]\label{rem:singularity}
The two divergence-theorem identities are genuinely different at the source.
In \eqref{eq:two-flux-ii} the field $X=x/\abs x^2$ has a non-integrable
$\abs x^{-1}$ growth at $0$ and a non-zero residual flux $2\pi$ through any
small circle; excising $B_\eps$ is mandatory, and the $2\pi$ is precisely
that residue (the planar Green-function source). In \eqref{eq:two-flux-iii}
the field $Y=e_r$ is bounded, $\divg Y=\abs x^{-1}$ is integrable, and the
inner-circle flux $2\pi\eps\to0$ vanishes; the excision is a convenience and
the resulting $I_E$ is a finite, well-defined number. We invoke
\cref{lem:two-flux} for $E=\Etil$, whose reduced boundary satisfies the
hypothesis $\partial^*\Etil\subset\{\abs{x-z}\ge\rho_i\}$ with $r_0=\rho_i$
by \eqref{eq:fold-away}, and whose interior contains $z$ since
$B_{\rho_i}(z)\subset\Etil$.
\end{remark}

\subsection{Sharp fold content}

The fold content is the deviation of $M$ from the graph value $1$. We bound
its total mass with no quantitative-isoperimetric square-root loss; this is
the crux of the sharp reduction.

\begin{proposition}[Sharp fold content]\label{prop:fold}
With the inset $\Etil$ and the trapping data \eqref{eq:fold-annulus},
\begin{equation}\label{eq:foldsharp}
   \int_{S^1}\bigl(M(\theta)-1\bigr)\,d\theta
   \ =\ 2\!\int_{\partial\Etil\cap\{\nu_r<0\}}\!
        \frac{\abs{\nu_r}}{\abs{x-z}}\,d\HH^1
   \ \le\ \frac1{\rho_i}\Bigl(\Exc+\frac{a\,\eta}{\rho_i\hat r}\Bigr)
   \ \le\ C\bigl(\Exc+a\eta\bigr),
\end{equation}
with $C$ absolute. In particular the integrand is non-negative, so
$M\ge1$ a.e., and
\begin{equation}\label{eq:foldangles}
   \bigl|\{\theta\in S^1:M(\theta)\ge3\}\bigr|
   \ \le\ \tfrac12\!\int_{S^1}\bigl(M-1\bigr)\,d\theta
   \ \le\ C\bigl(\Exc+a\eta\bigr).
\end{equation}
\end{proposition}

\begin{proof}
\emph{The fold identity.}
Subtract \eqref{eq:two-flux-ii} from \eqref{eq:two-flux-i} of
\cref{lem:two-flux} (legitimate by \cref{rem:singularity}); since
$\int_{S^1}1\,d\theta=2\pi$ equals the right-hand side of
\eqref{eq:two-flux-ii},
\[
   \int_{S^1}\bigl(M-1\bigr)\,d\theta
   =\int_{\partial\Etil}\frac{\abs{\nu_r}-\nu_r}{\abs{x-z}}\,d\HH^1
   =2\!\int_{\partial\Etil\cap\{\nu_r<0\}}\!
       \frac{\abs{\nu_r}}{\abs{x-z}}\,d\HH^1\ \ge0,
\]
using $\abs{\nu_r}-\nu_r=0$ where $\nu_r\ge0$ and $=2\abs{\nu_r}$ where
$\nu_r<0$. This is the first equality of \eqref{eq:foldsharp}. (Pointwise
$M(\theta)=2m(\theta)+1\ge1$ a.e.\ by the slice structure
\eqref{eq:fold-slice}, since every slice contains $(0,r_1)$; the content of
the identity is the \emph{global} bound on the excess $\int(M-1)$, which is
exactly twice the inward-pointing radial flux.)

\emph{From the radial fold flux to $P_r-I_{\Etil}$.}
On the folded part $\partial\Etil\cap\{\nu_r<0\}$ we bound
$\abs{x-z}\ge\rho_i$ (by \eqref{eq:fold-away}), so
\begin{equation}\label{eq:fold-step1}
   \int_{S^1}\bigl(M-1\bigr)\,d\theta
   \ \le\ \frac2{\rho_i}\!\int_{\partial\Etil\cap\{\nu_r<0\}}\!
          \abs{\nu_r}\,d\HH^1
   \ =\ \frac{2F_r}{\rho_i},\qquad
   F_r:=\!\int_{\partial\Etil\cap\{\nu_r<0\}}\!\abs{\nu_r}\,d\HH^1 .
\end{equation}
Write $P_r:=\int_{\partial\Etil}\abs{\nu_r}\,d\HH^1$ for the total radial
perimeter. Splitting according to the sign of $\nu_r$,
\[
   \int_{\partial\Etil}\nu_r\,d\HH^1
   =\int_{\{\nu_r\ge0\}}\abs{\nu_r}\,d\HH^1
    -\int_{\{\nu_r<0\}}\abs{\nu_r}\,d\HH^1
   =\bigl(P_r-F_r\bigr)-F_r=P_r-2F_r .
\]
By \eqref{eq:two-flux-iii}, the left side is $I_{\Etil}$, hence
\begin{equation}\label{eq:fold-Fr}
   2F_r=P_r-I_{\Etil} .
\end{equation}

\emph{Upper bound on $P_r$.}
Since $\abs{\nu_r}=\abs{\nu\cdot e_r}\le\abs\nu=1$ pointwise, the radial
perimeter is dominated by the perimeter:
\[
   P_r=\int_{\partial\Etil}\abs{\nu_r}\,d\HH^1
   \le\int_{\partial\Etil}1\,d\HH^1=\Per(\Etil)=2\pi\hat r+\Exc .
\]

\emph{The cancellation: lower bound on $I_{\Etil}$.}
Here is the sharp step. Because $\abs\Etil=\mu_E=\pi\hat r^2
=\abs{B_{\hat r}(z)}$, the indicators have equal mass, so
$\int(\mathbf1_{\Etil}-\mathbf1_{B_{\hat r}(z)})\,dx=0$ and we may subtract
\emph{any} constant from the kernel without changing the integral. Subtract
the constant $1/\hat r$:
\[
   I_{\Etil}-\int_{B_{\hat r}(z)}\frac{dx}{\abs{x-z}}
   =\int\bigl(\mathbf1_{\Etil}-\mathbf1_{B_{\hat r}(z)}\bigr)
        \frac{dx}{\abs{x-z}}
   =\int\bigl(\mathbf1_{\Etil}-\mathbf1_{B_{\hat r}(z)}\bigr)
        \Bigl(\frac1{\abs{x-z}}-\frac1{\hat r}\Bigr)\,dx,
\]
the last equality because $\int(\mathbf1_{\Etil}-\mathbf1_{B_{\hat r}(z)})
\cdot\hat r^{-1}\,dx=\hat r^{-1}(\abs\Etil-\abs{B_{\hat r}(z)})=0$. In polar
coordinates,
$\int_{B_{\hat r}(z)}\abs{x-z}^{-1}\,dx
=\int_0^{\hat r}\!\!\int_{S^1}s^{-1}\,s\,d\theta\,ds=2\pi\hat r$, so
\begin{equation}\label{eq:fold-Idiff}
   I_{\Etil}-2\pi\hat r
   =\int_{\Etil\triangle B_{\hat r}(z)}
       \pm\Bigl(\frac1{\abs{x-z}}-\frac1{\hat r}\Bigr)\,dx,
\end{equation}
the integrand being supported on the symmetric difference
$\Etil\triangle B_{\hat r}(z)$ (with sign $+$ on $\Etil\setminus
B_{\hat r}(z)$ and $-$ on $B_{\hat r}(z)\setminus\Etil$). Now
$\Etil\triangle B_{\hat r}(z)$ lies in the closed annulus
$\{\rho_i\le\abs{x-z}\le\rho_e\}$: indeed
$B_{\rho_i}(z)\subset\Etil\cap B_{\hat r}(z)$ (as $\rho_i\le\hat r$), so
neither $\Etil$ nor $B_{\hat r}(z)$ differs from the other inside
$B_{\rho_i}(z)$; and $\Etil\cup B_{\hat r}(z)\subset B_{\rho_e}(z)$ (as
$\hat r\le\rho_e$ and $\Etil\subset B_{\rho_e}(z)$), so neither differs
outside $B_{\rho_e}(z)$. On that annulus, with both $\abs{x-z}$ and
$\hat r$ in $[\rho_i,\rho_e]$,
\begin{equation}\label{eq:fold-kernel}
   \Bigl|\frac1{\abs{x-z}}-\frac1{\hat r}\Bigr|
   =\frac{\bigl|\hat r-\abs{x-z}\bigr|}{\hat r\,\abs{x-z}}
   \le\frac{\rho_e-\rho_i}{\hat r\,\rho_i}=\frac{\eta}{\rho_i\hat r}.
\end{equation}
Therefore, taking absolute values in \eqref{eq:fold-Idiff} and using
\eqref{eq:fold-kernel},
\begin{equation}\label{eq:Icancel}
   \bigl|I_{\Etil}-2\pi\hat r\bigr|
   \ \le\ \frac{\eta}{\rho_i\hat r}\,
          \bigl|\Etil\triangle B_{\hat r}(z)\bigr|
   \ =\ \frac{a\,\eta}{\rho_i\hat r}.
\end{equation}

\emph{Conclusion.}
Combining \eqref{eq:fold-Fr} with the bound on $P_r$ and the lower bound
$I_{\Etil}\ge2\pi\hat r-a\eta/(\rho_i\hat r)$ from \eqref{eq:Icancel},
\[
   2F_r=P_r-I_{\Etil}
   \le\bigl(2\pi\hat r+\Exc\bigr)
      -\Bigl(2\pi\hat r-\frac{a\eta}{\rho_i\hat r}\Bigr)
   =\Exc+\frac{a\eta}{\rho_i\hat r}.
\]
Inserting into \eqref{eq:fold-step1} yields the middle inequality of
\eqref{eq:foldsharp}. Finally, in the optimised inset regime of
\cref{thm:reduction}, \eqref{eq:close-inset} gives
$\mu_E=\pi-O(\sqrt{\dT})$, hence $\hat r\ge\tfrac34$ after shrinking
$\delta_\star$, and \cref{thm:annulus} gives
$\eta\le C\dT^{1/8}\le\tfrac14$. Thus $\rho_i\ge\hat r-\eta\ge\tfrac12$
(using
$\rho_e=\rho_i+\eta$ and $\rho_i\le\hat r\le\rho_e$). Therefore
$1/\rho_i\le2$ and $1/(\rho_i^2\hat r)\le1/(\tfrac14\cdot\tfrac34)\le6$,
giving the absolute bound $C(\Exc+a\eta)$. The fold-angle estimate \eqref{eq:foldangles} follows
because $M-1\ge2$ on $\{M\ge3\}$ (as $M=2m+1$ is odd, $M\ge3$ means
$m\ge1$), so $2\abs{\{M\ge3\}}\le\int_{\{M\ge3\}}(M-1)\le\int_{S^1}(M-1)$.
\end{proof}

\begin{remark}[Why the factor $\eta$ matters: the improvement over the
naive bound]\label{rem:eta-matters}
The naive estimate pulls $\abs{x-z}^{-1}\le\rho_i^{-1}$ out of the integral
defining $I_{\Etil}-2\pi\hat r$ before subtracting any constant, giving
\[
   \bigl|I_{\Etil}-2\pi\hat r\bigr|
   =\Bigl|\int\bigl(\mathbf1_{\Etil}-\mathbf1_{B_{\hat r}(z)}\bigr)
        \tfrac1{\abs{x-z}}\Bigr|
   \le\frac1{\rho_i}\,\abs{\Etil\triangle B_{\hat r}(z)}
   =\frac{a}{\rho_i}\le Ca,
\]
hence $\int(M-1)\le C(\Exc+a)$. The two estimates differ by the factor
$\eta=\rho_e-\rho_i$ on the $a$-term. This is decisive. By the quantitative
isoperimetric inequality (\cref{thm:fmp}), $a\le C_{\mathrm F}\sqrt\diso\,
\mu_E$ carries an unavoidable \emph{square-root} of the isoperimetric
deficit; with $\diso\sim\dT/\tau\sim\sqrt\dT$ at the boundary scale this
gives only $a\sim\dT^{1/4}$, and the naive fold content
$\Exc+a\sim\dT^{1/4}$ would cap the whole argument at $\dT^{1/4}$ closeness.
The cancellation \eqref{eq:Icancel} replaces $a$ by $a\eta$; since
$\eta\le C(\Exc+a)^{1/2}\le Ca^{1/2}\sim\dT^{1/8}$, the product
$a\eta\sim\dT^{3/8}$, and through \cref{cor:gd} the graph defect becomes
$\gd\sim\eta\,a\eta\sim\dT^{1/8}\cdot\dT^{3/8}=\dT^{1/2}$, matching the
closeness $O(\sqrt\dT)$. The single factor $\eta$ --- the legitimacy of
subtracting the constant $1/\hat r$, because $\abs\Etil=\abs{B_{\hat r}(z)}$,
together with the confinement of $\Etil\triangle B_{\hat r}(z)$ to the thin
$\eta$-annulus --- is exactly what removes the FMP square-root from the fold
content and makes the reduction sharp.
\end{remark}

\subsection{The graph defect}

The \emph{graph defect} of $\Etil$ about $z$,
\[
   \gd(\Etil):=\inf_{\rho\colon S^1\to(0,\infty)}
      \bigl|\Etil\,\triangle\,\{z+r\omega:0\le r<\rho(\omega)\}\bigr|,
\]
measures how far $\Etil$ is, in $L^1$, from being a radial graph about $z$.
The fold content controls it directly.

\begin{corollary}[Graph defect]\label{cor:gd}
The graph defect of $\Etil$ about the centre $z$ satisfies
\begin{equation}\label{eq:gd}
   \gd(\Etil)\ \le\ \rho_e\,\eta\!\int_{S^1}\!m(\theta)\,d\theta
   \ =\ \tfrac12\,\rho_e\,\eta\!\int_{S^1}\!\bigl(M-1\bigr)\,d\theta
   \ \le\ C\,\eta\,\bigl(\Exc+a\eta\bigr),
\end{equation}
with $C$ absolute.
\end{corollary}

\begin{proof}
Use the \emph{core profile} $\rho_0(\theta):=r_1(\theta)\in[\rho_i,\rho_e]$,
the innermost crossing radius from \eqref{eq:fold-slice}; it is defined for
a.e.\ $\theta$. Set
$G_0:=\{z+r e^{i\theta}:0\le r<\rho_0(\theta)\}$. By construction the slice
of $G_0$ in direction $\theta$ is $(0,\rho_0(\theta))=(0,r_1(\theta))
\subset\Etil_\theta$, so $G_0\subset\Etil$ (up to $\HH^1$-null directions),
and the slice of $\Etil\setminus G_0$ is exactly the union of caps
$(r_2,r_3)\cup\cdots\cup(r_{2m},r_{2m+1})$. By the polar-coordinate area
formula (Fubini on $\R^2\setminus\{z\}$),
\[
   \bigl|\Etil\setminus G_0\bigr|
   =\int_{S^1}\sum_{i=1}^{m(\theta)}\int_{r_{2i}}^{r_{2i+1}}\!r\,dr\,d\theta .
\]
On each ray with $m(\theta)\ge1$, all cap endpoints lie in $[\rho_i,\rho_e]$,
so the caps have total radial length $\sum_i(r_{2i+1}-r_{2i})\le
r_{2m+1}-r_2\le\rho_e-\rho_i=\eta$, and each cap sits at radius
$r\le\rho_e$; hence
\[
   \sum_{i=1}^{m(\theta)}\int_{r_{2i}}^{r_{2i+1}}\!r\,dr
   \ \le\ \rho_e\sum_{i=1}^{m(\theta)}(r_{2i+1}-r_{2i})
   \ \le\ \rho_e\,\eta\,\mathbf1_{\{m(\theta)\ge1\}}.
\]
Therefore
\[
   \gd(\Etil)\le\bigl|\Etil\triangle G_0\bigr|=\bigl|\Etil\setminus G_0\bigr|
   \le\rho_e\,\eta\,\bigl|\{\theta:m(\theta)\ge1\}\bigr|
   \le\rho_e\,\eta\!\int_{S^1}m\,d\theta,
\]
where the last step uses $\mathbf1_{\{m\ge1\}}\le m$ (as $m$ is a
non-negative integer). Since $M=2m+1$, $\int_{S^1}m\,d\theta=\tfrac12
\int_{S^1}(M-1)\,d\theta$, which is the displayed equality; and
\cref{prop:fold} together with $\rho_e\le C\hat r\le C$ (for
$\dT\le\delta_\star$) gives the final bound
$\le C\,\eta(\Exc+a\eta)$.
\end{proof}

\subsection{Realization as a measurable radial graph}

The core profile $\rho_0=r_1$ is, a priori, only \emph{measurable} in
$\theta$ (it is defined ray by ray from the one-dimensional slices, with no
modulus of continuity and no bounded variation available --- a cap on a ray
inserts a jump that the slice structure does not control). We do not attempt
to regularise it. Instead we record that the measurable core graph $G_0$ is
\emph{already} an admissible input to the nearly-spherical estimate of
\cref{sec:nearly}, whose hypotheses --- a bounded open star-shaped set
trapped between two concentric circles, with measurable radial profile and
small $L^\infty$ amplitude --- ask for no slope, no continuity, and no
bounded variation. This removes the regularity step entirely.

\begin{proposition}[Measurable radial-graph realization]\label{prop:realization}
Let $\rho_0=r_1\colon S^1\to[\rho_i,\rho_e]$ be the innermost-crossing core
profile of \cref{cor:gd}, and let
\[
   G_0:=\{z+r\omega(\theta):\ \theta\in S^1,\ 0\le r<\rho_0(\theta)\}
\]
be the associated core graph, taken in its open star-shaped representative
\textup{(}equivalently, replace $\rho_0$ by its lower-semicontinuous
envelope along rays, which alters $G_0$ only on an $\HH^2$-null set and hence
changes none of the quantities below\textup{)}. Then:
\begin{enumerate}[label=\textup{(\roman*)},leftmargin=2.4em]
\item\label{it:real-meas} $\rho_0$ is measurable with values in
$[\rho_i,\rho_e]$, and $G_0$ is a bounded open set, star-shaped about $z$,
with
\begin{equation}\label{eq:real-trap}
   B_{\rho_i}(z)\subset G_0\subset B_{\rho_e}(z);
\end{equation}
\item\label{it:real-amp} writing $\hat r_{G_0}:=\sqrt{\abs{G_0}/\pi}\in
[\rho_i,\rho_e]$ and $\varphi_0:=\rho_0/\hat r_{G_0}-1$, the amplitude obeys
\begin{equation}\label{eq:realization}
   \norm{\varphi_0}_{L^\infty(S^1)}\ \le\ \frac{\rho_e-\rho_i}{\rho_i}
   \ =\ \frac{\eta}{\rho_i}\ \le\ 2\eta;
\end{equation}
\item\label{it:real-close} $G_0\subset\Etil$ and
\begin{equation}\label{eq:realization-close}
   \abs{\Etil\,\triangle\,G_0}\ =\ \abs{\Etil\setminus G_0}
   \ \le\ C\,\eta\,(\Exc+a\eta).
\end{equation}
\end{enumerate}
Consequently $G_0$ satisfies the hypotheses of \cref{thm:graphG} on the
amplitude, the trapping ball, and the volume normalisation; only the
quadratic mode-one bound \eqref{eq:mode1hyp} remains to be supplied, which is
done at the optimal centre in \cref{sec:main}. No mollification is performed
and no continuity of $\rho_0$ is claimed or used.
\end{proposition}

\begin{proof}
\emph{\ref{it:real-meas}.} The profile $\rho_0=r_1$ is measurable: it is the
essential infimum, along the ray $R_\theta$, of the radii at which the slice
$\Etil_\theta$ first exits $\Etil$, and $\theta\mapsto r_1(\theta)$ is
measurable because for each $t$ the set $\{\theta:r_1(\theta)>t\}=\{\theta:
(0,t)\subset\Etil_\theta\ \text{up to }\HH^1\text{-null}\}$ is, up to a null
set, the slice projection of the open set $\Etil\cap(B_t(z)\setminus\set z)$,
hence measurable (as in the measurability of $R$ in \eqref{eq:Rdef}). By
\eqref{eq:fold-slice} every slice contains the innermost interval $(0,r_1)$
with $r_1\in[\rho_i,\rho_e]$, so $\rho_0\in[\rho_i,\rho_e]$. Passing to the
lower-semicontinuous envelope of $\rho_0$ along rays makes
$G_0=\{r<\rho_0\}$ open; this changes $\rho_0$, hence $\partial G_0$, only on
an $\HH^1$-null set of directions, so $\abs{G_0}$ and all symmetric
differences are unchanged. Then $B_{\rho_i}(z)\subset G_0$ because every ray
contains $(0,\rho_i)\subset(0,r_1)$, and $G_0\subset B_{\rho_e}(z)$ because
$\rho_0\le\rho_e$; this is \eqref{eq:real-trap}, and star-shapedness about
$z$ is immediate from the radial definition.

\emph{\ref{it:real-amp}.} From \eqref{eq:real-trap},
$\pi\rho_i^2\le\abs{G_0}\le\pi\rho_e^2$, so $\hat r_{G_0}=
\sqrt{\abs{G_0}/\pi}\in[\rho_i,\rho_e]$. For every $\theta$ both
$\rho_0(\theta)$ and $\hat r_{G_0}$ lie in $[\rho_i,\rho_e]$, whence
$\abs{\rho_0-\hat r_{G_0}}\le\rho_e-\rho_i=\eta$ and
$\norm{\varphi_0}_\infty=\norm{\rho_0-\hat r_{G_0}}_\infty/\hat r_{G_0}\le
\eta/\rho_i\le2\eta$, using $\rho_i\ge\tfrac12$ (\cref{rem:posterior}).

\emph{\ref{it:real-close}.} By construction the slice of $G_0$ in direction
$\theta$ is $(0,\rho_0(\theta))=(0,r_1(\theta))\subset\Etil_\theta$, so
$G_0\subset\Etil$ up to $\HH^1$-null directions, and
$\abs{\Etil\triangle G_0}=\abs{\Etil\setminus G_0}$. The bound
\eqref{eq:realization-close} is exactly \cref{cor:gd}: the slice of
$\Etil\setminus G_0$ is the union of caps
$(r_2,r_3)\cup\cdots\cup(r_{2m},r_{2m+1})$, of total area
$\le\rho_e\,\eta\int_{S^1}m\,d\theta=\tfrac12\rho_e\,\eta\int_{S^1}(M-1)\,
d\theta\le C\,\eta(\Exc+a\eta)$.
\end{proof}

\begin{remark}[Scales, and the role in the reduction]\label{rem:fold-scales}
Specialising to the inset at level $\hat v\sim\tau=\sqrt\dT$ (so
$\diso\le C\dT/\tau\le C\sqrt\dT$, $\Exc\le C\sqrt\dT$, and by \cref{thm:fmp}
$a\le C\sqrt\diso\le C\dT^{1/4}$), the annulus width is $\eta\le
C(\Exc+a)^{1/2}\le Ca^{1/2}\le C\dT^{1/8}$, an \emph{absolute} small
amplitude. Then \cref{cor:gd} gives
\[
   \gd(\Etil)\le C\eta(\Exc+a\eta)
   \le C\bigl(\dT^{1/8}\dT^{1/2}+\dT^{1/4}\dT^{1/4}\bigr)
   =C\bigl(\dT^{5/8}+\dT^{1/2}\bigr)\le C\sqrt\dT,
\]
and \cref{prop:realization} exhibits the measurable core graph $G_0$ of
amplitude $\norm{\varphi_0}_\infty\le2\eta\le C\dT^{1/8}$ with
$\abs{\Etil\triangle G_0}\le C\sqrt\dT$. These are precisely the scales fed
into the cutting estimate (\cref{lem:cutdeficit}) and the assembly
(\cref{thm:reduction}): the graph defect balances the level-carried closeness
$\abs{\Om\triangle\Etil}\sim\sqrt\dT$ at the optimal exponent, and the
measurable-profile form of \cref{thm:graphG} consumes $G_0$ directly, with no
intervening regularisation.
\end{remark}

\section{Cutting the caps lowers the deficit}
\label{sec:cut}

The construction of \cref{sec:inset,sec:annulus,sec:fold} produces from
$\Om$ a connected, hole-free inset $\Etil$ trapped in a thin annulus
$B_{\rho_i}(z)\subset\Etil\subset B_{\rho_e}(z)$, together with a radial core
profile whose graph
$G_0=\{z+r\omega:r<\rho_0(\omega)\}$ satisfies $G_0\subset\Etil$ and
$\abs{\Etil\setminus G_0}=\gd(\Etil)$ (\cref{cor:gd}). Manufacturing the
honest graph from $\Etil$ therefore amounts to \emph{excising} the radial
caps $\Etil\setminus G_0$, all of which lie in the collar
$\{\rho_i<\abs{x-z}<\rho_e\}$. The naive worry is that removing material can
only destroy torsion and hence inflate the scale-invariant deficit
$\dstar$. The opposite is true: caps in a thin collar carry very little
torsion, and the equal-area comparison ball shrinks faster than the torsion
is lost, so $\dstar$ \emph{decreases}. This section proves the two facts
that make this precise. They are sign-critical, so we prove them in full and
check the inequalities digit by digit.

Throughout, for a bounded open set $D\subset\R^2$ the \emph{torsion
function} $u_D\in H^1_0(D)$ is the weak solution of $-\Delta u_D=1$ in $D$,
$u_D=0$ on $\partial D$, and $T(D):=\int_D u_D$ is its torsional rigidity;
for a general set of finite measure $D$ we set $T(D):=T(\mathrm{int}\,D)$ on
the open representative produced by our construction, which is all we use.
Recall (\cref{thm:torsion-reg}) that $u_D$ is real-analytic in $D$ and, by
the comparison principle for the Laplacian, $u_D\ge0$. The variational
characterisation we use repeatedly is that $u_D$ minimises the Dirichlet
energy functional $J_D(v):=\tfrac12\int_D\abs{\nabla v}^2-\int_D v$ over
$v\in H^1_0(D)$, with
\begin{equation}\label{eq:torsion-energy}
   \min_{v\in H^1_0(D)}J_D(v)=J_D(u_D)=-\tfrac12\,T(D),
   \qquad
   \int_D\abs{\nabla u_D}^2=\int_D u_D=T(D),
\end{equation}
the second identity being the Euler--Lagrange relation tested against $u_D$
itself.

\subsection{Torsion loss under radial-cap cutting}

The next lemma is the analytic heart of the section: it controls the torsion
lost when one cuts radial caps off a set trapped in a thin collar. Unlike a
naive measure-based bound --- which is false in general, since excising a
zero-area set of positive capacity (a slit) changes the torsion --- the
estimate is variational and exploits the specific nested radial-cap geometry
produced in \cref{sec:fold}.

\begin{lemma}[Torsion loss under radial-cap cutting]\label{lem:torsion-stab}
Let $z\in\R^2$ and $0<\rho_i\le\rho_e$, and put
\begin{equation}\label{eq:tildeeta}
   \tilde\eta:=\tfrac14\bigl(\rho_e^2-\rho_i^2\bigr).
\end{equation}
Let $G\subset E\subset\R^2$ be bounded open sets with
\begin{equation}\label{eq:collar-hyp}
   B_{\rho_i}(z)\subset G\subset E\subset B_{\rho_e}(z),\qquad
   E\setminus G\subset\{\rho_i<\abs{x-z}<\rho_e\},
\end{equation}
and put $g:=\abs{E\setminus G}$. Then $0\le u_G\le u_E$ on $G$, the
difference $h:=u_E-u_G$ \textup(extended by $u_E$ on $E\setminus G$, equivalently
$h=u_E-\bar u_G$ where $\bar u_G\in H^1_0(E)$ is the extension of $u_G$ by
zero\textup) lies in $H^1_0(E)$ and is harmonic in $G$, and
\begin{equation}\label{eq:torsion-loss}
   0\ \le\ T(E)-T(G)
   \ =\ \int_{E\setminus G}u_E\ +\ \int_G h
   \ \le\ \tilde\eta\,g\ +\ \int_G h .
\end{equation}
Moreover the boundary-layer term $\int_G h$ is non-negative, satisfies the
energy identity
\begin{equation}\label{eq:energy-id}
   \int_G h\ \le\ T(E)-T(G)\ =\ \int_E\abs{\nabla h}^2,
\end{equation}
and obeys $0\le h\le\tilde\eta$ throughout $G$, whence the crude bound
$\int_G h\le\tilde\eta\,\abs G$.
\end{lemma}

\begin{proof}
\emph{Step A: the exact decomposition.}
Since $G\subset E$, the restriction $u_E|_G$ solves $-\Delta u_E=1=-\Delta
u_G$ in $G$ with $u_E\ge0=u_G$ on $\partial G$; by the comparison principle
on $G$ the difference $h:=u_E-u_G$ is then $\ge0$, i.e.\ $0\le u_G\le u_E$ on
$G$. Being a difference of two solutions of $-\Delta(\cdot)=1$, $h$ is
harmonic in $G$, with boundary trace $h=u_E$ on $\partial G$ (as $u_G=0$
there). Splitting $E=G\cup(E\setminus G)$ and using $u_G=0$ on $E\setminus G$,
\begin{equation}\label{eq:TEG-split}
   T(E)-T(G)=\int_E u_E-\int_G u_G
   =\int_{E\setminus G}u_E+\int_G(u_E-u_G)
   =\int_{E\setminus G}u_E+\int_G h,
\end{equation}
which is the equality in \eqref{eq:torsion-loss}; both summands are
$\ge0$, so $T(E)-T(G)\ge0$.

\emph{Step B: the cap term.}
Since $E\subset B_{\rho_e}(z)$, the comparison principle gives $u_E\le
u_{B_{\rho_e}(z)}$ on $E$: the difference $u_{B_{\rho_e}(z)}-u_E$ is harmonic
in $E$ (both solve $-\Delta(\cdot)=1$) and equals $u_{B_{\rho_e}(z)}\ge0$ on
$\partial E$, hence is $\ge0$ in $E$ by the minimum principle. The disc
torsion function is explicit,
\begin{equation}\label{eq:discu}
   u_{B_{\rho_e}(z)}(y)=\frac{\rho_e^2-\abs{y-z}^2}{4}.
\end{equation}
On the caps $E\setminus G\subset\{\rho_i<\abs{x-z}<\rho_e\}$ one has
$\abs{x-z}\ge\rho_i$, so
\begin{equation}\label{eq:barrier}
   u_E\ \le\ \frac{\rho_e^2-\abs{x-z}^2}{4}
   \ \le\ \frac{\rho_e^2-\rho_i^2}{4}\ =\ \tilde\eta
   \qquad\text{on } E\setminus G,
\end{equation}
and therefore $\int_{E\setminus G}u_E\le\tilde\eta\,g$. This is the main
term and is fully rigorous; combined with \eqref{eq:TEG-split} it gives the
inequality in \eqref{eq:torsion-loss}.

\emph{Step C: the boundary-layer term, energy form and crude bound.}
Let $\bar u_G\in H^1_0(E)$ denote the extension of $u_G$ by zero outside $G$;
then $h=u_E-\bar u_G\in H^1_0(E)$ agrees with the previously defined $h$ on
$G$ and equals $u_E$ on $E\setminus G$. Since $u_E$ minimises $J_E$ on
$H^1_0(E)$ and $\bar u_G$ minimises $J_E$ restricted to the subspace
$H^1_0(G)\subset H^1_0(E)$ (its minimum there being $J_E(\bar
u_G)=J_G(u_G)=-\tfrac12T(G)$), the parallelogram identity for the quadratic
functional $J_E$ gives the exact energy balance
\begin{equation}\label{eq:energy-balance}
   \tfrac12\int_E\abs{\nabla h}^2
   =J_E(\bar u_G)-J_E(u_E)
   =\bigl(-\tfrac12T(G)\bigr)-\bigl(-\tfrac12T(E)\bigr)
   =\tfrac12\bigl(T(E)-T(G)\bigr),
\end{equation}
using \eqref{eq:torsion-energy} for both $D=E$ and $D=G$. Hence
$T(E)-T(G)=\int_E\abs{\nabla h}^2$, the equality in \eqref{eq:energy-id};
combined with \eqref{eq:torsion-loss} and $\int_{E\setminus G}u_E\ge0$ this
gives $\int_G h\le T(E)-T(G)=\int_E\abs{\nabla h}^2$.

Finally, the bounds $0\le h\le\tilde\eta$ in $G$: non-negativity was shown in
Step A. For the upper bound, $h$ is harmonic in $G$ with boundary trace
$h=u_E$ on $\partial G$, and on $\partial G$ one has $u_E\le\tilde\eta$
--- indeed $\partial G\subset\overline{B_{\rho_e}(z)}$, so the barrier
\eqref{eq:discu} gives $u_E\le u_{B_{\rho_e}(z)}\le\tilde\eta$ on the part of
$\partial G$ lying in $\{\abs{x-z}\ge\rho_i\}$ (which contains
$\partial G\cap E$, the interior cap-floor boundary, since $E\setminus G$ is
confined to the collar by \eqref{eq:collar-hyp}), while $u_E=0$ on
$\partial G\cap\partial E$. By the maximum principle $h\le\tilde\eta$
throughout $G$, whence $\int_G h\le\tilde\eta\abs G$.
\end{proof}

The cap term $\tilde\eta\,g$ in \eqref{eq:torsion-loss} is exactly what the
deficit algebra of \cref{lem:cutdeficit} can afford. The boundary-layer term
$\int_G h$ is genuinely smaller in the application but requires care; we
isolate it next.

\begin{remark}[The boundary-layer term and the residual analytic point]
\label{gap:torsion-stab}
The cap term $\int_{E\setminus G}u_E\le\tilde\eta\,g$ of Step B is rigorous
and, as we show in \cref{lem:cutdeficit}, by itself drives the deficit down
once the harmonic boundary-layer term $\int_G h$ is controlled at the same
order. The sharp description of $\int_G h$ is, by Green's identity in $G$
\textup(legitimate when $\partial G$ admits an outward normal $\HH^1$-a.e.;
for the rough star-shaped $G$ below it is the limiting form along a smooth
exhaustion\textup),
\begin{equation}\label{eq:flux-rep}
   \int_G h=\int_G h\,(-\Delta u_G)
   =-\int_{\partial G}h\,\partial_\nu u_G\,d\HH^1
   =\int_{\partial G}u_E\,\abs{\partial_\nu u_G}\,d\HH^1
   =\int_\Gamma u_E\,\abs{\partial_\nu u_G}\,d\HH^1,
\end{equation}
where $\partial_\nu u_G\le0$ on $\partial G$ \textup(as $u_G\ge0=u_G$ on
$\partial G$\textup) and the integral is supported on the interior cap-floor
boundary $\Gamma:=\partial G\cap E$, because $u_E=0$ on the remainder
$\partial G\cap\partial E$. Since $u_E\le\tilde\eta$ on $\Gamma$ by
\eqref{eq:barrier}, this yields the clean bound
\begin{equation}\label{eq:flux-bound}
   \int_G h\ \le\ \tilde\eta\,\Phi,
   \qquad
   \Phi:=\int_\Gamma\abs{\partial_\nu u_G}\,d\HH^1,
\end{equation}
where $\Phi$ is the flux of $u_G$ through the interior cap-floor. The total
flux through the \emph{whole} boundary is exactly $\int_{\partial
G}\abs{\partial_\nu u_G}=\int_G(-\Delta u_G)=\abs G$, so $\Phi\le\abs G$
unconditionally; the substance is that $\Gamma$ is only the small,
fold-supported part of $\partial G$. For a boundary with a one-sided gradient
bound $\norm{\nabla u_G}_{L^\infty(\partial G)}\le C$ \textup(an $O(1)$
quantity for the disc-trapped torsion function, by the barrier $u_G\le
u_{B_{\rho_e}(z)}$ and interior estimates\textup) one has $\Phi\le
C\,\HH^1(\Gamma)$, so it would suffice to bound the length $\HH^1(\Gamma)$ of
the interior cap-floor by the fold content of \cref{prop:fold}. This is the
\emph{single} residual analytic point: the merely measurable core profile
$\rho_0=r_1$ of \cref{cor:gd} carries no slope control, so $\HH^1(\Gamma)$
--- which involves the radial derivative of $\rho_0$ --- is not controlled by
the fold content (a rough graph can have large perimeter over a small angular
sector). We therefore record the boundary-layer control as the explicit
scoped hypothesis
\begin{equation}\label{eq:flux-hyp}
   \int_G h\ \le\ \tilde\eta\,\Phi\ \le\ \tilde\eta\,g
\end{equation}
\textup(equivalently $\Phi\le g$: the flux through the interior cap-floor is
dominated by the excised area\textup), which holds whenever $\partial G$ has
the regularity making \eqref{eq:flux-rep} valid and $\HH^1(\Gamma)\le
C(\Exc+a\eta)$; under \eqref{eq:flux-hyp} the cutting estimate below is
unconditional. We do not claim \eqref{eq:flux-hyp} for an arbitrary
measurable star-shaped $G$, and we flag every downstream use of it
explicitly.
\end{remark}

\begin{remark}[The barrier constant]\label{rem:barrier}
The width of the collar enters only through
$\tilde\eta=\tfrac14(\rho_e^2-\rho_i^2)$, the maximal value of the disc
torsion function on the inner sphere. With $\rho_i=\hat r-O(\eta)$ and
$\rho_e=\hat r+O(\eta)$ one has $\tilde\eta=\tfrac14(\rho_e+\rho_i)
(\rho_e-\rho_i)\le\tfrac12\rho_e\,\eta\le C\eta$, so $\tilde\eta$ is of the
same order as the annulus width $\eta=\rho_e-\rho_i$. In the unit-scale
normalisation $\rho_e=1+\eta$, $\rho_i=1-\eta$ one gets exactly
$\tilde\eta=\eta$.
\end{remark}

\subsection{Cutting boundary material decreases \texorpdfstring{$\dstar$}{the scale-invariant deficit}}

Recall the scale-invariant relative deficit of a finite-measure set $D$,
\begin{equation}\label{eq:dstar-recall}
   \dstar(D)=\frac{T(B_D)-T(D)}{T(B_D)},\qquad
   T(B_D)=\frac{\abs D^2}{8\pi},
\end{equation}
where $B_D$ is any ball with $\abs{B_D}=\abs D$ (so
$T(B_D)=|B_D|^2/(8\pi)$ by scaling from $T(B_1)=\pi/8$). Equivalently,
$\dstar(D)=1-8\pi\,T(D)/\abs D^2$. We now show that excising radial caps
confined to a thin collar lowers $\dstar$, provided the torsion loss per unit
excised area stays below the torsion density of the larger set. The
hypothesis is stated directly as this \emph{loss-slope} condition, which is
exactly what the sign-critical algebra consumes and which the geometric input
of \cref{lem:torsion-stab} supplies.

\begin{lemma}[Cutting boundary material decreases $\dstar$]\label{lem:cutdeficit}
Let $z\in\R^2$, $0<\rho_i\le\rho_e$ and $\tilde\eta=\tfrac14(\rho_e^2-
\rho_i^2)$. Let $E,G$ be bounded open sets with
\begin{equation}\label{eq:cut-incl}
   B_{\rho_i}(z)\subset G\subset E\subset B_{\rho_e}(z),
   \qquad E\setminus G\subset\{\rho_i<\abs{x-z}<\rho_e\},
\end{equation}
put $V:=\abs E$, $g:=\abs{E\setminus G}=V-\abs G$, and suppose the
\emph{torsion-loss bound}
\begin{equation}\label{eq:cutcond}
   T(E)-T(G)\ \le\ \lambda\,g
   \qquad\text{with}\qquad
   \lambda\ \le\ \frac{T(E)}{\abs E}
\end{equation}
holds. Then $\dstar(G)\le\dstar(E)$.
\end{lemma}

\begin{proof}
Write $T:=T(E)$; then $\abs G=V-g$ and, since $\varnothing\ne B_{\rho_i}(z)
\subset G\subset E$, we have $0\le g\le V$ and $V>0$.

\emph{Step 1: lower bound for $T(G)$.} This is the hypothesis
\eqref{eq:cutcond} rearranged:
\begin{equation}\label{eq:TGlower}
   T(G)\ \ge\ T-\lambda\,g,
   \qquad \lambda\ \le\ \frac TV .
\end{equation}

\emph{Step 2: the deficit comparison reduces to a torsion inequality.} By
\eqref{eq:dstar-recall}, with $\abs E=V$ and $\abs G=V-g$,
\begin{equation}\label{eq:dstar-diff}
   \dstar(E)-\dstar(G)
   =\Bigl(1-\frac{8\pi\,T}{V^2}\Bigr)-\Bigl(1-\frac{8\pi\,T(G)}{(V-g)^2}\Bigr)
   =8\pi\Bigl[\frac{T(G)}{(V-g)^2}-\frac{T}{V^2}\Bigr].
\end{equation}
Thus $\dstar(G)\le\dstar(E)$ \emph{holds if and only if}
$\dfrac{T(G)}{(V-g)^2}\ge\dfrac{T}{V^2}$, i.e.\ (multiplying by
$(V-g)^2\ge0$) if and only if
\begin{equation}\label{eq:target}
   T(G)\ \ge\ T\,\frac{(V-g)^2}{V^2}\ =\ T\Bigl(1-\frac gV\Bigr)^2 .
\end{equation}
We now prove \eqref{eq:target}. (When $g=V$, $\abs G=0$ is impossible since
$G\supset B_{\rho_i}(z)$, so in fact $g<V$; but the argument below covers
$g\le V$ verbatim.)

\emph{Step 3: the sign-critical chain.} Expand the right-hand side of
\eqref{eq:target}:
\begin{equation}\label{eq:expand}
   T\Bigl(1-\frac gV\Bigr)^2
   =T-\frac{2Tg}{V}+\frac{Tg^2}{V^2}.
\end{equation}
Subtracting \eqref{eq:expand} from the lower bound \eqref{eq:TGlower} for
$T(G)$,
\begin{align}
   T(G)-T\Bigl(1-\frac gV\Bigr)^2
   &\ \ge\ \bigl(T-\lambda\,g\bigr)
     -\Bigl(T-\frac{2Tg}{V}+\frac{Tg^2}{V^2}\Bigr)
   \notag\\
   &\ =\ \frac{2Tg}{V}-\frac{Tg^2}{V^2}-\lambda\,g
   \ =\ g\Bigl(\frac{2T}{V}-\frac{Tg}{V^2}-\lambda\Bigr).
   \label{eq:bracket}
\end{align}
We bound the bracket from below. Since $0\le g\le V$ we have
$0\le g/V\le1$, hence $2-g/V\ge1$, and therefore
\begin{equation}\label{eq:bracket-key}
   \frac{2T}{V}-\frac{Tg}{V^2}
   =\frac TV\Bigl(2-\frac gV\Bigr)\ \ge\ \frac TV\ \ge\ \lambda,
\end{equation}
the last inequality being exactly the loss-slope hypothesis
\eqref{eq:cutcond} ($\lambda\le T/V=T(E)/\abs E$). Consequently the bracket
in \eqref{eq:bracket} is $\ge\frac TV-\lambda\ge0$, and since $g\ge0$ the
whole expression \eqref{eq:bracket} is $\ge0$. This is \eqref{eq:target},
which by Step 2 yields $\dstar(G)\le\dstar(E)$.
\end{proof}

\begin{remark}[Supplying the loss-slope hypothesis from
\cref{lem:torsion-stab}]\label{rem:supply}
\Cref{lem:torsion-stab} furnishes \eqref{eq:cutcond} as follows. Its
hypotheses coincide with \eqref{eq:cut-incl}, and \eqref{eq:torsion-loss}
gives the exact torsion loss $T(E)-T(G)=\int_{E\setminus G}u_E+\int_G h\le
\tilde\eta\,g+\int_G h$. Under the boundary-layer control \eqref{eq:flux-hyp}
of \cref{gap:torsion-stab} --- $\int_G h\le\tilde\eta\,g$ --- this reads
\begin{equation}\label{eq:loss-2eta}
   T(E)-T(G)\ \le\ 2\,\tilde\eta\,g,
\end{equation}
so \eqref{eq:cutcond} holds with the explicit slope $\lambda=2\tilde\eta$,
provided $2\tilde\eta\le T(E)/\abs E$. Even with only the crude bound
$\int_G h\le\tilde\eta\abs G$ of \cref{lem:torsion-stab} one obtains the loss
$\le\tilde\eta(g+\abs G)=\tilde\eta\abs E$, but that slope is
$\tilde\eta\abs E/g$, which is \emph{not} below $T/V$ at the relevant scales;
the gain from \eqref{eq:flux-hyp} is precisely the localisation of the
boundary-layer term to the cap region, replacing the slope $\tilde\eta\abs
E/g$ by the absolute $2\tilde\eta$.
\end{remark}

\begin{remark}[The loss-slope hypothesis \eqref{eq:cutcond} holds with room
to spare]\label{rem:roomtospare}
In the application $E=\Etil$ is the inset of \cref{def:inset}, of measure
$\mu_E=\abs\Etil=\pi-O(\sqrt\dT)\ge\pi/2$ and deficit
$\dstar(\Etil)\le\tfrac12$ for $\dT\le\delta_\star$. Then the torsion density
is an \emph{absolute} constant bounded below,
\[
   \frac{T(\Etil)}{\abs\Etil}
   =\bigl(1-\dstar(\Etil)\bigr)\frac{\abs\Etil}{8\pi}
   \ \ge\ \bigl(1-\tfrac12\bigr)\cdot\frac{\pi/2}{8\pi}
   \ =\ \frac1{32},
\]
using $T(\Etil)=(1-\dstar(\Etil))\,T(B_{\Etil})=(1-\dstar(\Etil))\,
\abs\Etil^2/(8\pi)$ from \eqref{eq:dstar-recall}. On the other hand the slope
$\lambda=2\tilde\eta$ supplied by \cref{rem:supply} vanishes:
$\tilde\eta\le\tfrac12\rho_e\,\eta\le C\eta\le C(\dT/\tau)^{1/4}\to0$ as
$\dT\to0$ (\cref{thm:annulus}, with $\tau=\sqrt\dT$ giving
$\eta\le C\dT^{1/8}$). Hence $\lambda=2\tilde\eta\le\tfrac1{32}\le
T(\Etil)/\abs\Etil$ for all $\dT\le\delta_\star$, and the loss-slope
condition \eqref{eq:cutcond} is satisfied with a fixed margin. The crucial
structural point is that the comparison is between a \emph{vanishing} slope
$\lambda=2\tilde\eta\to0$ and an \emph{absolute} torsion density $\ge1/32$,
so the inequality $\lambda\le T(\Etil)/\abs\Etil$ holds robustly; the
threshold $1/32$ is not sharp and is recorded only to fix $\delta_\star$.
This is exactly where the boundary-layer control \eqref{eq:flux-hyp} enters:
it is what turns the exact loss $\int_{E\setminus G}u_E+\int_G h$ into a bound
with absolute slope $2\tilde\eta$ rather than the much larger
$\tilde\eta\abs\Etil/g$.
\end{remark}

\begin{remark}[The application: cutting the caps]\label{rem:twouses}
\Cref{lem:cutdeficit} is used in the assembly (\cref{thm:reduction}) with
$E=\Etil$ and $G=G_0=\{z+r\omega:r<\rho_0(\omega)\}$ the measurable radial
core graph of \cref{cor:gd,prop:realization}, taken in its open star-shaped
representative. The inclusions $B_{\rho_i}(z)\subset G_0\subset\Etil\subset
B_{\rho_e}(z)$ hold by \cref{cor:gd,thm:annulus}, the difference
$\Etil\setminus G_0$ is the union of the radial caps, all contained in the
collar $\{\rho_i<\abs{x-z}<\rho_e\}$, and the exact torsion-loss decomposition
$T(\Etil)-T(G_0)=\int_{\Etil\setminus G_0}u_{\Etil}+\int_{G_0}h\le\tilde\eta\,g
+\int_{G_0}h$ of \cref{lem:torsion-stab} holds unconditionally. Under the
boundary-layer control \eqref{eq:flux-hyp} of \cref{gap:torsion-stab} the loss
is $\le2\tilde\eta\,g$ (\cref{rem:supply}); together with the absolute
torsion-density lower bound of \cref{rem:roomtospare}, the loss-slope
condition \eqref{eq:cutcond} holds with $\lambda=2\tilde\eta$, and
\cref{lem:cutdeficit} gives that cutting the caps lowers the deficit,
\[
   \dstar(G_0)\ \le\ \dstar(\Etil).
\]
This is the one place where the geometry of \cref{lem:cutdeficit} is needed:
the discarded material is precisely low-torsion boundary material in the
$\eta$-collar about the \emph{single} centre $z$. The deduction is
unconditional except for the single residual analytic point \eqref{eq:flux-hyp}
isolated in \cref{gap:torsion-stab} --- the localisation of the harmonic
boundary-layer term to the cap region; the cap term $\tilde\eta\,g$ and the
entire deficit algebra are rigorous.
\end{remark}

\begin{remark}[The other deficit step is \emph{not} an instance of
\cref{lem:cutdeficit}]\label{rem:transfer-not-cut}
The earlier reduction step --- passing from the full superlevel set
$\{u>\hat v\}$ to its largest connected component $\Etil$, i.e.\ discarding
the spillover components of total measure
$\abs{\{u>\hat v\}\setminus\Etil}\le D_I(\hat v)^2/(16\pi^3)$
(\cref{lem:spill}) --- is \emph{not} covered by \cref{lem:cutdeficit}, and we
do not claim it is. \Cref{lem:cutdeficit} requires
$G\subset E\subset B_{\rho_e}(z)$ with $E\setminus G$ inside one collar about
$z$; but the discarded spillover components are distinct level components,
which need not lie in the trapping ball $B_{\rho_e}(z)$ of $\Etil$ (they can
sit elsewhere in $\Om$), so the hypothesis $E\subset B_{\rho_e}(z)$ may fail
for $E=\{u>\hat v\}$. The correct tool for that step is the deficit transfer
\cref{lem:transfer} of \cref{sec:inset} combined with the \emph{exact}
additivity of torsion over connected components (the function $u-\hat v$ is
the common torsion function of every component of $\{u>\hat v\}$); this is
carried out in \cref{prop:insetcontrols}, yielding $\dstar(\Etil)\le
(\pi^2/\mu_E^2)\,\dstar(\Om)+O(\dT^2/\tau^2)$ with no appeal to the collar
geometry. Thus the full chain
\[
   \dstar(G_0)\ \le\ \dstar(\Etil)
   \ \le\ \frac{\pi^2}{\mu_E^2}\,\dstar(\Om)+O\!\Bigl(\frac{\dT^2}{\tau^2}\Bigr)
\]
combines \cref{lem:cutdeficit} (first inequality, this section) with
\cref{lem:transfer}/\cref{prop:insetcontrols} (second inequality,
\cref{sec:inset}); the two play distinct roles and rest on distinct
hypotheses.
\end{remark}

\begin{remark}[Why the deficit decreases — the mechanism]\label{rem:mechanism}
The single inequality \eqref{eq:bracket-key} is the entire content. Cutting
a cap of measure $g$ costs at most $\lambda\,g$ in torsion (\emph{linear} in
$g$, with a small slope $\lambda=2\tilde\eta\to0$ under \eqref{eq:flux-hyp}),
but it shrinks the equal-area comparison ball, lowering its torsion
$T(B_G)=\abs G^2/(8\pi)$ by $\approx\frac{2T(B_E)}{V}g=\frac{V}{4\pi}g$, i.e.\
by a slope $\approx\frac{2T}{V}$ after accounting for the deficit factor.
Because the unperturbed torsion density $T/V$ of a near-disc is an
\emph{absolute} constant of order $1$ while the collar slope
$\lambda=2\tilde\eta$ vanishes, the denominator $T(B_G)$ drops faster than the
numerator $T(G)$, and $\dstar=1-T/T(B)$ decreases. There is no square-root
loss and no constant that degrades as $\dT\to0$: the deficit algebra is exact
and the margin in \eqref{eq:cutcond} only widens. The exactness of the
torsion-loss \emph{decomposition} \eqref{eq:torsion-loss} is likewise
unconditional; the only inexact step is the localisation \eqref{eq:flux-hyp}
of the harmonic boundary-layer term, isolated in \cref{gap:torsion-stab}.
\end{remark}

\section{Perturbative stability of nearly spherical sets}
\label{sec:nearly}

This section proves, in full and self-containedly, the perturbative
torsional stability of small-amplitude radial graphs consumed by the
reduction of \cref{sec:main}. The result is a Fuglede-type estimate
\emph{for the torsion functional}: a planar set whose boundary is a graph of
small amplitude over a circle, after the natural volume and barycentre
normalisations, has scale-invariant torsional deficit controlling the square
of its Fraenkel asymmetry. We invoke no selection principle, no
spectral black box, and no quantitative isoperimetric inequality. The proof
is an explicit second-order torsion computation through the dual
(complementary) variational principle, with an admissible vector field built
by hand; its only nonclassical ingredient is a frequency-dependent damping of
the harmonic corrector that controls high-frequency boundary oscillations and
lets the hypothesis be merely $\|\varphi\|_{L^\infty}\le\eta_G$.

\subsection{Statement and standing normalisation}

Throughout, $S^1=\R/2\pi\Z$, $\omega(\theta)=(\cos\theta,\sin\theta)$, and for
$\varphi\in \mathrm{Lip}(S^1)$ with $\|\varphi\|_{L^\infty}<1$ and $z\in\R^2$
we write
\begin{equation}\label{eq:Gdef}
   G=G(z,\hat r_G,\varphi):=\set{z+r\omega(\theta):\ \theta\in S^1,\
   0\le r<\hat r_G\,(1+\varphi(\theta))}
\end{equation}
for the radial graph of amplitude profile $\varphi$, base radius
$\hat r_G=\sqrt{|G|/\pi}$ and centre $z$. We recall that, for a bounded set
$E$ of finite measure with torsion $T(E)=\int_E u_E$ ($-\Delta u_E=1$,
$u_E\in H^1_0(E)$),
\begin{equation}\label{eq:dstarAdef}
   \dstar(E)=\frac{T(B_E)-T(E)}{T(B_E)},\quad T(B_E)=\frac{|E|^2}{8\pi},
   \qquad
   \AF(E)=\inf_{x_0}\frac{|E\triangle B(x_0,r_E)|}{|E|},\ r_E=\sqrt{|E|/\pi},
\end{equation}
are the scale-invariant deficit and the Fraenkel asymmetry; both are
invariant under translations and dilations.

\begin{theorem}[Radial-graph torsion stability]\label{thm:graphG}
There are absolute constants
\begin{equation}\label{eq:graphG-const}
   \eta_G=\tfrac1{150},\qquad C_G=48,
\end{equation}
such that the following holds. Let $G=G(z,\hat r_G,\varphi)$ be a radial graph
as in \eqref{eq:Gdef} with $\varphi\in\mathrm{Lip}(S^1)$,
$\|\varphi\|_{L^\infty}\le\eta_G$, and barycentre at $z$. Then
\begin{equation}\label{eq:graphG}
   \AF(G)^2\ \le\ C_G\,\dstar(G).
\end{equation}
The statement is scale-invariant; $G$ need not have area $\pi$. It uses no
property of the original set $\Om$ and no selection.
\end{theorem}

Since $\dstar$ and $\AF$ are dilation- and translation-invariant, replacing
$G$ by $\hat r_G^{-1}(G-z)$ we assume for the proof
\begin{equation}\label{eq:norm}
   z=0,\qquad \hat r_G=1,\qquad\text{so}\qquad |G|=\pi,
\end{equation}
which fixes the volume constraint
\begin{equation}\label{eq:vol-constraint}
   |G|=\int_0^{2\pi}\frac{(1+\varphi)^2}{2}\d\theta=\pi
   \quad\Longleftrightarrow\quad
   \int_0^{2\pi}\bigl(2\varphi+\varphi^2\bigr)\d\theta=0 .
\end{equation}
Under \eqref{eq:norm}, $T(B_G)=\pi/8$, while $u_{B_1}=(1-|x|^2)/4$ and
$T(B_1)=\pi/8$. Expand $\varphi$ in its real Fourier series,
\begin{equation}\label{eq:fourier}
   \varphi(\theta)=a_0+\sum_{k\ge1}\bigl(a_k\cos k\theta+b_k\sin k\theta\bigr),
   \qquad
   \|\varphi\|_{L^2(S^1)}^2=\pi\Bigl(2a_0^2+\sum_{k\ge1}(a_k^2+b_k^2)\Bigr) .
\end{equation}
Lipschitz regularity gives $\sum_k k^2(a_k^2+b_k^2)<\infty$, so the series
below converge in the indicated norms.

\subsection{The dual (complementary) variational principle}

The torsion has two dual characterisations: the Dirichlet principle
$T(E)=\max_{w\in H^1_0(E)}\bigl(2\int_E w-\int_E|\nabla w|^2\bigr)$ (attained
at $w=u_E$), and the complementary principle below, which produces
\emph{upper} bounds on $T(E)$ — equivalently \emph{lower} bounds on the
deficit — from explicit divergence-constrained fields.

\begin{lemma}[Complementary principle]\label{lem:dual}
Let $E\subset\R^2$ be bounded open with finite perimeter and torsion function
$u_E\in H^1_0(E)$. For every $\sigma\in L^2(E;\R^2)$ with $\divg\sigma=-1$ in
$\mathcal D'(E)$,
\begin{equation}\label{eq:dual}
   T(E)\ \le\ \int_E|\sigma|^2,
\end{equation}
with equality iff $\sigma=\nabla u_E$.
\end{lemma}

\begin{proof}
Since $u_E\in H^1_0(E)$ and $\divg\sigma=-1$, the pairing
$\int_E\sigma\cdot\nabla u_E=-\langle\divg\sigma,u_E\rangle=\int_E u_E$ is
well defined, and $\int_E|\nabla u_E|^2=\int_E u_E=T(E)$ by testing
$-\Delta u_E=1$ with $u_E$. Hence
\[
   0\le\int_E|\sigma-\nabla u_E|^2
   =\int_E|\sigma|^2-2\!\int_E\sigma\cdot\nabla u_E+\!\int_E|\nabla u_E|^2
   =\int_E|\sigma|^2-T(E),
\]
which is \eqref{eq:dual}; equality forces $\sigma=\nabla u_E$ a.e.
\end{proof}

The disc field $\sigma_0:=\nabla u_{B_1}=-x/2$ satisfies $\divg\sigma_0=-1$
and $\int_{B_1}|\sigma_0|^2=T(B_1)=\pi/8$. We perturb it into an admissible
field on $G$.

\subsection{The comparison field}

Set the corrector coefficients
\begin{equation}\label{eq:gk}
   \alpha_k:=\tfrac12 m_k\,a_k,\quad \beta_k:=\tfrac12 m_k\,b_k,
   \qquad m_k:=\min\Bigl(1,\tfrac2k\Bigr)\in(0,1],\quad k\ge1 .
\end{equation}
The damping $m_k$ equals $1$ at the critical mode $k=2$ and decays like $2/k$
for $k>2$, so the corrector \emph{slope} stays bounded:
\begin{equation}\label{eq:slopebound}
   k|\alpha_k|=\tfrac12 m_k k\,|a_k|\le|a_k|,
   \qquad k|\beta_k|\le|b_k| .
\end{equation}
Define the harmonic corrector and its normal trace on $\partial B_1$,
\begin{equation}\label{eq:hC}
   h(r,\theta):=\sum_{k\ge1}r^k\bigl(\alpha_k\cos k\theta+\beta_k\sin k\theta\bigr),
   \qquad
   c(\theta):=\partial_r h(1,\theta)=\sum_{k\ge1}k\bigl(\alpha_k\cos k\theta+\beta_k\sin k\theta\bigr).
\end{equation}
By \eqref{eq:slopebound} and Lipschitz regularity, $\nabla h$ and $c$ are in
$L^2$, $\Delta h=0$ in $\R^2$, and
\begin{equation}\label{eq:cL2}
   \|c\|_{L^2(S^1)}^2=\pi\sum_{k\ge1}k^2(\alpha_k^2+\beta_k^2)
   \le\pi\sum_{k\ge1}(a_k^2+b_k^2)\le\|\varphi\|_{L^2}^2 .
\end{equation}
The comparison field on $G$ is
\begin{equation}\label{eq:sigma}
   \sigma(x):=
   \begin{cases}
      \nabla\bigl(u_{B_1}+h\bigr)(x), & x\in G,\ |x|\le1,\\[1mm]
      \Bigl(-\dfrac r2+\dfrac{c(\theta)}{r}\Bigr)\,\omega(\theta),
        & x=r\omega(\theta)\in G,\ |x|>1 .
   \end{cases}
\end{equation}
The collar line is used only where $\varphi(\theta)>0$, i.e.\ where $G$
protrudes beyond $B_1$.

\begin{lemma}[Admissibility]\label{lem:admissible}
The field $\sigma$ of \eqref{eq:sigma} lies in $L^2(G;\R^2)$ and satisfies
$\divg\sigma=-1$ in $\mathcal D'(G)$.
\end{lemma}

\begin{proof}
On $\{|x|<1\}$, $\divg\sigma=\Delta(u_{B_1}+h)=-1+0=-1$. On the collar
$\{|x|>1\}$, for a radial field $f(r)\omega$ one has
$\divg(f\omega)=\frac1r\frac{d}{dr}(rf)$; with $f=-\frac r2+\frac{c}{r}$ this
is $\frac1r\frac{d}{dr}\bigl(-\frac{r^2}2+c\bigr)=-1$. Across the interface
$\Gamma=\partial B_1\cap G$, the normal (radial) components match: from inside
$\sigma\cdot\omega=\partial_r(u_{B_1}+h)|_{r=1}=-\tfrac12+c(\theta)$, from the
collar $\sigma\cdot\omega=-\tfrac12+c(\theta)$. Hence the distributional
divergence carries no surface measure on $\Gamma$ and equals $-1$ on all of
$G$. Square-integrability: $\nabla u_{B_1}$ is smooth, $\nabla h\in
L^2(B_{1+\eta_G})$ by \eqref{eq:slopebound}, and the collar is a thin set of
width $\le\eta_G$ on which $|\sigma|^2\le2(\tfrac r2)^2+2c(\theta)^2$.
\end{proof}

\subsection{The exact second-order energy}

We compute $\int_G|\sigma|^2$ exactly and separate the quadratic form from a
genuinely cubic remainder. Write $p:=\varphi(\theta)$ and
\begin{equation}\label{eq:Phidef}
   \Phi_{\mathrm{col}}(r,\theta):=\Bigl(-\tfrac r2+\tfrac{c}{r}\Bigr)^2 r,
   \qquad
   \Phi_{\mathrm{in}}(r,\theta):=\Bigl[\bigl(-\tfrac r2+\partial_r h\bigr)^2
   +\bigl(\tfrac1r\partial_\theta h\bigr)^2\Bigr]r ,
\end{equation}
the polar energy densities of the collar and interior fields.

\begin{lemma}[Energy identity]\label{lem:energy}
With $\sigma$ as in \eqref{eq:sigma},
\begin{equation}\label{eq:energy-split}
   \int_G|\sigma|^2
   =\frac\pi8
     +\pi\sum_{k\ge1}k(\alpha_k^2+\beta_k^2)
     +\mathcal N,
   \qquad
   \mathcal N=\int_{\{p>0\}}\!\!\int_1^{1+p}\!\!\Phi_{\mathrm{col}}\,dr\,d\theta
     -\int_{\{p<0\}}\!\!\int_{1+p}^1\!\!\Phi_{\mathrm{in}}\,dr\,d\theta .
\end{equation}
\end{lemma}

\begin{proof}
Split $\int_G=\int_{B_1}-\int_{B_1\setminus G}+\int_{G\setminus B_1}$. Over
the full disc, $\int_{B_1}|\nabla(u_{B_1}+h)|^2
=\int_{B_1}|\nabla u_{B_1}|^2+2\int_{B_1}\nabla u_{B_1}\cdot\nabla
h+\int_{B_1}|\nabla h|^2$. The first term is $\pi/8$; the cross term vanishes,
$\int_{B_1}\nabla u_{B_1}\cdot\nabla h=-\int_{B_1}u_{B_1}\Delta
h+\int_{\partial B_1}u_{B_1}\partial_\nu h=0$ (as $\Delta h=0$,
$u_{B_1}|_{\partial B_1}=0$); and, by orthogonality of the harmonic modes,
$\int_{B_1}|\nabla h|^2=\pi\sum_k k(\alpha_k^2+\beta_k^2)$ since
$\int_{B_1}|\nabla(r^k\cos k\theta)|^2=\pi k$. The dimple
$-\int_{B_1\setminus G}|\nabla(u_{B_1}+h)|^2$ (where $p<0$) and the collar
$+\int_{G\setminus B_1}|\sigma|^2$ (where $p>0$) are precisely $\mathcal N$.
\end{proof}

\begin{lemma}[Quadratic form and cubic remainder]\label{lem:BL}
Assume $\|\varphi\|_{L^\infty}\le\tfrac1{150}$ and the corrector
\eqref{eq:gk}. Then
\begin{equation}\label{eq:quadform}
   \int_G|\sigma|^2
   =\frac\pi8
     +\frac14\!\int_{S^1}\!\varphi
     +\frac38\!\int_{S^1}\!\varphi^2
     +\pi\!\sum_{k\ge1}\!\bigl(k\alpha_k^2-k\alpha_k a_k\bigr)
     +\pi\!\sum_{k\ge1}\!\bigl(k\beta_k^2-k\beta_k b_k\bigr)
     +\mathcal R,
\end{equation}
where the remainder obeys the explicit bound
\begin{equation}\label{eq:Rbound}
   |\mathcal R|\ \le\ 5\,\|\varphi\|_{L^\infty}\bigl(\|\varphi\|_{L^2}^2+\|c\|_{L^2}^2\bigr)
   \ \le\ 10\,\|\varphi\|_{L^\infty}\,\|\varphi\|_{L^2}^2 ,
\end{equation}
the last step by \eqref{eq:cL2}.
\end{lemma}

\begin{proof}
We expand $\mathcal N$ from \eqref{eq:energy-split} to second order in $p$ and
collect the cubic tail into $\mathcal R$. The two boundary densities are
$C^\infty$ in $r$ on $[\tfrac12,2]\supset[1-\eta_G,1+\eta_G]$, with values and
$r$-derivatives at $r=1$
\begin{equation}\label{eq:traces}
   \Phi_{\mathrm{col}}(1)=\bigl(\tfrac12-c\bigr)^2,\quad
   \Phi_{\mathrm{in}}(1)=\bigl(\tfrac12-c\bigr)^2+s^2,\quad
   \Phi_{\mathrm{col}}'(1)=\tfrac34-c-c^2,
\end{equation}
where $s(\theta):=\tfrac1r\partial_\theta h(1,\theta)$,
$\|s\|_{L^2}=\|c\|_{L^2}$ (the tangential and normal traces of a harmonic
function on $\partial B_1$ have equal $L^2$ norm). For each $\theta$, Taylor's
theorem with integral remainder gives, on the collar branch $p>0$,
\[
   \int_1^{1+p}\!\Phi_{\mathrm{col}}\,dr
   =p\,\Phi_{\mathrm{col}}(1)+\tfrac12 p^2\,\Phi_{\mathrm{col}}'(1)
   +\int_1^{1+p}\tfrac12(1+p-r)^2\,\Phi_{\mathrm{col}}''(r)\,dr,
\]
and analogously on the dimple branch $p<0$ for $\Phi_{\mathrm{in}}$. We keep
the \emph{disc} part of the first two terms in the quadratic form and put the
$c,s$-dependent part and the third-order tail into $\mathcal R$. Concretely,
writing $\Phi_{\mathrm{col}}(1)=\tfrac14-c+c^2$ and
$\tfrac12\Phi_{\mathrm{col}}'(1)=\tfrac38-\tfrac c2-\tfrac{c^2}2$, the kept
quadratic contribution of the collar branch is
$\int_{\{p>0\}}\bigl(\tfrac14 p-pc+\tfrac38 p^2\bigr)$, and the residue
$\int_{\{p>0\}}\bigl(p\,c^2-\tfrac12 p^2(c+c^2)\bigr)+(\text{tail})$ goes to
$\mathcal R$; the dimple branch is identical except for the extra
$\int_{\{p<0\}}p\,s^2$ (from $\Phi_{\mathrm{in}}(1)$), also placed in
$\mathcal R$. Since $\Phi_{\mathrm{col}}(1)$ and $\Phi_{\mathrm{in}}(1)$ share
the kept part $\tfrac14-c$, the kept quadratic over the full circle is
\[
   \int_{S^1}\Bigl(\tfrac14\varphi-\varphi\,c+\tfrac38\varphi^2\Bigr)
   =\tfrac14\!\int\!\varphi+\tfrac38\!\int\!\varphi^2
     -\pi\!\sum_k k(\alpha_k a_k+\beta_k b_k),
\]
the last equality by \eqref{eq:hC} and orthogonality. Adding the harmonic
term $\pi\sum_k k(\alpha_k^2+\beta_k^2)$ from \eqref{eq:energy-split} gives
exactly the displayed quadratic form in \eqref{eq:quadform}.

\emph{Remainder bound.} Each residue is controlled pointwise by $|p|\le\eta_G$
and the traces $c,s$. The second-order $c$-residues satisfy
\[
   \Bigl|\int_{S^1}\!\bigl(\varphi\,c^2-\tfrac12\varphi^2(c+c^2)\bigr)\Bigr|
   +\Bigl|\int_{\{p<0\}}\!\varphi\,s^2\Bigr|
   \le \eta_G\!\int\!\bigl(c^2+s^2\bigr)+\tfrac12\eta_G^2\!\int\!\bigl(|c|+c^2\bigr)
   \le 2\eta_G\,\|c\|_{L^2}^2,
\]
using $\|s\|_{L^2}=\|c\|_{L^2}$, $\eta_G\le\tfrac1{150}$, and
$\int|c|\le\sqrt{2\pi}\|c\|_{L^2}\le\|c\|_{L^2}^2/\eta_G$ when
$\|c\|_{L^2}\ge\sqrt{2\pi}\,\eta_G$ (otherwise the term is itself
$\le2\eta_G\|c\|_{L^2}^2$ trivially). For the Taylor tails, on the relevant
$r$-interval $[1-\eta_G,1+\eta_G]$ one has the bounds
$|\Phi_{\mathrm{col}}''|\le 2(1+c^2)$ and $|\Phi_{\mathrm{in}}''|\le
2(1+c^2+s^2)$ (direct differentiation of \eqref{eq:Phidef}, using $r\le
1+\eta_G\le\tfrac{151}{150}$ and, on the dimple, $r\le1$ so $r^{k-1}\le1$,
whence the corrector contributes only through $c,s$ up to an absolute
factor), giving
\[
   |\text{tail}|\le\tfrac16\!\int|p|^3\bigl(1+c^2+s^2\bigr)
   \le\tfrac16\eta_G\!\int\varphi^2+\tfrac16\eta_G\!\int\varphi^2(c^2+s^2)
   \le\tfrac16\eta_G\,\|\varphi\|_{L^2}^2+\tfrac13\eta_G\,\|c\|_{L^2}^2 .
\]
Collecting, $|\mathcal R|\le\eta_G\bigl(\tfrac16\|\varphi\|_{L^2}^2+(2+\tfrac13)\|c\|_{L^2}^2\bigr)
\le 5\,\eta_G(\|\varphi\|_{L^2}^2+\|c\|_{L^2}^2)$, which is \eqref{eq:Rbound}.
\end{proof}

\begin{remark}[On the damping $m_k$]\label{rem:damping}
The undamped corrector $\alpha_k=\tfrac12 a_k$ minimises the quadratic form
mode by mode and gives the sharp coefficient $\tfrac\pi8(3-2k)$, but its
normal trace has $\|c\|_{L^2}^2=\tfrac\pi4\sum_k k^2(a_k^2+b_k^2)$, an
$H^1$-quantity \emph{not} controlled by $\|\varphi\|_{L^2}$; then the
remainder would carry an unbounded factor of $k$. The damping
\eqref{eq:gk} caps the slope, $\|c\|_{L^2}\le\|\varphi\|_{L^2}$ by
\eqref{eq:cL2}, so $\mathcal R$ is genuinely cubic \emph{uniformly in the
frequency $k$}, at the cost of only a bounded factor in the leading
coefficient (preserved exactly at the critical mode $k=2$). This is the sole
place high-frequency oscillations are handled, and it is what allows the
hypothesis to be merely $\|\varphi\|_{L^\infty}\le\eta_G$.
\end{remark}

\subsection{Coercivity of the deficit}

We first read off the per-mode quadratic coefficient.

\begin{lemma}[Per-mode deficit coefficient]\label{lem:permode}
Using the volume identity
$\tfrac14\int_{S^1}\varphi+\tfrac38\int_{S^1}\varphi^2=-\tfrac14\int_{S^1}\varphi^2$
(immediate from \eqref{eq:vol-constraint}, since
$\int\varphi=-\int\varphi^2$), the net coefficient of
$a_k^2$ in $\tfrac\pi8-\int_G|\sigma|^2$ (and likewise of $b_k^2$) is
\begin{equation}\label{eq:Dk}
   D_k:=\pi\Bigl(\tfrac12 k m_k-\tfrac14 k m_k^2-\tfrac14\Bigr),
   \qquad m_k=\min(1,2/k).
\end{equation}
For $k\ge2$ one has
\begin{equation}\label{eq:Dk-floor}
   D_2=\frac\pi4,\qquad D_k=\pi\Bigl(\tfrac34-\tfrac1k\Bigr)\ \ge\ \frac{5\pi}{12}
   \quad(k\ge3),
   \qquad\text{hence } D_k\ge\frac\pi4\ \ \forall k\ge2,
\end{equation}
while $D_1=0$.
\end{lemma}

\begin{proof}
By \eqref{eq:gk}, $k\alpha_k^2=\tfrac14 k m_k^2 a_k^2$ and $k\alpha_k
a_k=\tfrac12 k m_k a_k^2$, so the contribution of the $k$-th cosine mode to
the corrector sum in \eqref{eq:quadform} is $\pi(k\alpha_k^2-k\alpha_k
a_k)=\pi(\tfrac14 k m_k^2-\tfrac12 k m_k)a_k^2$. The disc-collar quadratic
$\tfrac38\int\varphi^2$ and the linear $\tfrac14\int\varphi$ combine, by the
volume identity above, into $-\tfrac14\int\varphi^2$, whose
$k$-th cosine contribution is $-\tfrac14\pi a_k^2$. Hence the coefficient of
$a_k^2$ in $\int_G|\sigma|^2-\tfrac\pi8$ is
$\pi(\tfrac14 k m_k^2-\tfrac12 k m_k+\tfrac14)$, and its negative is $D_k$ as
in \eqref{eq:Dk}. For $k=2$ ($m_2=1$): $D_2=\pi(1-\tfrac12-\tfrac14)=\tfrac\pi4$.
For $k\ge3$ ($m_k=2/k$): $\tfrac12 k m_k=1$, $\tfrac14 k m_k^2=\tfrac1k$, so
$D_k=\pi(1-\tfrac1k-\tfrac14)=\pi(\tfrac34-\tfrac1k)\ge\pi(\tfrac34-\tfrac13)=\tfrac{5\pi}{12}$.
For $k=1$ ($m_1=1$): $D_1=\pi(\tfrac12-\tfrac14-\tfrac14)=0$.
\end{proof}

\begin{lemma}[Normalisation kills modes $0$ and $1$]\label{lem:normmodes}
Under \eqref{eq:vol-constraint} and barycentre $0$, with
$\|\varphi\|_{L^\infty}\le\eta_G\le\tfrac1{150}$ and
$\rho^2:=\sum_{k\ge2}(a_k^2+b_k^2)$,
\begin{equation}\label{eq:lowmodes}
   2a_0^2+a_1^2+b_1^2\ \le\ \tfrac1{10}\,\rho^2,
   \qquad\text{hence}\qquad
   \|\varphi\|_{L^2}^2=\pi\bigl(2a_0^2+a_1^2+b_1^2+\rho^2\bigr)\le\tfrac{11}{10}\pi\,\rho^2 .
\end{equation}
\end{lemma}

\begin{proof}
By \eqref{eq:vol-constraint}, $2\pi a_0=-\int\varphi^2$, so
$|a_0|\le\tfrac1{2\pi}\|\varphi\|_{L^\infty}\!\int|\varphi|\le\tfrac1{2\pi}\eta_G\sqrt{2\pi}\,\|\varphi\|_{L^2}$,
whence $2a_0^2\le\tfrac{\eta_G^2}\pi\|\varphi\|_{L^2}^2$. For the barycentre,
$\int_G x\,dx=\int_0^{2\pi}\tfrac{(1+\varphi)^3}3\,\omega\,d\theta$; its linear
part is $\int\varphi\,\omega=\pi(a_1,b_1)$, and setting the moment to $0$ gives
$\pi^2(a_1^2+b_1^2)\le\bigl(\tfrac13\!\int(3|\varphi|^2+|\varphi|^3)\bigr)^2
\le(2\!\int\varphi^2)^2\le4\eta_G^2(\sqrt{2\pi}\|\varphi\|_{L^2})^2
=8\pi\eta_G^2\|\varphi\|_{L^2}^2$. Thus
$2a_0^2+a_1^2+b_1^2\le\bigl(\tfrac{\eta_G^2}\pi+\tfrac{8\eta_G^2}\pi\bigr)\|\varphi\|_{L^2}^2
=\tfrac{9\eta_G^2}\pi\|\varphi\|_{L^2}^2$. Inserting
$\|\varphi\|_{L^2}^2=\pi(2a_0^2+a_1^2+b_1^2+\rho^2)$ and solving,
$2a_0^2+a_1^2+b_1^2\le\tfrac{9\eta_G^2}{1-9\eta_G^2}\rho^2$. With
$\eta_G\le\tfrac1{150}$ the factor is $\le\tfrac1{10}$, giving
\eqref{eq:lowmodes}.
\end{proof}

\begin{proposition}[Deficit coercivity]\label{prop:coercive}
For $\eta_G=\tfrac1{150}$, any normalised radial graph \eqref{eq:norm} with
$\|\varphi\|_{L^\infty}\le\eta_G$ and barycentre $0$ obeys
\begin{equation}\label{eq:coercive}
   \dstar(G)\ \ge\ \rho^2=\sum_{k\ge2}(a_k^2+b_k^2).
\end{equation}
\end{proposition}

\begin{proof}
Grouping \eqref{eq:quadform} by mode through \cref{lem:permode} (which already
incorporates the volume identity
$\tfrac14\int\varphi+\tfrac38\int\varphi^2=-\tfrac14\int\varphi^2$),
\begin{equation}\label{eq:Tub}
   \frac\pi8-\int_G|\sigma|^2
   =\sum_{k\ge1}D_k(a_k^2+b_k^2)-\mathcal R
   \ \ge\ \frac\pi4\,\rho^2-|\mathcal R|,
\end{equation}
discarding the (vanishing) $k=1$ term and using $D_k\ge\tfrac\pi4$ for $k\ge2$
from \eqref{eq:Dk-floor}. By \eqref{eq:Rbound}, \eqref{eq:cL2} and
\cref{lem:normmodes},
\[
   |\mathcal R|\le 10\,\eta_G\,\|\varphi\|_{L^2}^2
   \le 10\,\eta_G\cdot\tfrac{11}{10}\pi\,\rho^2=11\pi\,\eta_G\,\rho^2
   =\frac{11\pi}{150}\rho^2\le\frac\pi8\rho^2,
\]
since $\tfrac{11}{150}<\tfrac18$. Therefore
$\tfrac\pi8-\int_G|\sigma|^2\ge(\tfrac\pi4-\tfrac\pi8)\rho^2=\tfrac\pi8\rho^2$.
By \cref{lem:dual}, $T(G)\le\int_G|\sigma|^2$, so
$\tfrac\pi8-T(G)\ge\tfrac\pi8\rho^2$; dividing by $T(B_G)=\pi/8$ gives
$\dstar(G)\ge\rho^2\ge\tfrac12\rho^2$.
\end{proof}

\begin{remark}[Quantitative margin]\label{rem:numerics}
The threshold $\eta_G=\tfrac1{150}$ is dictated only by the generous
remainder constant $10$ in \eqref{eq:Rbound}; the true admissible amplitude is
far larger. A direct evaluation of the explicit admissible energy
$\int_G|\sigma|^2$ — itself a rigorous upper bound for $T(G)$ by
\cref{lem:dual}, needing no further estimate — gives the stronger
$\dstar(G)\ge\tfrac{8}\pi\cdot 0.73\,\rho^2\ge\tfrac2\pi\rho^2$ already for
$\|\varphi\|_{L^\infty}\le\tfrac1{13}$, with the infimum of
$\dstar(G)/\rho^2$ attained at the single mode $k=2$ (the ellipse), the
classical Saint--Venant extremal. We adopt the
clean lower bound \eqref{eq:coercive}, $\dstar(G)\ge\rho^2$, which already has
the sharp small-amplitude constant $1$ at $k=2$ doubled by the
volume-normalisation gain; the value of $C_G$ below uses only this bound.
\end{remark}

\subsection{The asymmetry upper bound and conclusion}

\begin{lemma}[Asymmetry bound]\label{lem:fuglede}
For a normalised radial graph \eqref{eq:norm} with
$\|\varphi\|_{L^\infty}\le\tfrac1{150}$ and barycentre $0$,
\begin{equation}\label{eq:asym}
   \AF(G)^2\ \le\ \frac{16}{\pi}\,\rho^2 .
\end{equation}
\end{lemma}

\begin{proof}
As $|G|=\pi$, the concentric unit ball is admissible in
\eqref{eq:dstarAdef}, so $\AF(G)\le\pi^{-1}|G\triangle B_1|$ and, in polar
form,
\[
   |G\triangle B_1|=\int_0^{2\pi}\Bigl|\varphi+\tfrac12\varphi^2\Bigr|\d\theta
   \le\bigl(1+\tfrac12\eta_G\bigr)\!\int_0^{2\pi}\!|\varphi|
   \le\tfrac{301}{300}\sqrt{2\pi}\,\|\varphi\|_{L^2}.
\]
Hence $\AF(G)^2\le\pi^{-2}(\tfrac{301}{300})^2 2\pi\|\varphi\|_{L^2}^2\le
\tfrac3\pi\|\varphi\|_{L^2}^2$, and by \cref{lem:normmodes}
$\|\varphi\|_{L^2}^2\le\tfrac{11}{10}\pi\rho^2$, so
$\AF(G)^2\le\tfrac{33}{10}\rho^2\le\tfrac{16}\pi\rho^2$ (as
$\tfrac{33}{10}<\tfrac{16}\pi\approx5.09$).
\end{proof}

\begin{proof}[Proof of \cref{thm:graphG}]
After the dilation and translation reducing to \eqref{eq:norm} with
barycentre $0$, \cref{prop:coercive} gives $\dstar(G)\ge\rho^2$ and
\cref{lem:fuglede} gives $\AF(G)^2\le\tfrac{16}\pi\rho^2$. Therefore
\[
   \AF(G)^2\le\frac{16}\pi\rho^2\le\frac{16}\pi\,\dstar(G)\le 48\,\dstar(G),
\]
since $\tfrac{16}\pi\approx5.09<48$. This is \eqref{eq:graphG} with
$\eta_G=\tfrac1{150}$, $C_G=48$ (any constant $\ge\tfrac{16}\pi$ works; we
report the safe value $48$).
\end{proof}

\begin{remark}[Sharpness and what is used]\label{rem:sharp}
The exponent $2$ in \eqref{eq:graphG} is sharp: for the pure mode
$\varphi=a\cos2\theta$ (a small ellipse), after normalisation
$\AF(G)^2\sim c_1 a^2$ and $\dstar(G)\sim c_2 a^2$ as $a\to0$ with
$c_1,c_2>0$, so $\AF(G)^2/\dstar(G)$ has a finite positive limit and no
exponent larger than $2$ can hold. The proof uses only the two variational
characterisations of torsion (\cref{lem:dual}; elementary), harmonicity of
$r^k e^{\mathrm i k\theta}$, the divergence theorem on the disc, Parseval, and
the volume/barycentre normalisations — \emph{no} quantitative isoperimetric
inequality, \emph{no} selection or symmetrisation, and \emph{no} boundary
regularity beyond the Lipschitz graph structure.
\end{remark}

\section{Proof of the main theorem}
\label{sec:main}

We now assemble the constructions of \S\S4--7 with the perturbative
nearly--spherical estimate of \S8. Throughout, $\Om\subset\R^2$ is bounded
open with $\abs\Om=\pi$, $u\in H^1_0(\Om)$ is its torsion function, $\dT=\dT(\Om)=
\tfrac\pi8-T(\Om)\ge0$ is the Saint--Venant deficit, and $\dstar(\Om)=
\tfrac8\pi\dT$ is the scale--invariant relative torsion deficit. We write
$C$ for an absolute constant whose value may change from line to line, and
$\eta_G,C_G$ for the constants of the radial--graph theorem
\cref{thm:graphG}, which depend only on the absolute mode-one constant $K$ of
\eqref{eq:mode1-scoped} below and are therefore themselves absolute. The
threshold $\delta_\star>0$ is absolute and may be shrunk finitely many times.

The section has two results. \Cref{thm:reduction} reduces a near--extremal
$\Om$ to a measurable small--amplitude radial graph $G$ that inherits the
torsion deficit and is close to $\Om$ in measure; this is the entire payload
of \S\S4--7. \Cref{thm:main} then feeds $G$ into \cref{thm:graphG} and
transfers the resulting asymmetry bound back to $\Om$ by the triangle
inequality.

\subsection{The reduction}

The construction selects an inset level $\hat v\sim\sqrt\dT$, passes to the
trapping annulus, cuts the radial folds to make an honest graph, and
optimises the free scale $\tau$ at $\tau=\sqrt\dT$. We carry $\tau$ symbolic
through the estimates and set $\tau=\sqrt\dT$ only at the end, so that the
balance producing the exponent $\tfrac12$ is visible.

\begin{theorem}[Sharp reduction to a graph inset]\label{thm:reduction}
There is an absolute $\delta_\star>0$ such that the following holds. Let
$\Om\subset\R^2$ be bounded open with $\abs\Om=\pi$ and $\dT\le\delta_\star$,
and put
$\tau:=\sqrt\dT$. Then there exist a point $z\in\R^2$ and a bounded open
star-shaped radial graph
\[
   G=\{z+r\omega:\ 0\le r<\rho(\omega),\ \omega\in S^1\},
   \qquad \rho\colon S^1\to(0,\infty)\ \text{measurable},
\]
with $\hat r_G:=\sqrt{\abs G/\pi}$, such that
\begin{enumerate}[label=\textup{(\roman*)},leftmargin=2.6em]
\item \textup{(small amplitude)}
   $\ \norm{\dfrac{\rho}{\hat r_G}-1}_{L^\infty(S^1)}\ \le\ C\,\dT^{1/8}
   \ \le\ \eta_G$;
\item \textup{(controlled deficit)} $\ \dstar(G)\ \le\ C\,\dstar(\Om)$;
\item \textup{(closeness)} $\ \abs{\Om\triangle G}\ \le\ C\,\sqrt\dT$.
\end{enumerate}
Every constant is absolute. The construction uses only the two level
deficits $D_H,D_I$ (through \cref{thm:level-identity,lem:profile,lem:muprime,%
thm:annulus,lem:two-flux}); it invokes no selection principle, no transport,
and no Serrin/Magnanini--Poggesi symmetry estimate.
\end{theorem}

\begin{proof}
Throughout we abbreviate $\diso=\diso(\Etil)$, $\Exc=P(\Etil)-2\pi\hat r$,
$a=\abs{\Etil\triangle B_{\hat r}(z)}$ and $\eta=\rho_e-\rho_i$, with the
notation of \cref{def:inset,thm:annulus}.

\emph{Step 0: the inset.} Apply \cref{lem:level} with the given $\tau\in
(0,\tfrac1{16})$ (legitimate since $\tau=\sqrt\dT\le\sqrt{\delta_\star}<
\tfrac1{16}$ once $\delta_\star<\tfrac1{256}$): there is a regular value
$\hat v\in(\tau,2\tau)$ with
\begin{equation}\label{eq:m-good}
   D_H(\hat v)+D_I(\hat v)\le A_0\,\frac\dT\tau,
   \qquad
   \mu_B(\hat v)-\mu(\hat v)\le A_0\,\frac\dT\tau,
\end{equation}
$A_0$ absolute. Let $\Etil_0$ be the connected inset of \cref{def:inset} (the
component of $\{u>\hat v\}$ of largest measure); by \cref{lem:noholes} it is
connected and has no interior holes, and $\mu_E=\abs{\Etil_0}$,
$\hat r=\sqrt{\mu_E/\pi}$. By \cref{lem:close,lem:spill},
\begin{equation}\label{eq:m-close}
   \abs{\Om\triangle\Etil_0}\ \le\ 8\pi\tau+A_0\frac\dT\tau
       +\frac{A_0^2}{16\pi^3}\,\frac{\dT^2}{\tau^2}.
\end{equation}
Since $\mu_E=\abs{\Etil_0}=\pi-\abs{\Om\triangle\Etil_0}$ (as $\Etil_0\subset\Om$),
\eqref{eq:m-close} gives $\mu_E\ge\pi-(8\pi\tau+A_0\dT/\tau+A_0^2\dT^2/
(16\pi^3\tau^2))$; for $\delta_\star$ small (so that the right--hand bound is
$\le\pi/2$, which is exactly the standing hypothesis \eqref{eq:halfmass} of
\cref{sec:inset} at $\tau=\sqrt\dT$, with the higher--order spillover term
absorbed) we obtain $\mu_E\ge\pi/2$, hence $\hat r=\sqrt{\mu_E/\pi}\ge1/\sqrt2$
and, for the annulus produced below, $\rho_i\ge\tfrac12$ after possibly
shrinking $\delta_\star$.

\emph{Step $0'$: hole-filling.} Pass from $\Etil_0$ to its saturation
$\Etil:=E^h$ of \cref{def:fill}, with all bounded complementary components
filled in (\cref{conv:fill}). By \cref{lem:fill}\ref{it:fill-top} the set
$\Etil$ is bounded, open, connected and \emph{globally hole-free} ---
$\R^2\setminus\overline\Etil$ connected --- \emph{unconditionally}, with no
hypothesis on $\partial\Om$; this discharges the exterior-connectedness
convention \cref{lem:ext-conn} required by the trapping annulus and the
inscribed-ball lemma. The filled mass obeys $H:=\abs{\Etil\setminus\Etil_0}
\le D_I(\Etil_0)^2/(64\pi^2\mu_E)\le C(\dT/\tau)^2$
(\cref{lem:fill}\ref{it:fill-key},\ref{it:fill-scales}), so at $\tau=\sqrt\dT$,
$H\le C\dT$, of the same higher order as the spillover. By
\cref{lem:fill}\ref{it:fill-scales} all inset scales are preserved with the
same rates: $\diso(\Etil)\le\diso(\Etil_0)$, $\Exc(\Etil)\le\Exc(\Etil_0)+CH$,
$a(\Etil)\le a(\Etil_0)+CH$, $\dstar(\Etil)\le\dstar(\Etil_0)+CH/\mu_E$, and
the closeness deteriorates only by $H$:
\begin{equation}\label{eq:m-fillclose}
   \abs{\Om\triangle\Etil}\ \le\ \abs{\Om\triangle\Etil_0}+H
   \ \le\ 8\pi\tau+A_0\frac\dT\tau+\frac{A_0^2}{16\pi^3}\frac{\dT^2}{\tau^2}+CH.
\end{equation}
From here on $\Etil$ denotes the hole-filled inset, $\mu_E:=\abs\Etil$ (now
$\mu_E+H$ of the previous step, still $\ge\pi/2$), $\hat r:=\sqrt{\mu_E/\pi}$,
and we run Steps~1--4 on it; all of \cref{sec:annulus,sec:fold,sec:cut} apply
to it directly (\cref{conv:fill}).

\emph{Step 1: the scales (carrying $\tau$ symbolic).} From the inset control
\eqref{eq:diso-inset} of \cref{prop:insetcontrols} (which includes the
spillover term $4\pi S$ and absorbs it),
\begin{equation}\label{eq:m-diso}
   \diso\ \le\ C\,\frac\dT\tau,
   \qquad
   \Exc=2\pi\hat r\,\diso\ \le\ C\,\frac\dT\tau,
\end{equation}
using $\hat r\le1$. The sharp quantitative
isoperimetric inequality \cref{thm:fmp}, in the form
$a=\abs{\Etil\triangle B_{\hat r}(z)}\le C_{\mathrm F}\sqrt{\diso}\,\abs\Etil$,
gives
\begin{equation}\label{eq:m-a}
   a\ \le\ C\,\sqrt{\diso}\ \le\ C\Bigl(\frac\dT\tau\Bigr)^{1/2}.
\end{equation}
The trapping annulus \cref{thm:annulus}, applied to the hole-filled inset
$\Etil=E^h$ (globally hole-free unconditionally by \cref{lem:fill}\ref{it:fill-top},
\cref{conv:fill}), yields
$B_{\rho_i}(z)\subset\Etil\subset B_{\rho_e}(z)$ with
\begin{equation}\label{eq:m-eta}
   \eta=\rho_e-\rho_i\ \le\ C\,(\Exc+a)^{1/2}.
\end{equation}
By \cref{cor:gd} the radial--graph defect of $\Etil$ about $z$ obeys
\begin{equation}\label{eq:m-gd}
   \gd(\Etil)\ \le\ C\,\eta\,(\Exc+a\eta).
\end{equation}

\emph{Step 2: optimise $\tau=\sqrt\dT$ and re--derive every exponent.} Set
$\tau=\sqrt\dT$, so $\dT/\tau=\sqrt\dT$. We re--derive each scale
independently from \eqref{eq:m-diso}--\eqref{eq:m-gd}:
\begin{align}
   \diso\ &\le\ C\,\frac\dT\tau\ =\ C\,\dT^{1/2}, &
   \Exc\ &\le\ C\,\dT^{1/2}, \label{eq:s-diso}\\
   a\ &\le\ C\Bigl(\frac\dT\tau\Bigr)^{1/2}\ =\ C\,\dT^{1/4}, \label{eq:s-a}\\
   \eta\ &\le\ C\,(\Exc+a)^{1/2}\ \le\ C\,a^{1/2}\ \le\ C\,\dT^{1/8},
       \label{eq:s-eta}
\end{align}
where in \eqref{eq:s-eta} we used that $a\sim\dT^{1/4}$ dominates
$\Exc\sim\dT^{1/2}$ as $\dT\to0$, so $\Exc+a\le 2a\le C\dT^{1/4}$ and
$(\Exc+a)^{1/2}\le C\dT^{1/8}$. For the graph defect, \eqref{eq:m-gd} gives
\begin{equation}\label{eq:s-gd}
   \gd(\Etil)\ \le\ C\,\eta\,(\Exc+a\eta)
   \ \le\ C\bigl(\dT^{1/8}\!\cdot\dT^{1/2}+\dT^{1/4}\!\cdot\dT^{1/4}\bigr)
   \ =\ C\bigl(\dT^{5/8}+\dT^{1/2}\bigr)\ \le\ C\,\dT^{1/2},
\end{equation}
the dominant term being $a\eta\cdot\eta$ with
$a\,\eta^2=\dT^{1/4}\cdot\dT^{1/4}=\dT^{1/2}$. As a cross--check against the
heuristic $\gd\sim\dT/\tau$: at $\tau=\sqrt\dT$ one has
$\dT/\tau=\dT^{1/2}$, matching \eqref{eq:s-gd}. Likewise the closeness of the
hole-filled inset \eqref{eq:m-fillclose} reads
\begin{equation}\label{eq:s-close}
   \abs{\Om\triangle\Etil}\ \le\ 8\pi\,\dT^{1/2}+A_0\,\dT^{1/2}
       +\frac{A_0^2}{16\pi^3}\,\dT+CH\ \le\ C\,\dT^{1/2},
\end{equation}
i.e.\ both contributions $\tau=\dT^{1/2}$ (the level height) and
$\dT/\tau=\dT^{1/2}$ (the level smoothing) coincide at the optimum, while the
spillover $\dT^2/\tau^2=\dT$ and the filled mass $H\le C\dT$ are genuinely
higher order. This is precisely the
balance $\tau\sim\dT/\tau$, solved by $\tau=\sqrt\dT$; any other power of
$\tau$ would leave one of the two terms larger than $\dT^{1/2}$.

\emph{Step 3: realisation as the measurable core graph $G$.} We take $G:=G_0$
the measurable open star-shaped core graph of \cref{prop:realization}, with
profile $\rho:=\rho_0\colon S^1\to[\rho_i,\rho_e]$ measurable. No
mollification or continuity is required: the nearly-spherical estimate
\cref{thm:graphG} admits measurable profiles. By
\cref{prop:realization}\ref{it:real-close} and \cref{cor:gd},
\begin{equation}\label{eq:s-EtilG}
   \abs{\Etil\triangle G}\ =\ \abs{\Etil\setminus G_0}
   \ \le\ C\,\eta(\Exc+a\eta)\ \le\ C\,\dT^{1/2},
\end{equation}
the last step being \eqref{eq:s-gd}, and
$\norm{\rho/\hat r_G-1}_{L^\infty}\le2\eta$ by
\cref{prop:realization}\ref{it:real-amp} with $\hat r_G=\hat r_{G_0}=
\sqrt{\abs G/\pi}\in[\rho_i,\rho_e]$.

\emph{Step 4: verification of (i), (ii), (iii).}

\noindent\textup{(i)} \emph{Amplitude.} By
\cref{prop:realization}\ref{it:real-amp} the profile $\rho=\rho_0$ obeys
$\norm{\rho/\hat r_G-1}_\infty\le2\eta$ with $\hat r_G=\sqrt{\abs G/\pi}\in
[\rho_i,\rho_e]$, and by \eqref{eq:s-eta} $\eta\le C\dT^{1/8}$, so
$\norm{\rho/\hat r_G-1}_\infty\le C\dT^{1/8}$. Shrinking $\delta_\star$ so
that $C\,\delta_\star^{1/8}\le\eta_G$ gives
$\norm{\rho/\hat r_G-1}_\infty\le\eta_G$, which is (i). Here it is essential
that $\eta\sim\dT^{1/8}$ is an \emph{absolute} smallness (decaying with
$\dT$), so the hypothesis of \cref{thm:graphG} is met for $\dT$ small,
\emph{regardless} of the cruder $L^1$--roundness $a\sim\dT^{1/4}$.

\noindent\textup{(iii)} \emph{Closeness.} By \eqref{eq:s-close} and
\eqref{eq:s-EtilG}, with $\Etil=E^h$ so that $\abs{\Om\triangle\Etil}$ already
includes the filled mass $H\le C\dT$ (\eqref{eq:m-fillclose}),
\[
   \abs{\Om\triangle G}\ \le\ \abs{\Om\triangle\Etil_0}+H+\abs{\Etil\triangle G}
   \ =\ \abs{\Om\triangle\Etil}+\abs{\Etil\triangle G}
   \ \le\ C\,\dT^{1/2}+C\,\dT^{1/2}\ =\ C\,\sqrt\dT,
\]
which is (iii); the filled mass $H\le C\dT$ is negligible.

\noindent\textup{(ii)} \emph{Deficit.} Here $G=G_0=
\{z+r\omega:r<\rho_0(\omega)\}$ is the measurable core graph of
\cref{cor:gd,prop:realization}, so $\rho_0\in
[\rho_i,\rho_e]$, $B_{\rho_i}(z)\subset G_0\subset\Etil\subset B_{\rho_e}(z)$,
and $\Etil\setminus G_0$ (the radial caps) lies in the collar
$\{\rho_i<\abs{x-z}<\rho_e\}$. The torsion-loss decomposition
\eqref{eq:torsion-loss} of \cref{lem:torsion-stab} holds unconditionally,
\[
   T(\Etil)-T(G_0)=\int_{\Etil\setminus G_0}u_{\Etil}+\int_{G_0}h
   \ \le\ \tilde\eta\,g+\int_{G_0}h,
   \qquad g=\abs{\Etil\setminus G_0},
\]
with $\tilde\eta=\tfrac14(\rho_e^2-\rho_i^2)\le\eta\,\rho_e\le C\dT^{1/8}$ and
the cap term $\int_{\Etil\setminus G_0}u_{\Etil}\le\tilde\eta\,g$ rigorous. The
torsion density is an absolute constant bounded below,
\[
   \frac{T(\Etil)}{\abs\Etil}
   =\bigl(1-\dstar(\Etil)\bigr)\frac{\abs\Etil}{8\pi}
   \ \ge\ \bigl(1-\tfrac12\bigr)\frac{\pi/2}{8\pi}=\frac1{32},
\]
using $\dstar(\Etil)\le\tfrac12$ (from Step~5 below, equivalently
\cref{prop:insetcontrols}\,\eqref{eq:dstar-inset}) and $\abs\Etil\ge\pi/2$.
Under the boundary-layer control \eqref{eq:flux-hyp} of
\cref{gap:torsion-stab} --- the single residual analytic point, namely the
localisation $\int_{G_0}h\le\tilde\eta\,g$ of the harmonic boundary-layer
term to the cap region --- the loss-slope hypothesis \eqref{eq:cutcond} holds
with $\lambda=2\tilde\eta\le\tfrac1{32}\le T(\Etil)/\abs\Etil$
(\cref{rem:supply,rem:roomtospare}), so \cref{lem:cutdeficit} applies to
$G=G_0\subset\Etil$ and gives directly, with no further regularisation step
(since $G=G_0$),
\begin{equation}\label{eq:s-defchain}
   \dstar(G)\ =\ \dstar(G_0)\ \le\ \dstar(\Etil).
\end{equation}
This is the only place in the reduction where a residual analytic point is
invoked; everything else --- the exact decomposition, the cap term, and the
deficit algebra --- is unconditional.

\emph{Step 5: transfer of the deficit from $\Om$ to $\Etil$.} The deficit of
the raw component $\Etil_0$ is controlled by the inset control
\eqref{eq:dstar-inset} of \cref{prop:insetcontrols},
$\dstar(\Etil_0)\le(\pi^2/\mu_E^2)\,\dstar(\Om)\le C\,\dstar(\Om)$; and the
hole-filling raises it only by the higher-order term $CH/\mu_E\le C\dT$
(\cref{lem:fill}\ref{it:fill-scales}, eq.~\eqref{eq:fill-dstar}), so
\begin{equation}\label{eq:s-transfer}
   \dstar(\Etil)\ \le\ \dstar(\Etil_0)+\frac{CH}{\mu_E}
   \ \le\ \frac{\pi^2}{\mu_E^2}\,\dstar(\Om)+C\dT\ \le\ C\,\dstar(\Om).
\end{equation}
We stress that the transfer to $\Etil_0$ is established in \cref{sec:inset}
\emph{without} appeal to the cutting principle \cref{lem:cutdeficit}: the
passage from the super--level set $\{u>\hat v\}$ to its largest component
$\Etil_0$ is handled there by the exact additivity of torsion over the
(disjoint) level components and the Saint--Venant bound on the spillover
torsion (\cref{prop:insetcontrols}, eqs.~\eqref{eq:Tadd}--\eqref{eq:Tspill}),
both internal to \cref{sec:identity,sec:inset}; and the hole-filling step
\eqref{eq:fill-dstar} uses only domain monotonicity of torsion and the area
bound $H\le C\dT$, again internal to \cref{sec:inset}. In particular the
spillover components need \emph{not} lie in the trapping collar of $\Etil$ ---
so \cref{lem:cutdeficit}, whose hypothesis requires the excised set to sit in
a thin concentric collar, is \emph{not} invoked here, and there is no circular
dependence on \cref{sec:cut} (see \cref{rem:fwd}). With $\mu_E\ge\pi/2$
(Step~0) the prefactor is $\pi^2/\mu_E^2\le4$, so for $\delta_\star$ small
\eqref{eq:s-transfer} gives $\dstar(\Etil)\le 5\,\dstar(\Om)$; in particular
$\dstar(\Etil)\le\tfrac12$ once $\dstar(\Om)\le\tfrac1{10}$, which justifies
the value used in Step~4. Inserting \eqref{eq:s-transfer} into
\eqref{eq:s-defchain} and using $\dT=\tfrac\pi8\dstar(\Om)\le C\dstar(\Om)$,
\[
   \dstar(G)\ \le\ \dstar(\Etil)+C\dT\ \le\ C\dstar(\Om)+C\dstar(\Om)
   \ =\ C\,\dstar(\Om),
\]
which is (ii). This completes the proof.
\end{proof}

\subsection{The main theorem}

\begin{theorem}[Sharp quantitative Saint--Venant, selection--free]
\label{thm:main}
There is an absolute constant $C$ such that every bounded open set
$\Om\subset\R^2$ with $\abs\Om=\pi$ satisfies
\begin{equation}\label{eq:main}
   \AF(\Om)^2\ \le\ C\,\dT(\Om).
\end{equation}
\end{theorem}

\begin{proof}
\emph{Trivial regime.} If $\dT\ge\delta_\star$, then since $\abs\Om=\abs B=\pi$
for the optimal disc $B$ we have $\abs{\Om\triangle B}\le\abs\Om+\abs B=2\pi$,
so $\AF(\Om)\le2$ and
$\AF(\Om)^2\le4\le(4/\delta_\star)\,\dT$. Thus \eqref{eq:main} holds with
$C=4/\delta_\star$ in this regime.

\emph{Near--extremal regime.} Assume $\dT\le\delta_\star$ with $\delta_\star$
the threshold of \cref{thm:reduction}, and let $G=\{z+r\omega:r<\rho(\omega)\}$
be the radial graph it provides, with $\hat r_G=\sqrt{\abs G/\pi}$ and
amplitude $\norm{\rho/\hat r_G-1}_\infty\le C\dT^{1/8}\le\eta_G$.

\Cref{thm:graphG} now applies directly to the graph $G$ as it is produced by
\cref{thm:reduction}, radial about the annulus centre $z$, with no
barycentre recentring: its hypotheses are the volume normalisation
$\hat r_G=\sqrt{|G|/\pi}$ (built into \cref{thm:reduction}), the $L^\infty$
amplitude bound $\|\rho/\hat r_G-1\|_\infty\le\eta_G$ (part~(i)), and the
quadratic mode-one bound \eqref{eq:mode1hyp}, i.e.\ $\|P_1\varphi\|_{L^2}\le
K\|\varphi\|_{L^2}^2$ for the profile $\varphi=\rho/\hat r_G-1$ in polar
coordinates about $z$. The first two hold by \cref{thm:reduction}. The third
is supplied by the choice of centre, as we now record.

\begin{remark}[The mode-one bound at the $L^1$-optimal centre]%
\label{gap:barycentre}
The centre $z$ is the \emph{exact} $L^1$-optimal (Fraenkel) centre of the
inset $\Etil$ (\cref{lem:zinterior}), an interior point with a definite
inscribed ball. Exact $L^1$-optimality is a first-order extremality condition
on the comparison ball: for the optimal centre the directional derivative of
$p\mapsto|\Etil\triangle B_{\hat r}(z+p)|$ vanishes, which is the statement
that the translation mode of $\partial\Etil$ carries no first-order
$L^1$-content. This is precisely the mechanism behind the quadratic mode-one
bound \eqref{eq:mode1hyp}: it forces the first Fourier mode of the radial
profile to be of second order in the amplitude, exactly as the barycentre
condition does through \cref{rem:mode1}, but \emph{at the centre $z$ already
in hand}, with no recentring of the set. What remains to make this fully
quantitative for the measurable core graph $G$ (as opposed to $\Etil$) is the
passage of the first-order extremality from the $L^1$-content of $\partial
\Etil$ to the $L^2$ mode-one coefficient of the profile of $G$, with the
collar error $\abs{\Etil\triangle G}\le C\sqrt\dT$ of \cref{thm:reduction}
absorbed. We isolate this as the explicit scoped hypothesis:
\begin{equation}\label{eq:mode1-scoped}
   \|P_1\varphi\|_{L^2(S^1)}\ \le\ K\,\|\varphi\|_{L^2(S^1)}^2
   \qquad(\varphi=\rho/\hat r_G-1\text{ about the }L^1\text{-optimal }z),
\end{equation}
for an absolute $K$. Under \eqref{eq:mode1-scoped} the argument below is
unconditional. This replaces the earlier, false-as-stated requirement of a
barycentre recentring --- which need not preserve single-valuedness of the
graph under an $L^\infty$ amplitude bound alone --- by a precise condition at
the existing centre, and removes the recentring obstruction entirely from
\cref{thm:graphG}. By \cref{lem:shift-asym} the asymmetry side needs no
mode-one hypothesis at all, and by \cref{rem:mode1-residual} the bound
\eqref{eq:mode1-scoped} is equivalent, for the present purpose, to the quartic
translation estimate $\dstar(G)\ge c\,\|P_1\varphi\|_{L^2}^4$
\eqref{eq:quartic-residual} --- the flatness of the translation direction of
the torsional deficit; this is the precise form of the residual.
\end{remark}

\emph{Asymmetry of $G$.} By \cref{thm:graphG} (with the mode-one bound
\eqref{eq:mode1-scoped}) and \cref{thm:reduction}(ii),
\begin{equation}\label{eq:m-AG}
   \AF(G)^2\ \le\ C_G\,\dstar(G)\ \le\ C_G\cdot C\,\dstar(\Om)
   \ =\ C\,\dstar(\Om)\ =\ C\,\tfrac8\pi\,\dT\ \le\ C\,\dT,
\end{equation}
so $\AF(G)\le C\sqrt\dT$.

\emph{Transfer to $\Om$ by the triangle inequality.} Let $B_G=B(y,\hat r_G)$
be the disc of area $\abs G$ realising the asymmetry of $G$, i.e.\ centred at
the optimal $y$ so that $\abs{G\triangle B_G}=\AF(G)\,\abs G$, and let
$B=B(y,1)$ be the disc of area $\pi$ \emph{concentric} with $B_G$. Since
$B_G$ and $B$ are concentric discs of areas $\abs G$ and $\pi$, their
symmetric difference is the annulus of area
\begin{equation}\label{eq:m-discs}
   \abs{B_G\triangle B}\ =\ \bigl|\,\abs G-\pi\,\bigr|
   \ =\ \bigl|\,\abs G-\abs\Om\,\bigr|\ \le\ \abs{\Om\triangle G}.
\end{equation}
Because $\abs\Om=\pi=\abs B$, the disc $B$ is admissible in the definition of
$\AF(\Om)$, so by the triangle inequality for the symmetric difference,
\begin{align*}
   \AF(\Om)\ \le\ \frac{\abs{\Om\triangle B}}{\pi}
   &\ \le\ \frac{\abs{\Om\triangle G}+\abs{G\triangle B_G}
            +\abs{B_G\triangle B}}{\pi}\\
   &\ \le\ \frac1\pi\Bigl(\abs{\Om\triangle G}+\AF(G)\,\abs G
            +\abs{\Om\triangle G}\Bigr),
\end{align*}
using \eqref{eq:m-discs} and $\abs{G\triangle B_G}=\AF(G)\abs G$. With
$\abs G=\pi+O(\sqrt\dT)\le2\pi$, $\abs{\Om\triangle G}\le C\sqrt\dT$ by
\cref{thm:reduction}(iii), and $\AF(G)\le C\sqrt\dT$ by \eqref{eq:m-AG},
\begin{equation}\label{eq:m-final}
   \AF(\Om)\ \le\ \frac1\pi\bigl(C\sqrt\dT+\AF(G)\,\abs G+C\sqrt\dT\bigr)
   \ \le\ C\sqrt\dT+C\,\AF(G)\ \le\ C\sqrt\dT.
\end{equation}
Squaring \eqref{eq:m-final} gives $\AF(\Om)^2\le C\,\dT$, which is
\eqref{eq:main} in the near--extremal regime. Combining the two regimes
proves the theorem.
\end{proof}

\begin{remark}[What was used, what was not, and the two sharp ingredients]
\label{rem:audit}
The reduction \cref{thm:reduction} draws on the two level deficits and
nothing else of substance:
\begin{itemize}[leftmargin=1.5em]
\item $D_H+D_I$ enters through the level identity
   $\int_0^m(D_H+D_I)\,dv=4\pi\dT$ (\cref{thm:level-identity}), which powers
   both the level selection \cref{lem:level} and the deficit transfer
   \cref{lem:transfer}, and through the pointwise bound $-\mu'\ge4\pi$
   (\cref{lem:muprime}) and the profile identity (\cref{lem:profile}) that
   carry closeness;
\item $D_I$ enters through the isoperimetric/annulus inputs
   (\cref{thm:annulus,thm:fmp}) and the multi--component spillover bound
   (\cref{lem:spill}).
\end{itemize}
At no point is the Brasco--De~Philippis--Velichkov selection principle used,
nor any mass--transport (moving--ball, kinetic) argument, nor a Serrin /
Magnanini--Poggesi symmetry estimate. The only nonlinear stability input is
the \emph{perturbative} nearly--spherical bound \cref{thm:graphG}, applied to
a graph whose amplitude is already absolutely small.

Two ingredients make the reduction \emph{sharp}, i.e.\ produce the exponent
$\tfrac12$ rather than $\tfrac14$:
\begin{enumerate}[label=\textup{(\alph*)},leftmargin=1.8em]
\item the fold--content bound \cref{prop:fold}, in which subtracting the
   constant $1/\hat r$ (legitimate because $\abs\Etil=\abs{B_{\hat r}(z)}$)
   and confining $\Etil\triangle B_{\hat r}(z)$ to the $\eta$--annulus
   replaces $a$ by $a\eta$, so the graph defect avoids the
   quantitative--isoperimetric square--root loss; without it the fold content
   would carry the full $a\sim\sqrt\diso$ and cap the result at
   $\dT^{1/4}$ closeness;
\item the cutting lemma \cref{lem:cutdeficit}: excising the low--torsion
   boundary caps to manufacture the graph $G$ \emph{decreases} the
   scale--invariant deficit $\dstar$, because the equal--area comparison ball
   shrinks faster (by $\gd/4$ in $\tilde\eta$--weighted terms) than the
   torsion is lost.
\end{enumerate}
Finally, the optimisation $\tau=\sqrt\dT$ is the balance between the two
$\tau$--dependent contributions to closeness: the level height $\sim\tau$
grows in $\tau$, while the smoothing/graph defect $\sim\dT/\tau$ decays in
$\tau$; equating them, $\tau\sim\dT/\tau$, forces $\tau=\sqrt\dT$, at which
both equal $\sqrt\dT$ and the amplitude is $\eta\sim(\dT/\tau)^{1/4}=
\dT^{1/8}$, the absolute smallness the graph theorem requires.
\end{remark}


\bibliographystyle{amsplain}
\bibliography{refs}

\end{document}
